\documentclass[12pt]{article}
\usepackage{jheppub}
\usepackage{amssymb,amsfonts}
%\usepackage{draftfil}
%\usepackage{showkeys}
%\usepackage{epsfig}
%gives names of refs and cites
%If you do not have the msbm fonts, delete the following 10 lines
\font\mybb=msbm10 at 10pt
%\font\mybb=msbm12 at 12pt

\renewcommand{\arraystretch}{1.2}

% \usepackage[square, comma, compress,numbers]{natbib}
\usepackage{natbib}
\setcitestyle{square, comma, numbers, compress}
\usepackage{wrapfig}
\usepackage{epsfig}
\usepackage[T1]{fontenc}
\usepackage[utf8]{inputenc}
\usepackage{lmodern}
\usepackage{float}
\usepackage{placeins}
\usepackage{dsfont}
\usepackage{arydshln}
\usepackage{caption}

\usepackage[dvipsnames]{xcolor}
\usepackage{tikz}
%%%%%%%%%%%%%%%%%%%%%%%%%%    My Macros    %%%%%%%%%%%%%%%%%%%%%%%%%
%%%%%%%%%%%%%%%%%%%%%%%%%%%%%%%%%%%%%%%%%%%%%%%%%%%%%%%%%%%%%%%%%%%%

\def\be{\begin{eqnarray}}
\def\ee{\end{eqnarray}}
\newcommand{\nn}{\nonumber}
\newcommand\para{\paragraph{}}
\newcommand{\ft}[2]{{\textstyle\frac{#1}{#2}}}
\newcommand{\eqn}[1]{(\ref{#1})}
\newcommand{\pl}[1]{\frac{\partial L}{\partial{#1}}}
\newcommand{\ppp}[2]{\frac{\partial {#1}}{\partial {#2}}}
\newcommand{\ph}[1]{\frac{\partial H}{\partial{#1}}}

\newcommand\balpha{\mbox{\boldmath $\alpha$}}
\newcommand\bbeta{\mbox{\boldmath $\beta$}}
\newcommand\bgamma{\mbox{\boldmath $\gamma$}}
\newcommand\bomega{\mbox{\boldmath $\omega$}}
\newcommand\blambda{\mbox{\boldmath $\lambda$}}
\newcommand\bmu{\mbox{\boldmath $\mu$}}
\newcommand\bphi{\mbox{\boldmath $\phi$}}
\newcommand\bpsi{\mbox{\boldmath $\psi$}}
\newcommand\bzeta{\mbox{\boldmath $\zeta$}}
\newcommand\bsigma{\mbox{\boldmath $\sigma$}}
\newcommand\bepsilon{\mbox{\boldmath $\epsilon$}}
\newcommand\btau{\mbox{\boldmath $\tau$}}
\newcommand\beeta{\mbox{\boldmath $\eta$}}
\newcommand\btheta{\mbox{\boldmath $\theta$}}
\newcommand\valpha{\vec{\alpha}}
\newcommand\vg{\vec{g}}
\newcommand{\ra}{\rightarrow}
\newcommand{\ket}[1]{|#1\rangle}
\newcommand{\bra}[1]{\langle#1|}

\def\Dslash{\,\,{\raise.15ex\hbox{/}\mkern-12mu D}}
\def\Dbarslash{\,\,{\raise.15ex\hbox{/}\mkern-12mu {\bar D}}}
\def\delslash{\,\,{\raise.15ex\hbox{/}\mkern-9mu \partial}}
\def\delbarslash{\,\,{\raise.15ex\hbox{/}\mkern-9mu {\bar\partial}}}
\def\pslash{\,\,{\raise.15ex\hbox{/}\mkern-9mu p}}
\def\calDslash{\,\,{\raise.15ex\hbox{/}\mkern-12mu {\cal D}}}
\newcommand\Bprime{B${}^\prime$}
\newcommand{\sign}{{\rm sign}}
\newcommand{\D}{{\cal D}}

\newcommand{\p}{\partial}



\newcommand{\co}{\color{olive}}
\newcommand{\cblue}{\color{blue}}
\newcommand{\cred}{\color{red}}

\newcommand\bte{\tilde{\bf e}}
\newcommand{\NN}{{\mathbb N}}
\newcommand{\Z}{{\mathds Z}}
\newcommand{\ZZ}{{\mathbb Z}}
\newcommand{\Q}{{\mathds Q}}
\newcommand{\QQ}{{\mathbb Q}}
\newcommand{\R}{{\mathds R}}
\newcommand{\RR}{{\mathbb R}}
\newcommand{\C}{{\mathds C}}
\newcommand{\CC}{{\mathbb C}}
\newcommand{\PP}{{\mathbb P}}
\newcommand{\e}{\,{\rm e}}
\newcommand{\CP}{{\bf CP}}
\newcommand{\Tr}{{\rm Tr}}
\newcommand{\bu}{{\bf u}}
\newcommand{\bk}{{\bf k}}
\newcommand{\bx}{{\bf x}}
\newcommand{\bp}{{\bf p}}
\def\implies{\Rightarrow}
\def\To{\longrightarrow}
\def\longlongrightarrow{\relbar\joinrel\relbar\joinrel\rightarrow}
\def\ridiculousrightarrow{\relbar\joinrel\relbar\joinrel\relbar%
\joinrel\rightarrow}
\def\calN{\mathcal{N}}
\def\calL{\mathcal{L}}
\def\calV{\mathcal{V}}
\def\calF{\mathcal{F}}
\def\calO{\mathcal{O}}
\def\calR{\mathcal{R}}
\def\calH{\mathcal{H}}
\def\calE{\mathcal{E}}

\def\underarrow#1{\mathrel{\mathop{\longrightarrow}\limits_{#1}}}
\def\onnearrow#1{\mathrel{\mathop{\nearrow}\limits^{#1}}}
\def\undernearrow#1{\mathrel{\mathop{\nearrow}\limits_{#1}}}
\def\onarrow#1{\mathrel{\mathop{\longrightarrow}\limits^{#1}}}
\def\onArrow#1{\mathrel{\mathop{\longlongrightarrow}\limits^{#1}}}
\def\OnArrow#1{\mathrel{\mathop{\ridiculousrightarrow}\limits^{#1}}}
\def\lae{\mathrel{\mathop{\smash{\lower .5 ex \hbox{$\stackrel<\sim$}}}}}
\def\lae{\mathrel{\mathop{\smash{\lower .5 ex \hbox{$\stackrel>\sim$}}}}}
\def\eqq{\stackrel?=}

\newcommand{\upket}{|\uparrow\rangle}
\newcommand{\downket}{|\downarrow\rangle}
\newcommand{\rightket}{|\rightarrow\rangle}
\newcommand{\leftket}{|\leftarrow\rangle}

\def\theequation{\thesection.\arabic{equation}}

\newcommand{\tvp}{\tilde{\varphi}}

\newcommand{\BZ}[1]{{\color [rgb]{1,0.3,1} [{BZ:} #1]}}
\newcommand{\RM}[1]{{\color{Red} [{RM:} #1]}}

\newcommand{\alyssa}{\textcolor{magenta}}
%%%%%%%%%%%%%%%%%%%%%%%%%%%%%%%%%%%%%%%%%%%%%%%%%%%%%


\title{A Course on Quantum Mechanics}

\author{Mukesh Bansal, Harry Cliff, Erin Crawley, Madelyn Leembruggen, and David Tong}



\abstract{These notes start to flesh out each lecture. They contain something resembling a script. It still needs a lot of polishing -- at the moment,  it's just a summary of things that I think would need saying to get the story across, together with what I would want to write on the board (or using powerpoint slides if necessary). Hopefully it also gives us a better sense of the time that would be needed for each section.}


\begin{document}
\pagestyle{plain} \setcounter{page}{1}
\newcounter{bean}
\baselineskip16pt \setcounter{section}{0}


\maketitle




%%%%%%%%%%%%%%%%%%%%%   NOTES START  %%%%%%%%%%%%%%%%%%%%%%%%%%%
%%%%%%%%%%%%%%%%%%%%%%%%%%%%%%%%%%%%%%%%%%%%%%%%%%%%%%%%%%%%%%%%%


\newpage
\noindent
{\Large \bf Section 1: The Basics}



\section{The Wavefunction}

{ \color{blue}Things, like this, in blue, are comments. Obviously we need a broad, rousing, visionary introduction. It would be easy to spend the whole 30 minutes doing this, and we might leave this open as a possibility. But starting with a kind of ramped-up pop science lecture might also might give a misleading impression about the general tone of the course. I like the idea of using the hydrogen atom demo to do this work for us, before jumping into something more concrete.}



 
 \subsection{The Hydrogen Atom}
 
The goal of these lectures is to start to get to grips with quantum mechanics, to understand its richness and strangeness and why it's the single most important development in all of science. We will need to build up quite a lot of mathematical technology over these lectures. But we're going to start by explaining where we're going. The problem that we want to finish this course by solving. 

\para
This is the same problem that motivated Heisenberg in 1925 when he was trapped on the island of Helgoland in the North Sea, desperate to escape his hayfever and fevered atomic dreams. It is the spectral lines of atoms.

\para
Shine a light at a gas of atoms and it will absorb only certain wavelengths. Turn the light off, and the atoms will emit light in those same wavelengths. it's like each atom, or each molecule, comes with a fingerprint, a bar code, uniquely specifying its identity.  This is called the spectrum of the atom. It's an extremely useful fact -- it allows us to detect what gases are made of without getting too close. It allowed us to figure out what the Sun and faraway stars are made of -- or at least their atmospheres -- without actually needing to visit. It's currently being used to detect the atmospheres of exoplanets -- planets that orbit other stars and there is some, possibly optimistic hope, that it might ultimately allow us to detect life on other planets, if they leave some biosignature in the atmosphere.

\para
But there is also a conceptual issue with these spectral lines. Which is that they're discrete. And, at least before this point, nothing else in physics was discrete. Physics was about fluids sloshing backwards and forwards, or optics or the motion of planets. Despite appearances, there's nothing discrete here. There's nothing discrete about any of this. Everything is continuous. Even where you think there's discreteness, it's largely illusory. When I was young there were nine planets in our solar system. Now there are eight. And while astronomers will bend over backwards trying to define a planet -- it has to be suitably round and clear its own path, blah, blah, blah -- the reality is there's stuff out there that ranges in size from a grain of sand to a giant ball of gas like jupiter and if you want to say what's a planet and what isn't then you have to draw a line somewhere and that's always going to be slightly arbitrary.  Because the world we live in is continuous.

\para
Until you get to the atomic level. Atoms themselves are discrete -- or, perhaps said better, their constituents are discrete. Electron and protons and neutron (or better still, quarks). But here, in the spectral lines, you get a very visual imprint of that discreteness. Why do atoms only shine at certain frequencies? How is this possible? How can we explain it and, better yet, predict the pattern of this barcode? ({\color{blue} It might be worth showing more than just the spectrum of hydrogen at this point.})

\para
This was the starting point for quantum mechanics, and will be the end point for this course. In fact, it's really what "quantum" means. For a physicist, quantum is synonymous with discreteness. (It comes from the latin meaning "how much" or "how great".)  One of the many great surprises of quantum mechanics, is that you start with this discreteness -- an absolutely certainty about what atoms are doing -- yet it leads, ultimately, into a world where nearly everything is uncertain. Where particles have a delocalised existence and can be in two places at the same time, a smeared wave of probability. A world where, in the extreme, cats can be alive and dead. So there's always this underlying tension in quantum mechanics, bordering on a paradox: something are certain, others much more fluid, more uncertain. 


\para
{\color{blue}There's an assumption here that this story isn't entirely new to the audience....that they've heard of the idea of a spectrum, even if long ago. The goal here is to highlight why it's surprising.}

\para
{\color{red} I'll write my estimation of times needed  in red. Maybe 10-15 minutes on this, with a demo. Perhaps slightly fleshed out. If we encourage questions, I could well imagine that it stretches to 15.}

\subsection{The Wavefunction}

We're going to start with the mathematics of quantum mechanics. Now, this isn't just a small tweak to the laws of physics. It's a root and branch overhaul. It's perhaps best thought of as changing the operating system of the universe. If you've got some law of physics -- say a law of gravity, or Hooke's law for a spring bouncing backwards and forwards -- then you have choice: you can either plug it into the old classical operating system. That's Newton's laws of motion. That will give you an approximate answer, that holds to wonderful accuracy when the objects involved are big -- tennis balls or planets. Or actual springs. Alternatively, you can plug it into the new quantum operating system. That gives the true answer. That's what we're going to learn.

\para
First, let me remind you of the classical operating system. This is really the equation that we all learned in school
%
\be F=ma\ .\nn\ee
%
This is Newton's second law. It tells us how a particle, with mass $m$, responds to a force $F$ -- it responds by accelerating with acceleration $a$. If the particle is moving in three-dimensions, then this is really a vector equation.

\para
{\color{olive} Anything that I write on a board, or present on slides, will be in this rather nice olive colour. To start, we'll just have:
\be  \mbox{Newton's second law: $F=ma$.}\nn\ee}
%
{\color{blue}I'd like to introduce notation for vectors by subsequently altering this equation. My standard notation is to draw a wiggle line under a vector but I was told recently that this is a weird English habit. Is that right? Obviously bold isn't possible on a blackboard, although it would work on slides, while the notation where you use a stupid arrow on top is extremely annoying if you're actually calculating anything. Either way, we need to decide.}

\para
This is, obviously, an important equation! The most important thing about it is that it tells us that a particle responds by accelerating. The force doesn't directly determine the velocity of a particle -- it determines its acceleration. It's important, in part, because it's counterintuitive. If you push something with a constant force, then it feels like it moves with a constant speed. {\color{blue}This is one experiment I can actually do! A book on a table would do the job.} But that's because there's all sorts of complicated friction things going on. As Galileo first realised, you should treat friction as yet another force that you add to this mix. What's really happening is that forces cause things to accelerate.

\para
For our purposes, the most important fact is the following. Newton's law tells you about neither the position of the particle, not its velocity. Instead it tells you about the acceleration. That means that if you want to solve this equation, and figure out what the particle is doing, then you need to specify that as extra information. You only need to specify this at some initial time, say $t_0$ -- after that, the equation will do the rest -- it tells you that it accelerates in such a such direction, and moves here. But you need to know the starting point. We call that starting point {\it the state} of the system. If you know the initial state, and you can solve the equation, then you know everything.

\para
{\color{olive}In classical mechanics, the state of a particle at some initial time is given by it's position ${\bf x}$ and it's velocity ${\bf v}$. The equation $F=ma$ then tells you the state at all later times.}

\para
This is the most basic fact about classical mechanics. And it's wrong! Meaning that this isn't the way our world works. Even the most basic things get overhauled in quantum mechanics. Instead, in the quantum world, the state of a particle is given by the {\it wavefunction}

\para
{\color{olive} In quantum mechanics, even this basic statement is wrong! The state of a particle at some time is given by the {\it wavefunction} $\psi(\bx)$. Importantly, this is a complex-valued function: $\psi(\bx) \in \mathbb{C}$.}  

\para
This means that the value of the wavefunction varies in space -- at every point $\bx$, you need to specify a complex number.  {\color{blue} The audience will need to know what a complex number is. Just starting a list here of things we assume. I'm happy to explain/remind what the $\mathbb{C}$ symbol means but the general idea of a complex number needs to be familiar.} 

\para
Already, this is a huge departure from the classical world. In the classical world, you need to write down six numbers to specify the state of a particle: three to give the coordinates in space, $(x,y,z)$, and three to give its velocity  -- let's say $(u,v,w)$, the speeds in the $x$, $y$ and $z$-directions. But in the quantum world, you need to specify an infinite amount of data to say what the state of a particle is -- a function's worth. You need to associate a number to every point in space. And there are an infinite number of points in space! Moreover, it has to be a complex number-- the theory doesn't work with just real numbers. Complex numbers really sit at the heart of physics. 

\para
But there's also a sense in which this quantum state contains {\it less} information than the classical state. Because we don't say anything about the velocity of the particle. That has to be encoded, somehow, in the wavefunction too. 

\para
There's a long and winding story about {\it why} the wavefunction is the right description of the quantum state, some of which we'll get to in this course. Other parts will have to wait for subsequent courses. But, for now, we're going to take it as an axiom of quantum mechanics, At any given time, this object -- the wavefunction -- is what we will want to compute. 

\para
So what is it?! Well, the answer is both simple and extremely subtle. The wavefunction is telling us the probability that the particle lives at some point. Now probabilities have to be between zero and one. {\color{blue} They need to know this much about probability!} And, most notably, probabilities are not complex numbers. So a better statement is that:

\para
{\color{olive}The probability density $P(\bx)$ for the particle to sit at point $\bx$ is given by 
%
\be P(\bx) = |\psi(\bx)|^2\ .\nn\ee}
%
The probability is the mod-squared of the wavefunction. {\color{blue} Again, this idea from complex numbers needs to be familiar.} Now there's a word "density" that snuck in here. What does that mean? The probability to be at any particular point $\bx$ is always zero. What's the probability that the $x$ coordinate is exactly $1.00123445437$. Well, the probability is zero because you have to specify the point to ridiculous accuracy. So instead, we ask: what's the probability that the particle sits in some region of space that I'll call $V$. Then:

\para
{\color{olive} Probability $P_V$ of sitting in some region $V$ is
%
\be P_V = \int_V  P(\bx) \,d^3x = \int_V |\psi(\bx)|^2 \,d^3x\ .\nn\ee
%
This is called the Born rule.}

\para
That is, you integrate the probability distribution over the region $V$ -- which is specified, tucked away at the bottom of the integral sign. {\color{blue} This needs the audience to know about volume integrals. Not necessarily how to do them, but at least what the notation means.}

\para
So we see that, right from the outset, probability is part of quantum mechanics. That uncertainty is there from the beginning. It is named after Max Born, who first understood how to think about the wavefunction. 

\para
It's worth saying a few words about this. In all other areas of physics -- and, indeed, other areas of life  -- probability is all about our ignorance. If you throw a dice then, from the moment it leaves your hand, there's nothing uncertain about what will happen. Its initial conditions are set, and from then on it's just following deterministic laws of physics. Every bounce, every spin, is predetermined by the initial conditions and by this equation here -- $F=ma$. There's nothing random about it. But, in general, we don't know the initial conditions that well, and these equations are hard to solve, we just admit our ignorance and fold that into probability -- we say that the probability it will land on any given number is 1/6. 

\para
The same is true of the weather, or our health, or anything else. Probability is summarising our ignorance. Gain more knowledge and you can say more. Gain perfect knowledge and that probability evaporates completely. 
 
\para
But that's not what's going on with quantum mechanics. This probability is real -- it's an inherent property of the world. There's no avoiding it. No secret information that, if only we could learn, would eliminate the probabilisitic nature of things. . This has been a very hard-fought lesson to take from quantum mechanics -- with arguments going back and forth over the decades. Many eminent scientists were deeply unhappy about this. Schrodinger didn't like it. Einstein didn't like it. But the wonderful thing about science is that, no matter how eminent, you don't get the last work. Nature does. And nature has told us, unequivocally, that the universe is, at heart, random. 

\para
There's much more that should be unsettling about this. I've told you that the right description of a particle is as a smeared wave of probability -- or more accurately as a smeared wave of complex square-rooted probability, whatever that means. But that doesn't tally with what we see as particles. We can take photographs of subatomic particles leaving tracks -- electrons and positrons and muons and pions -- and they never look like smeared clouds of probability, spreading jello-like throughout space, a little bit here and a little bit over there. Instead they look like, well, particles. {\color{blue}Show a cloud chamber photo here?} How to reconcile this with our fundamental description in terms of a wavefunction? Well, that too is something that we're going to have to come to terms with as this course progresses.

\para
{\color{red} My guess is that this, in total, is at least 30 minutes up to here. I've written little, but said a lot. Perhaps 30 minutes uninterrupted. I can imagine it could be double that with an interactive audience with many questions. Either way, I think this is a good place to stop the first lecture.} 

\para
{\color{blue} I'm not sure that there's anything concrete enough in this lecture to set problems, at least not before we get to the normalisation of the wavefunction. Perhaps plotting the Gaussian wavefunction and it's probability distribution to get a sense of what a localised distribution looks like in preparation for the next lecture?  I like the suggested simulation though -- it's hinting at wavefunction collapse without ever getting there. Although ultimately, it's going to start with a function (presumably necessarily real for graphical reasons) and end up with a histogram that looks like a block version of that function-squared.}





\newpage
\section{Normalisation and Superposition}

In this lecture we're going to talk more about the wavefunction, and what the implications are if we say that it captures the state of a quantum system. 

\subsection{Normalisation}

{\color{olive} \be P(\bx) = |\psi(\bx)|^2\ .\nn\ee}
We've met the wavefunction $\psi(\bx)$, and its interpretation as a probability distribution -- by taking the modulus squared. But if we're going to interpret something as a probability, then that comes with certain requirements. First, the probability should always be a real, positive number. Well, that's taken care of by taking the modulus of the wavefunction -- it's guaranteed to be real and positive. Next, the particle has to be {\it somewhere}. That means that if we ask "what's the probability that the particle is anywhere in the universe" then the answer has to "one". (Quantum mechanics is weird, but it's not nuts...there's still logic to it!). Mathematically this means

\para
{\color{olive} The particle has to exist somewhere, so if we integrate over all space 
\be \int_{\mathbb{R}^3} P(\bx) =  \int_{\mathbb{R}^3} |\psi(\bx)|^2 = 1\ .\nn\ee
%
This is called the {\it normalisation} of the wavefunction.}
%
Here there's no restriction to the integral -- that volume $V$ that was dangling below the integral sign has been replaced by $\mathbb{R}^3$, which means all of three-dimensional space. 

\para
Now, in general, this isn't such a big deal. Suppose that you have a wavefunction $\Psi(\bx)$ -- that's a capital $\psi$ -- which isn't normalised, but instead obeys:

\para
{\color{olive} Suppose that some candidate wavefunction $\Psi(\bx)$ obeys
%
\be \int_{\mathbb{R}^3} |\Psi(\bx)|^2 = N< \infty\ .\nn\ee
%
We say that $\Psi(\bx)$ is {\it normalisable}. We can then construct a normalised wavefunction simply by 
%
\be \psi(\bx) = \frac{1}{\sqrt{N}} \Psi(\bx)\ .\nn\ee} 
%
Sometimes it's useful to work with non-normalised (but normalisable) wavefunctions and do the normalisation at the end. For a wavefunction to be normalisable we must have:

\para
%
{\color{olive} Normalisable requires $\psi(\bx) \ra 0$ as $|\bx| \ra \infty$.}
\para
We could be more precise here -- worrying about how quickly the wavefunction dies out as we go off to infinity in space. But this will suffice for now. 

\para
{\color{olive} In contrast, wavefunctions that are not normalisable obey
%
\be \int_{\mathbb{R}^3} |\psi(\bx)|^2 = \infty\ .\nn\ee
%
These are not good states in quantum mechanics.}  Annoyingly, even though they're not good states it turns out that functions like this do rear their head from time to time, as we will see. 

\para
There's one last thing we need to know, which I'll again state as an axiom. If two states differ only by a constant complex phase then they're actually the same state. Recall that a {\it complex phase} is a complex number with modulus one. 

\para
{\color{olive} Comment: two states $\psi(\bx)$ and $e^{i\alpha} \psi(\bx)$ should be viewed as the same state. The probability distribution $P(x) = |\psi(\bx)|^2$ doesn't change. Neither, it turns out, does anything else we can measure. }

\para
That ``anything else" caveat is important. For example, you could consider multiplying the wavefunction by a phase that depends on the position of the particle. But, as we will soon see, that very much changes the state. It's only constant complex phases that give you the same state. 

\para
{\color{olive} Note: multiplying by, say $e^{i\alpha x}$ with $x$ the position of the particle also leaves $|\psi(\bx)|^2$ unchanged but this {\it does} change the quantum state...as we will later see.}


\subsection{Superposition}

So the state of a system is described by the wavefunction. And the wavefunction should be normalisable but, as we've seen, that's not a big deal: provided that it's normalisable then it can always be normalised. Now, that's the only restriction on states. Any normalised wavefunction is a viable state of the particle. Could be nice and sensible - a little lump of probability in some well-defined region of space -- it could be spread out and ridiculous. Anything is possible. 

\para
Now comes a big conceptual leap. Which I'll go slowly and write it down.

\para
{\color{olive} Suppose that we have two different states of a particle, $\psi_1(\bx)$ and $\psi(2)$} -- two different wavefunctions. Then I can always add them together to get another state. {\color{olive} Then the combination $\psi_3(\bx) = \alpha \psi_1(\bx) + \beta\psi_2(\bx)$ is also a normalisable wavefunction for any constant $\alpha,\beta\in \C$, so corresponds to a state. This is called the principle of superposition. You can add two states together to get a third.}


\para
Strictly, we should prove this. We would need to show that if $\psi_1(\bx)$ and $\psi_2(\bx)$ are normalisable, then so too is $\psi_3(\bx)$. It's not necessarily normalised -- but it is normalisable and so can be normalised to correspond to a state. There's a slightly clunky formal proof, but it boils down to the fact that if $\psi_1(\bx)$ and $\psi_2(\bx)$ are normalisable, then they decay away sufficiently quickly as you go off to infinity, and so too must $\psi_3(\bx)$. ({\color{blue} It requires some familiarity with the mathematics for the audience to get this argument. I wouldn't want to give the actual proof in the lectures -- it's in the book -- as I'm not sure it adds any intuition beyond this. Might we do a mini-explainer: a short 2-3 minute video going through the proof? Or maybe even that would be too much of a distraction.} \alyssa{This proof will be included in a Tutorial for Lecture 2.} \href{https://drive.google.com/file/d/15D_utjt-mdNEDmhK1WxOPFAAx-HqOa1U/view?usp=drive_link}{[Lecture 2, Tutorial 1]}


\para
The principle of superposition is where a lot of the magic of quantum mechanics happens. Let's suppose that I have a state of a particle that looks very much like what a particle should look like -- which, in quantum mechanics, means some lump of probability that's localised in some small region of space. That would correspond to a wavefunction something like this:

\para
{\color{olive} The wavefunction 
%
\be \psi(\bx) = \frac{1}{\sqrt{N}} e^{-a(\bx-{\bf X})^2}\ .\nn\ee
%
describes a particle localised near the point ${\bf X}$.} {\color{blue} Should have a picture on screen, or perhaps I just draw it} {\color{olive} For what it's worth, we need $N=(\pi/2a)^{3/2}$ for the wavefunction to be normalised. But we could also consider a wavefunction for single particle given by
%
\be \psi(\bx) = \frac{1}{\sqrt{N'}} \left(e^{-a(\bx_1-{\bf X_1})^2} + e^{-a(\bx_2-{\bf X_2})^2}\right)\ .\nn\ee
%
This describes a particle that sits partly at ${\bf X}_1$ and partly at ${\bf X}_2$. The particle is in two places at the same time!}

\para
This isn't just a glitch of the mathematics. We can now routinely do experiments where particles -- typically single atoms -- are placed in a superposition and sit in two places at the same time. We usually use rubidium or caesium atoms, because they're particularly amenable to being nudged by lasers, but it's quite possible to put a single atom in two places at the same time. The current record is roughly a centimetre apart. Which, if you're an atom, is a long way.  We then do even stranger things, like get the atom to interact with its ghostly partner. That's what's called, with mild understatement, interference. All of this is possible because states in quantum mechanics are much more subtle than in the classical world - they're wavefunctions.

\para
You may also have heard of quantum computers . Those two derive their power from, among other things, the principle of superposition. A qubit doesn't have to be 0 or 1 -- it can be a little bit 0 and a little bit 1 and that gives it an advantage of boring classical bits, allowing certain computations to be done exponentially faster. It's all boils down to the principle of superposition. 


\para
There is an important mathematical implication of this. States in quantum mechanics are examples of vectors. Now that might sound odd. You might think that vectors are things that have a length and a direction. That's certainly what we learn in school. But a mathematician's definition of a vector is that it's a class of objects where you can multiply them by numbers, or add them together, or multiply them by numbers and then add them together, and you get the same kind of object. That's what mathematicians call a vector space. So you can do this will actual vectors -- add two vectors together and you get a third. Multiply a vector by a number and you get another vector. You can also do this with just numbers -- that's not surprising: a number is just a one-dimensional vector. But here we're learning that functions -- and in particular, normalisable functions -- are examples of vectors. You can add two normalisable functions and you get a third. Or, said physically, you can add two states in quantum mechanics and you get a third. The novelty is that functions are examples of infinite dimensional vectors! 

\para
As we proceed through this course, we're going to see that the relationship between functions and vectors really sits at the heart of what quantum mechanics is about. Mathematically, at least, it will probably be the big conceptual leap. 

\para
{\color{red} I'm writing more in this lecture than the last one. If there's a blackboard, I suspect that this will fill 30 minutes. It will be faster on powerpoint. Still, with an interactive audience I doubt we'll get much change out of half an hour.}


{\color{blue}\subsection{Problems} I like the idea of asking them to normalise various wavefunctions. We could do two quadratics on an interval, but one with a maximum at the origin and the other with a minimum at the origin, so the particle prefers to be at the boundary. Hopefully it is more revealing than it is messy.

\para
Perhaps a deceptively simple question: we could consider a particle on a circle with wavefunction $\psi(x) = e^{i nx/R }$ with $n\in \Z$. First get them to normalise. Then get them to ask what's the probability that the particle sits at some particular point on the circle -- of course, it's the uniform probability distribution for all values of $n$. So different quantum states can have the same probability distribution} \alyssa{Currently, this is set as the practice problem for Lecture 2.} \href{https://drive.google.com/file/d/1Q7qUD2UxkJ8PmS6yvz7MZLnbiHyQ6ZD6/view?usp=drive_link}{[Lecture 2 Practice Problem]}

\newpage
\section{The Schr\"{o}dinger Equation}

So far we've talked about the state in quantum mechanics -- that's the wavefunction -- but we haven't talked about how it depends on time. That's like saying that the state of a particle in classical mechanics is described by position and velocity but not giving any hint of how it changes. In other words, we need to the quantum version of Newton's $F=ma$. And the quantum version of $F=ma$ is the Schr\"{o}dinger equation. 


\para
{\color{olive} The wavefunction evolves in time, so we should write $\psi(\bx,t)$. It is governed by the Schr\"{o}dinger equation 
\be i\hbar\ppp{\psi}{t} = \hat{H}\psi\ .\nn\ee
%
Here $\hbar$ is Planck's constant. It takes the value
%
\be \hbar&\approx& 1.05 \times 10^{-34} \ {\rm J\,s} \nn\\ &\approx& 6.68 \times 10^{-16}\ {\rm eV\,s}\ .\ee}
%
You may have heard that this is Planck's constant divided by $2\pi$ -- but it's better to use $\hbar$ than $h/2\pi$. hbar is the characteristic signature of quantum mechanics -- see it in any equation and it's jumping up and down and telling you that quantum mechanics is at play. The value of $\hbar$ is important for any applications, but conceptually we need to remember the units, or dimensions. It has units of energy-times-time. That's the same as the units of angular momentum. That will be important. 
%

\para
{\color{olive} $\hbar$ has units energy $\times$ time, the same as angular momentum}

\para
The Schr\"{o}dinger equation is a differential equation. It's telling us how the wavefunction changes with time. And the meat of it is this capital $H$ with a strange hat on. That's called the Hamiltonian, and it captures the kind of force that the particle experiences

\para
{\color{olive} $\hat{H}$ is the {\it Hamiltonian}. It is given by
%
\be \hat{H} = -\frac{\hbar^2}{2m} \nabla^2 + V(\bx)\nn\ee
%
where 
%
\be \nabla^2 = \ppp{^2}{x^2} + \ppp{^2}{y^2} + \ppp{^2}{z^2}\ee
%
is the {\it Laplacian} and $V(\bx)$ is the potential energy.}

\para
So there's a lot to unpack here. What's going on. First, the Hamiltonian is an example of what mathematicians call an {\it operator}. That's what the hat means and we're going to retain that notation throughout these lectures. Anything that has a hat is an operator. That means that it's just waiting for a function to come along, from the right, so it can do something to it. In this case, the Hamiltonian does two things. First it differentiates the function -- the wavefunction $\psi$ in this case. And it differentiates it with respect to the spatial coordinates $x$, $y$, and $z$, as dictated by this upside-down triangle. The upside down triangle is the hallmark of vector calculus -- how you differentiate functions in three-dimensions and we'll meet it again later, but for now you can just take this as the definition of upside-down-triangle-squared. It's called the Laplacian. So the the first thing the Hamiltonian does is it differentiates the wavefunction, twice with respect to $x$, twice with respect to $y$, and twice with respect to $z$ and then adds them all up. 

\para
The second thing the Hamiltonian does is multiply by this function $V(\bx)$. That function is the potential energy and that's what you have to input into the Schr\"{o}dinger equation to say what kind of problem you're dealing with. It could the potential energy for a sort of three-dimensional spring, or the potential energy for an electron moving under the electrostatic attraction of a proton
%

\para
{\color{olive} Examples of potential energy include the 3d harmonic oscillator
%
\be V(\bx) = -k r\nn\ee
%
where $r = \sqrt{x^2 + y^2 + z^2}$ and $k$ is the spring constant. Or the Coulomb energy for an electron orbiting a proton
%
\be V(\bx) = -\frac{1}{4\pi \epsilon_0} \frac{1}{r}\ .\nn\ee}
%
Here $\epsilon_0$ is a number that captures the strength of the electrostatic force. It's got a pretentious name -- permittivity of free space but it's just a constant of nature that tells you how strong electric forces are. For each difference choice of the potential energy, we have to solve the Schr\"{o}dinger equation afresh. In some sense, that's the whole purpose of these lectures. To solve this equation in different situations and understand what it's telling us. 

\para
There's already something interesting here. In classical mechanics, we have to solve Newton's second law $F=ma$. And there you first have to decide on what force you're dealing with. And there are lots of different kinds of forces -- gravitational, electrostatic, spring-things. And also things like friction. In fact, in many every day situations friction is one of the most dominant forces. 

\para
Now, some of those forces are fundamental. And some, like friction, are kind of fake forces. Perhaps emergent force is a better phrase. There's no friction on a fundamental level. Instead, it comes about because the particle you're interested in is interacting in some complicated way with some other material that it's rubbing against, and exciting the vibrations of the atom in that material. That means that, if you ignore all those small atomic vibrations, when you have friction energy does not seem to be conserved.

\para
From the perspective of $F=ma$, the equation couldn't care less about where the force came from -- conserves energy, doesn't conserve energy -- it really doesn't matter. Plug anything in and off you go.

\para
But you can only write down the Schr\"{o}dinger equation for forces that conserve energy! You can't write down the Schrodinger equation for a friction force -- if you want to do it, you somehow need to include all those other gazillions of atoms that cause the friction in the first place. The entire framework of quantum mechanics pushes you to consider only certain very special types of forces. Which, it turns out, are those forces that hold at a fundamental level. Forces that conserve energy. 

\para
{\color{olive} The force is given in terms of the potential energy by
%
\be {\bf F} = -\nabla V = -\left(\begin{array}{c}  \partial V/\partial x \\ \partial V/\partial y  \\ \partial V/\partial z \end{array}\right)\ .\nn\ee
%
This is a {\it conservative force}.}
%
There's the upside-down triangle again. It's definition is given on the right -- it's a way of differentiating with respect to $x$, $y$ and $z$ to turn a function like the potential energy into a vector like the force. Here it's called the gradient of the function. 


\para
Forces that can be written like this are called conservative forces -- precisely because they conserve energy. There's one other kind of force that we can write down the Schrodinger equation for, which is the magnetic force. We won't do it in these lectures, but it is possible. 

\para
{\color{olive} In classical mechanics, the conserved energy is
%
\be E = \frac{1}{2}m{\bf v}^2 + V(\bx) = \frac{1}{2m} {\bf p}^2 + V(\bx)\nn\ee
%
where ${\bf p} = m{\bf v}$ is the momentum}. 

\para
Now this looks suspiciously like the Hamiltonian. Something divided by $2m$ plus the potential energy. In fact, the Hamlitonian is really a quantity that first appears in classical mechanics -- it's named after the Irish mathematician William Rowan Hamiltonan -- that's very closely related to the energy. And you should think of this strange thing called the Hamiltonian as the energy, and we'll see in more detail exactly what that means as we go along.

\para
{\color{olive} In quantum mechanics, the Hamiltonian $\hat{H}$ is related to the energy. But that suggests that momentum is somehow related to differentiation!
\be {\bf p} = \pm i \hbar \nabla \ ?!\nn\ee
%
What could that mean? We will see as we go along.}

\para
{\color{blue} I like the idea of having a simulation here. We can take the 1d harmonic oscillator and first show a classical ball rolling back and forth. And then solve the Schrodinger equation using a coherent state, which rolls back and forth in the same way. That gives some nice intuition for what's going on. But then we should do the same thing for an anharmonic oscillator with $V(x) = x^2 + \epsilon x^4$ for some small choice of $\epsilon\approx 0.1$ or something. Now the classical ball will roll back and forth in a way that should be indistinguishable. But I'm pretty sure that the quantum wavefunction will break up and do funky things -- would be nice to see.}

\subsection{Conservation of Probability}

Before we solve any specific problem, there's a general property of the Schrodinger equation that we can look at. Because the wavefunction -- or at least the wavefunction-mod-squared -- has the interpretation as a probability. And probabilities have to sum to one. Which means that the way the wavefunction evolves has to be such that this probability interpretation continues. The sum -- or really the integral -- of the probability has to be one for all times. That's something that we can check. 

\para
To keep things simple, I'll do this for a particle moving in just one dimension, rather than in three dimensions. {\color{blue} The difficulty in doing it in three-dimensions is that you need to invoke the divergence theorem at the very end. I think it's best to duck this issue here. Maybe it's worth a mini-explainer going through the same derivation in three-dimension?} \alyssa{This derivation is included in the Tutorial for Lecture 3.} \href{https://drive.google.com/file/d/12JSixdQ0R_mRoL3O6Sr0HWOza6faMNb5/view?usp=drive_link}{[Lecture 3, Tutorial 1]}

\para
{\color{olive} The probability density is $P=|\psi|^2 = \psi^\star\psi$. We have
%
\be \ppp{P}{t} = \ppp{\psi^\star}{t}\psi + \psi^\star\ppp{\psi}{t}\ .\nn\ee
%
For a particle moving in one-dimension, we have $\psi(x,t)$ and the Schrodinger equation is
%
\be i\hbar\ppp{\psi}{t} = -\frac{\hbar^2}{2m} \ppp{^2\psi}{x^2} + V(x)\psi\ .\ee
%
The complex conjugate is
%
\be -i\hbar\ppp{\psi^\star}{t} = -\frac{\hbar^2}{2m} \ppp{^2\psi^\star}{x^2} + V(x)\psi^\star\ .\ee}
%
Substitute in and you'll find that those $V(x)$ terms cancel. We're left with
%
{\color{olive} 
\be \ppp{P}{t} &=& \frac{i\hbar}{2m} \left(\psi^\star \ppp{^2\psi}{x^2} - \psi \ppp{^2\psi^\star}{x^2}\right) 
\nn\\ &=& \frac{i\hbar}{2m} \ppp{}{x}\left(\psi^\star \ppp{\psi}{x}- \psi \ppp{\psi^\star}{x}\right)
\ .\ee
%
So we have
%
\be \ppp{P}{t} + \ppp{J}{x} = 0\nn\ee
%
where 
%
\be J = -\frac{i\hbar}{2m} \left(\psi^\star \ppp{\psi}{x}- \psi \ppp{\psi^\star}{x}\right)\ .\ee
%
is the probability current.}

\para
So this calculation doesn't say that the probability density $P$ doesn't change with time -- that would be too restrictive, because it can flow from one place to another. Instead, it says that if $P$ does change with time then it has to follow this equation. Equations of this form appear all over the place in physics. It's called the {\it continuity equation}. It's the equation that says that something is conserved {\it locally}. If the probability changes in some region of space, then can't just reappear miles away. Instead, it must leak into a neighbouring region of space. It's the current $J$ that captures how probability flows. A continuity equation appears anytime in physics that we have some conserved quantity -- electric charge, or energy -- because the equation captures the combination of conservation and locality. 

\para
{\color{olive} Consider the probability in some region $-L \leq x\leq L$. 
%
\be P_L = \int_{-L}^{+L} P(x,t) \,dx\ .\nn\ee
%
This changes with time as
%
\be \ppp{P_L}{t} = \int_{-L}^{+L} \ppp{P}{t} \,dx = -\int_{-L}^{+L} \ppp{J}{x} \,dx = -\left[J(L,t) - J(-L,t)\right]\ .\nn\ee 
}
That's the key result -- because we have a total derivative here, the change in probability in a region is only because there's a current leaking out of the region -- at $x=L$ and $x=-L$ in this case. 

\para
{\color{olive} In particular, when $L\rightarrow \infty$, $P_L\ra 1$. Now $\partial P_L/\partial t=0$ because $J\ra 0$ as $x\ra \infty$. So the total probability is conserved: it always integrates to one.}

\para
{\color{red} I suspect that this is somewhat over 30 minutes. It's probably under 30 minutes without the conservation of probability. Depending on the simulation, it might be best to push the conservation of probability calculation to he next lecture.}

{\color{blue} 
\subsection{Problems}
Show that the units of angular momentum are energy times time. \alyssa{This will be the practice problem for Lecture 3.} \href{https://drive.google.com/file/d/1Qu2bBqQa2FQBaBoZ9YU335qluru4K50T/view?usp=drive_link}{[Lecture 3 Problem Statement]}

\para
We could get them to re-run the conservation of probability calculation in three-dimensions, although they'll need to know the divergence theorem to get to the end. \alyssa{This will be covered in a Tutorial for Lecture 3.} \href{https://drive.google.com/file/d/12JSixdQ0R_mRoL3O6Sr0HWOza6faMNb5/view?usp=drive_link}{[Lecture 3, Tutorial 1]}
}


\subsection{Tangent: Erwin Schrodinger}

{\color{blue}Such a lovely man. I wonder if there's a different way to spin this, rather than focussing on the more unsavoury parts of his life. Schrodinger was playing catch up when he developed his equation. The key insight came to Heisenberg in the summer of 1925, sitting on Helgoland. But -- and I find this completely remarkable -- he had no idea what a matrix is. He'd followed his logic to find that quantum mechanics needs objects for which $xp\neq px$ and he even derived the rule for multiplying together infinite-dimensional matrices. But he'd never heard of matrices. Happily, back in Gottingen, his postdoc advisor Max Born knew what they were. So Born, Jordan (another lovely man) and Heisenberg essentially put together what we now see as quantum mechanics.

\para
As far as I can tell, if Schrodinger was motivated by this success at all, it was because he too didn't know what a matrix was and had no intention of learning. He was ancient (38 I think!) and didn't want to mess around with any fancy new math. He wrapped this up in philosophically minded ramblings about a lack of visualisation but, honestly, he just didn't want to embrace the new. 

\para
He gave a seminar on de Broglie's ideas which were rightly criticised for being too naive. If you think something is a wave, it's not enough to say what it's wavelength is: you've got to say how it waves. How it moves. That's what drove him to his year of astonishing creativity, writing down his equation and reproducing the work that the Gottingen gang had done using operators. Maybe this is a more fun story to tell?}

\subsection{Explainer: Hamiltonian Mechanics}

{\color{blue}I think we could certainly do with a refresher on why having a conservative force of the form ${\bf F} = -\nabla V$ means that energy $E={\bf p}^2/2m + V(\bx)$ is conserved. This means the calculation 
%
\be \frac{dE}{dt} = \frac{1}{m} {\bf p}\cdot \dot{\bf p}+ \nabla V \cdot \dot{\bx}\ .\ee
%
Then use Newton's second law to write $\dot{\bf p} = m{\bf a} = {\bf F} = -\nabla V$ to get $\dot{E} = 0$. But it does need some familiarity with the chain rule, including using $\nabla$ when differentiating $V(\bx)$ so it might be worth spending some time explaining what's going on -- $\bx$ is really $\bx(t)$ and so when we differentiate with respect to time, blah, blah.

\para
One possibility is to include the brief calculation above in the lecture and have an explainer that goes more slowly through the partial differentiation side of things. 


\para
All of this seems more urgent that explaining Hamilton's equations  in an explainer -- i.e.  thinking of phase space and splitting Newton's equations second order equation into two first order equations. We could include an explainer on this too, although I suspect that it would all look a bit trivial -- there wouldn't be any subtle aspects of Hamiltonian mechanics where the momentum is anything other than $p=mv$. (I think introducing a magnetic field is the first time such subtleties arise.).}

\newpage
\section{Collapse and the Double Slit Experiment}

I want to do two things in this lecture, both of them rather famous aspects of quantum mechanics. First, I want to tell you about the collapse of the wavefunction. Then I want to tell you about the double slit experiment.


\subsection{Collapse of the Wavefunction}

The collapse of the wavefunction is one of the more unsettling aspects of quantum mechanics. The wavefunction evolves through the Schrodinger equation. Until it doesn't. Until, that is, you look at it. Or, in the more usual words, until an observation is made. Then the probabilistic side of quantum mechanics comes to the fore.

\para
{\color{olive} Suppose that you have a particle in a superposition, like this
%
\be \psi(\bx) = \frac{1}{\sqrt{N'}}\left(e^{-a(\bx - {\bf X}_1)^2} + e^{-a(\bx - {\bf X}_2)^2}\right)\ .\nn\ee}
%
The particle is sitting in a ghostly superposition of ${\bf X}_1$ and ${\bf X}_2$. As I mentioned before, this isn't a statement about our ignorance -- it's not that it's somewhere and we just don't know where. It really is in two places at the same time. We can do experiments -- in fact related to the double slit experiment that we'll soon see, to demonstrate that it's in two places at the same time. But that begs the question. Why don't we just look? Put a detecting apparatus at ${\bf X}_1$ and another at ${\bf X}_2$ and just see where it is. What happens?

\para
{\color{olive} Suppose that we detect the particle to be in the vicinity of ${\bf X}_1$. The the wavefunction changes discontinuously to become something like
%
\be \psi(\bx) = \frac{1}{\sqrt{N}} e^{-a(\bx - {\bf X}_1)^2}\ .\nn\ee
This is the collapse of the wavefunction.}

\para
The wavefunction no longer has any support near the other position ${\bf X}_2$. It can't have because there's no chance the particle is over there. Yet  we know it's right here. 

\para
The collapse of the wavefunction is one of the more disconcerting aspects of quantum mechanics. Most of the time, the wavefunction evolves through the Schrodinger equation. There's nothing random about this, nothing probabilisitic. The Schrodinger equation is as deterministic as anything in classical mechanics. But then you choose to look -- to make an observation. And this is where the probabilities rear their heads. The wavefunction collapses -- of all the potentialities in the wavefunction, it picks one, randomly, determined by probability. And the wavefunction changes, abruptly, discontinuously, to reflect that.

\para
In a previous lecture I told that you the wavefunction described the inherent state of the system -- not our knowledge about it. But it's also true, that gaining knowledge about the particle affects it. It changes the state. So one lesson to take from this is that you can't observe a system without changing it. In looking, you change the state of the system.You might think that you've gained knowledge and reduced your uncertainty about where the particle is but, as we will see, the wavefunction contains other information beyond the position of the particle and you've necessarily upset that as well.

\para
The fact that the wavefunction can evolve in two different ways -- smoothly by the Schrodinger equation, and abruptly by measuring -- has made quite a lot of people upset. They don't like it. The more philosophically minded you are, the more upset you tend to be. People say words like ontic and epistemic and get really very annoyed. We'll talk about some of these issues -- what's called interpretations of quantum mechanics -- later when we have more under our belt. But, for now, I'm just going to illustrate some of this with a simple experiment. Perhaps the most famous experiment in quantum mechanics. This is the double slit.

\subsection{The Double Slit Experiment}

{\color{blue} No script for this as it would be intertwined too much with the demo. But we should definitely do a demo! First the classical case with balls. Obviously they only go through one slit or the other. Then the classical case for a wave. Now it matters whether a single slit is open or both slits are open, but that's ok because the wave is an extended object -- it literally goes through both slits at the same time. This would be a good point to change some parameters, as suggested in the powerpoint slide: slit width, slit separation, wavelength.

\para
Then the big surprise. The quantum particles is a combination of the two -- we see diffraction, but we also only detect a particle in particular position. It's described by the wavefunction until we look. If we look after it's passed, then we get collapse of the wavefunction and the diffraction pattern only builds up over time, probabilistically. If we look as the particle goes through a slit, then we collapse the wavefunction earlier and the diffraction pattern disappears.

\para
There's a more precise story here which I didn't go into in the book. Look more closely at the particle by sending in shorter wavelength light and you disturb it more. This needs the Heisenberg uncertainty principle to be done properly. But perhaps we could give some hint of this here as a prelude to the uncertainty principle that comes later.}

{\color{red} This lecture is probably less than 30 minutes, although it depends how long we dwell on the double slit experiment.  That might mean that we move the conservation of probability to this lecture. Either way, combined they probably add up to about an hour.} 

\subsection{Problems}

{\color{blue} This is a very conceptual lecture and I can't think of any problems. There's a suggestion to do the classical interference pattern. The only way I know how to do this is Frauenhofer something or other and was quite involved, thinking about how phases of waves differ along different paths. Such a problem would need quite a bit of handholding. It might be better as an explainer. } \alyssa{Classical interference will be included in the Tutorial for Lecture 4.} \href{https://drive.google.com/file/d/1SRmv_wbY5g6IYCpOA7z6lYqfwOQ5H9WR/view?usp=drive_link}{[Lecture 4, Tutorial 1]}


\newpage
\noindent
{\Large \bf Section 2: A Particle on a Line}


\section{A Particle in 1d}

We've got the Schrodinger equation, which is the foundational equation of quantum mechanics. And we now have to do two things. The first is to solve it in various situations. And the second is to understand what on earth it's telling us about the world we live in. How to interpret it. What do the solutions mean? 

\para
And there's really two ways that we could approach this. A logical way would be to list all the axioms of quantum mechanics up front -- tell you about Hilbert spaces and eigenstates and so on. Or we could take a pedagogical approach, where we go a bit slower and start to first solve some simple examples and build some intuition for the Schrodinger equation, before we lay everything  on the table and explain what's going on. So I'm going to take that second approach. We'll spend a few lectures solving the Schrodinger equation in some simple examples, and trying to piece together what it means. That will require us to take some leaps of faith in places, but we'll then backtrack later and fill in the gaps.

\para
To this end, we're going to look at a very restrictive class of examples. We'll take a single quantum particle -- just one -- and allow it move on a line. So it doesn't get to roam in 3d space. It's restricted just to live in one dimension.

\para
{\color{olive}A classical particle moving in 1d is described by its position $x(t)$ which obeys Newton's second law
\be m\ddot{x} = -\frac{dV}{dx}\nn\ee
%
with $V(x)$ the potential energy.} {\color{blue} Maybe a very quick simulation here to explain that you plot the potential and the particle rolls backwards and forwards like a ball on a hill. But only if it's super quick and doesn't get in the way.}

\para
{\color{olive}A quantum particle is described by a wavefunction $\psi(x,t)$ which obeys the Schrodinger equation
%
\be i\hbar \ppp{\psi}{t} = -\frac{\hbar^2}{2m}\ppp{^2\psi}{x^2} + V(x)\psi\ .\nn\ee
%
We want to solve this.}

\para
So what are we going to do. This is an example of a partial differential equation, or a PDE -- it means that you're differentiating with respect to more than one variable...time here, and space here and it first it looks quite hard. But there's a straightforward way to make it easier.

\para
{\color{olive} We look for solutions that take the form
%
\be \psi(x,t) = e^{-i\omega t} \psi(x)\nn\ee
%
for some choice of {\it frequency} $\omega$.} Note that our notation might be a little confusing as I've given the wavefunction $\psi(x,t)$ and $\psi(x)$ the same name. I could -- perhaps should -- have given this guy a different name, but $\psi$ is such a common name for the wavefunction and, from now on, it's just this guy $\psi(x)$ that we'll work with. This means that the spatial part of the wavefunction doesn't change with time, but it comes with an overall complex phase $e^{-i\omega t}$ -- and that phase is rotating with time with a frequency $\omega$. It circles around once every $1/\omega$.

\para
{\color{olive} The general Schrodinger equation takes the form
%
\be i\hbar \ppp{\psi}{t} = \hat{H}\psi\ .\nn\ee
%
This is sometimes called the {\it time-dependent Schrodinger equation}. Substitute in the ansatz above} ...do you guys know the word ansatz? It's German for guess... {\color{olive} and this becomes
%
\be \hat{H}\psi = E \psi\nn\ee
%
where $E=\hbar\omega$. This is the {\it time-independent Schrodinger equation}. Written in full, it is
%
\be -\frac{\hbar^2}{2m} \frac{d^2\psi}{dx^2} + V(x) \psi = E\psi\ .\nn\ee
%
We need to solve this for both $\psi(x)$ and for $E$. We will see that $E$ has the intepretation of the energy of the particle.}


\para
If you know a little about matrices and eigenvalues, then this form of the Schrodinger equation looks very much like an eigenvalue equation. We're going to come back to that later. And if you don't know about eigenvalues then there'll be an explainer later. \alyssa{Finding eigenvalues will be included in the Tutorial for Lecture 5.} \href{https://drive.google.com/file/d/1TrNUJWBwGnIj0qXJB3SdRJRH1n03Gh1Q/view?usp=drive_link}{[Lecture 5, Tutorial 1]}

\para
You might worry that we've made this specific choice $\psi(x,t) = e^{-i\omega t} \psi(x)$ and this is too restrictive. That perhaps we're missing some other solutions. It will turn out that if we can find these solutions, they give us everything by simply adding together different solutions of this form. There's nothing more. 

\para
{\color{olive} Solutions of the form $\psi(x,t) = e^{-i\omega t}\psi(x)$ are called {\it stationary states} or {\it energy eigenstates}. }


\subsection{The Free Particle}

{\color{olive} The simplest system is a free particle with $V(x)=0$.} This is a particle that doesn't experience a force. It's simple but, as we will see, it's also annoyingly subtle. In fact, it has subtleties that don't appear in more complicated problems, essentially because there's the possibility that the wavefunction runs off to infinity. Nonetheless, we can solve it. 

\para
{\color{olive} We have
\be -\frac{\hbar^2}{2m} \frac{d^2\psi}{dx^2} = E\psi\nn\ee
%
The solutions are
%
\be \psi(x) = e^{ikx} \nn\ee
%
for any choice of $k$ and for $E=\hbar^2 k^2/2m$. These are called {\it plane wave} states -- they're waves spread out everywhere in space.} Differentiate this twice and you'll pull down two powers of $ik$ and you'll see that this works. You need that factor of $i$ in the exponent because there's a minus sign here in the Schrodinger equation.

\para
Now comes a subtlety. The wavefunction is complex. So this is a solution whether $k$ is real or imaginary. Mathematically it solves the Schrodinger equation for any $k$. But suppose that $k$ is imaginary -- so it's $i$ times some real number. Then the factor of $i$ cancels in the exponent and this is the exponential of some number times x. But that's going to blow up at either $x\ra +\infty$ (if that number is positive) or $x\ra -\infty$ (if that number is negative).  These are non-normalisable 

\para
{\color{olive} We will only consider solutions with $k\in \R$. Even these states aren't normalisable because
%
\be \int_{-\infty}^{+\infty} |\psi|^2\,dx = \int_{-\infty}^{+\infty} dx = \infty\ .\nn\ee
%
} Strictly, this means that we shouldn't consider these states at all! This is exactly the annoying subtlety that I mentioned above.  {\color{olive}This is an annoying subtlety. We'll duck this issue for now and come back to it later.

\para
Classically, a free particle has energy $E= \ft12 mv^2 = p^2/2m$. Compare to the result above suggests that the momentum is
%
\be p = \hbar k\ .\nn\ee
%
Said differently, our solution is
%
\be \psi(x) = \cos(kx) + i\sin(kx)\nn\ee
%
These have wavelength
%
\be \lambda = \frac{2\pi}{|k|}\ .\nn\ee
%
So the momentum is related to the wavelength of the wavefunction
%
\be |p| = \frac{2\pi\hbar}{\lambda}\ .\nn\ee
%
$\lambda$ is called the de Broglie wavelength.}  This was an idea that predates the Schrodinger equation -- in fact it was the idea that inspired Schrodinger's work. That a particle with momentum $p$ has an associated wavelength given by this formula. Here we see that this is a simple result of the Schrodinger equation. And it's good intuition to have. A wavefunction that oscillates with some wavelength $\lambda$ will correspond to a particle with momentum $p$ given by this formula. 

\subsection{A Particle on a Circle}

We're going to continue to look at the free particle, but we'll change the rules of the game slightly. Suppose that the particle lives on a circle of radius $R$. {\color{blue} simple picture here}. That means that if it moves, at some point it comes back to its starting point.  

\para
{\color{olive} For a free particle on a circle of radius $R$, the wavefunction should be a periodic function
%
\be \psi(x+ 2\pi R) = \psi(x)\ .\nn\ee
} 
This is really the statement that the wavefunction should be single-valued. The solutions are the same as before, but not we get this extra condtion.

\para
{\color{olive} We require 
%
\be e^{ik(x+2\pi R)} = e^{ikx} \ \ \ \Longrightarrow \ \ \ e^{2\pi ikR} = 1 \ \ \ \Longrightarrow \ \ \ k = \frac{n}{R} \nn\ee
%
with $n\in \Z$.} This fancy $\Z$ means the integers, both positive and negative i.e. $0,\pm 1,\pm 2$.

\para
This is the first time that we see the ``quantum'' in quantum mechanics. Something has to be discrete. Notice that we didn't put discreteness into our equations at all! It's only when we solve them that we find that things have to be quantised -- come in bundles. 

\para
{\color{olive} The momentum is quantised
%
\be p = \hbar k = \frac{\hbar n}{R}\ .\nn\ee
%
The energy is also quantised
%
\be E = \frac{\hbar^2 k^2}{2m} = \frac{\hbar^2n^2}{2mR^2}\nn\ee
%
with $n\in \Z$.} The lowest energy state is $n=0$. This is just a constant wavefunction, spread everywhere around the circle. That corresponds to a particle that doesn't move -- it has momentum $p=0$ and energy $E=0$. If you want to make the particle move then you need to inject a minimum amount of energy. The slowest it can move is $n=\pm 1$ with  momentum $p=\hbar/R$ which needs energy $E = \hbar^2/2m R^2$. The smaller the circle, the bigger the energy needed to make it move. If the circle is big enough, the gap between different allowed momenta becomes tiny. At some point, this becomes experimentally indistinguishable from a continuum of momenta.

\para
{\color{olive} The collection of all possible energies is called the {\it spectrum} of the Hamiltonian.}


\para
There's a simple intuitive way to see how the quantisation arises. {\color{olive} Note: the only allowed states are those for which an integer number of de Broglie wavelengths that can fit around the circle
%
\be \lambda = \frac{2\pi R}{|n|} \ \ \ n\in \Z\nn\ee
}
We can see this if we plot the wavefunctions around the circle. It's actually simpler to visualise if we plot just the real part of the wavefunction. {\color{blue} Show some simple wavefunctions...the higher $n$ ones are much easier to visualise.}

\para
For the particle on the circle, we no longer have any problem with normalisability. The wavefunction isn't normalised, but it is normalisable. {\color{olive} We have
%
\be \int_0^{2\pi R} |\psi|^2 \,dx = 2\pi R\ .\nn\ee
%
So the normalised wavefunction is
%
\be \psi(x) = \frac{e^{ikx}}{\sqrt{2\pi R}}\ .\ee
%
But, for any state} (regardless of the choice of $k$ {\color{olive} the probability distribution around the circle is
%
\be P(x) = |\psi(x)|^2 = \frac{1}{2\pi R}\ .\nn\ee
} 
So the probability distribution doesn't distinguish the different states... they have different momenta and different energies but, as far as the probability is concerned, each of them is equally smeared over the whole circle.

\para
{\color{red} I think this is a long lecture -- certainly longer than 30 minutes. Maybe closer to 40, perhaps more. As always, I have big error bars on these estimates.  But I think that the calculations reinforce one another well, so it makes a nice coherent story.}


\subsection{Problem: A Particle in a Box}
\alyssa{This entire subsection has been written into a practice problem for Lecture 5.} \href{https://drive.google.com/file/d/1KrovFFhp2vBc-8qHwI2lAaeIfoWElGgC/view?usp=drive_link}{[Lecture 5 Practice Problem]}

{\color{blue} The particle in a box is really covering the same conceptual ground as the particle on a circle. So one option --- which I think is my preferred option -- is to turn this whole section into one long problem, albeit with a lot of handholding. Nonetheless, I've written it here as a lecture.}

\para
There's a very similar story if we put our 1d particle in a box with infinite walls that it can never escape from. We do this by considering the following potential
%
\para
{\color{olive} Put the particle in a box of width $L$ with potential
%
\be V(x) = \left\{\begin{array}{cc} 0 \ \ \ & 0<x<L\\ \infty \ \ \ &{\rm otherwise}\end{array}\right.\nn\ee
%
}That infinite potential  might seem like overkill but, as we'll see later, quantum particles are slippery and we need to work hard to trap them. That infinite potential is just our way of saying that there's no way in hell the particle is getting out of the box. {\color{olive} Inside the box, we want to solve the free Schrodinger equation
%
\be -\frac{\hbar^2}{2m} \frac{d^2\psi}{dx^2} = E\psi\nn\ee
%
but now with the added requirement that $\psi$ vanishes outside the box and, in particular, 
%
\be \psi(x=0) = \psi(x=L)  = 0 \ .\ee
%
}Again, the solutions are the same cos + i-times-sin wavefunctions. But they can't quite be that because the wavefunction has to vanish at the end. And if sine vanishes then cos doesn't and vice versa. So in fact the solutions that do the trick are: {\color{olive} The solutions are
%
\be \psi(x) = \sqrt{\frac{2}{L}}\sin kx\  \ \ \ {\rm with} \ \ \ k=\frac{\pi n}{L} \nn\ee
%
where $n\in Z$.} This time I've taken the liberty of normalising the wavefuntion for you -- that's what the $\sqrt{2/L}$ is doing. {\color{blue} If we're setting this as a problem then get them to do this.} Again, we see that the momentum, and hence the energy, is quantised. You can check that the wavefunction vanishes at $x=0$ but at $x=L$ it's given by $\sin(n\pi)$ which also vanishes. 


\subsection{Explainer, Simulations, Tangents}

It was suggested to put eigensystems as an explainer. I think that can probably wait until later. But something on solving separable equations would be useful. The manipulations that are done  in this lecture really just need a fluency with differentiation, nothing more. I don't know if we want to talk about uniqueness of solutions to differential equations and such like, or leave that to the next lecture where we highlight the linearity of the equations.

\para
A simulation of the wavefunctions for a particle on a circle would be good. We should plot the real part of $\psi(x)$. But then put the time-dependence back in and plot the real part of $\psi(x,t)$. 

\para
And, yes, a tangent on standing waves could work well here.





\newpage
\section{The Gaussian Wavepacket}

In the last lecture, we looked at the free particle on a line. This should be the simplest possible quantum mechanical system because it's a particle with no force acting on it, yet we found that it was annoyingly subtle. The energy eigenstate takes the form $\psi(x) = e^{ikx}$ and is non-normalisable, so can't have any probability interpretation. What are we supposed to do?

\para
In this lecture, we're going to look more closely at the simple free particle. We'll look for different solutions to the Schrodinger equation which do have a probability interpretation. And, along the way, we'll learn more about how to interpret the wavefunction and how to extract some of its properties. These are all things that we will formalize later, but by looking at the humble free particle we'll be able to build some intuition for what's going on.

\para
First, something important about the Schrodinger equation:

{\color{olive} The Schrodinger equation 
%
\be i\hbar \ppp{\psi}{t} = -\frac{\hbar^2}{2m} \frac{d^2\psi}{dx^2} + V(x) \psi\nn\ee
%
is linear.} This means that the function $\psi(x,t)$ that we're solving for appears just once in each term. There's no $\psi^2$ or $e^\psi$ or any other horrible function of $\psi$. It's just one $\psi$ in each term. In some terms it's differentiated, in others it's multiplied by something, but there's only ever one. {\color{olive} This means that if $\psi_1(x,t)$ and $\psi_2(x,t)$ are both solutions, then so too is $\psi_1(x,t) + \psi_2(x,t)$.} We can add as many solutions together as we like, or add them with different coefficients in front, and we will still get a solution. 


\para
It's reasonably rare in physics that equations have this property. There are some other places -- the Maxwell equations of electromagnetism are linear like this. But it's not true of the Navier-Stokes equations that describe how water flows, or the Einstein equations for gravity. This is one way -- perhaps the only way -- in which quantum mechanics is easier than classical physics.

\para
{\color{olive} For the free particle, with $V(x)=0$, we have the (non-normalisable) solutions 
%
\be \psi_k(x,t) = e^{-iE_kt/\hbar} e^{ikx} \ \ \ {\rm with}\ \ \ E_k=\frac{\hbar^2k^2}{2m}\ .\ee
%
for any $k\in \R$.} I've put a little subscript $k$ on the wavefunction and the energy to show that they're determined by the choice of $k$. As we saw, these are non-normalisable, meaning that we shouldn't really count them as physical solutions because they don't have any probability interpretation. Nonetheless, they're still useful. In particular, we can use the linearity of the Schrodinger equation to take sums of these solutions and, hopefully, get one that is normalisable.

\para
Because $k$ can take any real value, we can integrate rather than sum. (An integral is just an infinite sum.) {\color{olive} By linearity,
%
\be \psi(x,t) = \int_{-\infty}^{+\infty} dk\ A(k) \,\psi_k(x,t)\nn\ee
%
will also be a solution for any choice of $A(k)$.} What's more, if we pick good choices of the function $A(k)$, then the resulting wavefunction could be normalisable, and so a good quantum state. 

\para
{\color{olive} We will look at the wavefunction with
%
\be A(k) = \exp\left(-\frac{(k-k_0)^2}{2\sigma^2} +\frac{k_0^2}{2\sigma^2}\right)\ .\ee
%
} {\color{blue} Have a plot of the function $A(k)$ here. There was a suggestion to have a simulation of the Gaussian distribution...I'm not quite sure what this means but it sounds like a good idea! (A comment: the power of $\sigma$ is annoyingly wrong in the book -- I've fixed it here. Also, it's not quite the usual standard deviation.)} This is a function that looks like this -- it's peaked at $k=k_0$ and has width given by $\sigma$. There's that extra bit on the end -- the $e^{k_0^2/2\sigma}$ -- that's just a constant that looks annoying but will turn out to simplify things later on.

\para
So what's the intuition behind this? {\color{olive} Recall that the momentum is 
%
\be p = \hbar k\ .\ee
%
The wavefunction $\psi(x,t)$ has a range of different momenta spanning 
%
\be  (k-k_0)^2 \lesssim \mbox{a few times $\sigma^2$}\nn\ee
%
} So we expect this particular solution to correspond to a particle with momentum close to $\hbar k_0$, but with a spread of momentum determined by this number $\sigma$. And, remember, this is a solution to the Schrodinger equation for any choice of $k_0$ and $\sigma$. 

\para
If we substitute this function $A(k)$ into the integral for the wavefunction, we find....{\color{olive} The wavefunction takes the form
%
\be \psi(x,t) &=& \int_{-\infty}^{+\infty} dk \ \exp\left(-\frac{k^2-2kk_0}{2\sigma^2} -i\frac{\hbar t k^2}{2m} + ikx\right)\nn\\
&=& \int_{-\infty}^{+\infty} dk\ \exp\left(-\frac{\alpha}{2}\left( k-\frac{\beta}{\alpha}\right)^2 + \frac{\beta^2}{2\alpha}\right)
\nn\ee
%
} In the first line, the terms are $A(k)$, then $e^{-iEt/\hbar}$ then $e^{ikx}$. In the second line, I've completed the square {\color{olive} with
%
\be \alpha = \frac{1}{\sigma^2}  + \frac{i\hbar t}{m}\  \ \ {\rm and}\ \ \ \beta =  \frac{k_0}{\sigma^2}+ ix\ .\nn\ee
%
}Now we have to do the integral in this equation. This is known as the Gaussian integral. It's pretty well known, and there's a problem that walks you through it. 

{\color{olive} The Gaussian integral is given by 
%
\be\int_{-\infty}^{+\infty} dk\ e^{-\alpha(k-\gamma)^2/2} = \sqrt{\frac{2\pi}{\alpha}}\ .\ee
%
for any $\alpha$ with ${\rm Re}(\alpha)>0$ and any $\gamma$. Using this, our wavefunction becomes
%
\be \psi(x,t) = \sqrt{\frac{2\pi}{\alpha(t)}}\exp\left(-\frac{1}{2\alpha(t)}\left(x-\frac{ik_0}{\sigma^2}\right)^2\right)\nn\ee
%
} This is our final answer. Something rather nice has happened. We started with a Gaussian spread of momenta given by $A(k)$. And, after doing the integral over the different momenta -- really the different wavenumbers $k$ -- we get a wavefunction that is again a Gaussian, but this time in position. It's the same Gaussian form of roughly "e-to-the-minus-x-squared", with some prefactor and some offset. {\color{olive} This is a Gaussian in space. It is called the {\it Gaussian wavepacket}.}

\para
So our next task is to try to interpret this wavefunction. The first thing to note is:

{\color{olive} The wavefunction is normalisable (but not normalised)}.

This follows because it decays exponentially quickly for large positive or negative $x$, so integrating over $x$ will give a finite number. We'll compute the normalisation soon.

\para
The wavefunction doesn't have definite energy or momentum, because we built it by integrating over a range of momenta. We say that there is some uncertainty in the momentum which we write as:

{\color{olive} The uncertainty in the momentum is
%
\be \Delta k^2 \sim \sigma^2\nn\ee
%
}But there's also some uncertainty in the position -- this wavefunction describes the probability of a particle siting at various points. It's a little confusing because there's that factor of $i$ in the $ik_0/\sigma^2$ term and we'll understand what that's doing shortly. But the wavefunction drops off very quickly when $x\gg \alpha(t)$. So 

{\color{olive} Naively, the uncertainty in the position is
%
\be \Delta x^2 \sim \alpha(t) \ .\ee
%
But again $\alpha(t)\in\C$ so it's not clear how to interpret this. If we look at $t=0$, we have
%
\be \Delta x^2 \sim \frac{1}{\sigma}\ .\ee
%
This is our first hint of the Heisenberg uncertainty relation. The uncertainty in the position can be reduced only if we increase the uncertainty in the momentum, and vice versa.} We'll see a more precise version of the Heisenberg uncertainty relation in a few lectures time. 


\para
It's a bit tricky to interpret the wavefunction because it's complex. What does it mean that $x$ is peaked at this imaginary value? How should we think of the spread $\alpha(t)$ as being complex? We can duck these issues if we look at the probability density

{\color{olive} We can compute the probability density
%
\be P(x,t) = C|\psi(x,t)|^2 &=& \frac{2\pi C}{|\alpha|} \,\exp\left(-\frac{1}{2\alpha}( x -ik_0/\sigma^2)^2 -\frac{1}{2\alpha^\star}(x +ik_0/\sigma^2)^2 \right)\nn\\ &=&  
\frac{2\pi C e^{k_0^2/\sigma^2}}{|\alpha|} \exp\left(-\frac{1}{\sigma^2|\alpha|^2}\left(x-\frac{\hbar k_0t}{m}\right)^2\right)\nn\ee
%
where $C$ is a normalisation factor. } You need a little bit of algebra to get to the second line. {\color{olive}You can check that the probability is correctly normalised if we take 
%
\be C = \frac{e^{-k_0^2/\sigma^2}}{\sqrt{4\sigma^2 \pi^{3}}} \ .\ee
%
}
It's much simpler to interpret the probability density. First, we can see that the peak of the distribution is where the expression  $x-\hbar k_0t/m$ vanishes. This means that
%
{\color{olive} The peak of the distribution moves with speed
%
\be v = \frac{\hbar k_0}{m}\ .\ee
%
This is the expected because we summed over wavefunctions with momentum peaked near $p_0=\hbar k_0$. Moreover, the wavefunction spreads out over time, with width
%
\be |\alpha|^2 = \frac{1}{\sigma^2} + \frac{\hbar^2t^2}{m^2}\ .\nn\ee
%
}
Some of the lessons we've learned here are specific to the free particle. In particular, the fact that there is no proper quantum state with a fixed energy (i.e. a stationary state) is unusual -- that's not normally the case. But other lessons carry over to all quantum systems. You can always take sums (or integrals) of a collection of energy eigenstates and get a new time-dependent solution. That's what we've done here and, even though the energy eigenstates weren't normalisable, we can build good normalisable quantum states -- these Gaussian wavepackets. They're not the only states we can build -- we made a specific choice of that function we called $A(k)$ when doing the integral -- but it is a nice choice because it captures what we naively think of as a particle -- something localised in space, here with this Gaussian profile. One surprise, however, is that a free particle doesn't like to be localised. The wavefunction spreads out, getting broader over time. That's a general property of quantum particles -- they want to be de-loclised.

\para
Another general property is the pay-off between the spread in momentum and the spread in position. For this choice of Gaussian wavepacket, we saw that taking a small spread in momentum -- that means small $\sigma$ -- results in a large spread in position -- one that looks like $1/\sigma^2$, and vice-versa. In this example, you can only pin down the particles momentum at the cost of knowing less about its position. That's the essence of something deeper that we'll prove later, known as the Heisenberg uncertainty relation. 

{\color{red} I think this is probably 30+ minutes}

{\color{blue} In the book, I take a first look at some expectation values $<x>$ and $<p>$. I'm not sure this is the right place to do it for these lectures. We can always revisit this decision later, but I might postpone this until it's done properly as part of the fuller formalism of quantum mechanics.}


\subsection{Problem: The Gaussian Integral}

{\color{blue}Should walk them through how to do the Gaussian integral here. There's the question of whether we want to have just real parameters in the exponent, or complex. Note that the Gaussian wavepacket needs a complex $\alpha$ and $\gamma$. Can also get them to compute the normalisation factor $C$ for the probablity. } \alyssa{Computing the Gaussian integral and the normalisation factor will be the practice problem for Lecture 6.} \href{https://drive.google.com/file/d/1xDxOLH_lZp7oQLy4JM21t9CUUGyPKqXj/view?usp=drive_link}{[Lecture 6 Problem Statement]}


\newpage
\section{The Harmonic Oscillator}

In this section, we're going to discuss the quantum harmonic oscillator. I don't think it's overstating the case to say that the harmonic oscillator is, by some margin, the most important example in all of quantum mechanics.  A harmonic oscillator is just a fancy name for a spring bouncing backwards and forwards. It's Hooke's law, if you like. And the reason it's important is twofold. First we can solve it. Which might not sound like a big deal now, but as physics gets more complicated the number of things that we can actually solve gets smaller and smaller until eventually, it's just the harmonic oscillator that's left standing. 

\para
The second reason is that there's whole bunch of techniques in physics which take anything in the universe and try to make it look like a harmonic oscillator. We give these techniques fancy sounding names -- like Fourier analysis or perturbation theory or Feynman diagrams. But, at the end of the day, it's converting whatever complicated equations you have into harmonic oscillators. Sometimes it's an infinite number of harmonic oscillators, so it's not without its complications, but harmonic oscillators nonetheless. I don't think it's unreasonable to suggest that a large part of theoretical physics boils down to the art of making things look like a bunch of springs bouncing back and forth. 

\para
All of which means, we should understand it and understand it well. 



\para
The classical harmonic oscillator is just the fancy name we give to a spring bouncing back and forth. \alyssa{There will be a Tutorial on the classical harmonic oscillator for Lecture 7.} \href{https://drive.google.com/file/d/1ilchItBj6ydShAizOPrLgBx9beY4JqNd/view?usp=drive_link}{[Lecture 7, Tutorial 1]}

{\color{olive} The classical harmonic oscillator obeys the equation
%
\be \ddot{x} = -\omega^2 x\nn\ee
%
This is Hooke's law for a spring. The most general solution has the form 
%
\be x(t) = A\cos(\omega(t-t_0)) \nn\ee
%
with $A$ and $t_0$ integration constants. The constant $\omega$ tells you the frequency at which the particle bounces back and forth. 

\para
This has a conserved energy given by
%
\be E_{\rm classical} = \frac{1}{2}m\dot{x}^2 + \frac{1}{2}m\omega^2x^2\nn\ee
%
with a quadratic potential.} 

{\cblue Draw potential here} Note the potential includes a factor of the mass $m$, just like the kinetic term, but this ensures that it drops out of the equation of motion. 

\para
{\co We want to solve the Schrodinger equation for the quantum harmonic oscillator
%
\be -\frac{\hbar^2}{2m}\frac{d^2\psi}{dx^2} + \frac{1}{2}m\omega^2 x^2\psi = E\psi\nn\ee
%
}
As always, our goal is to find all solutions to this equation. That means that we should find all possible normalisable wavefunctions $\psi(x)$ and their associated energies $E$. 
There's various ways to do this. I'm going to show you a slightly slick way. It involves a trick, but it's a trick that's well worth knowing as it appears in similar forms in many other places. The way I'm going to present this is slightly less rigorous than it could be. I'll show you all the solutions to the Schrodinger equation, but I won't prove to you that they're all the solutions. If you want to do better, you'll need to use different methods and these are provided in the book. 


\para
The first part of the trick is you stare hard at the Schrodinger equation and wait for inspiration to strike. (This is a good trick to pull off if you can manage it!) What that means in this case is {\co we notice the following solution
%
\be \psi_0(x) = e^{-m\omega x^2/2\hbar} \ \ \ {\rm with}\ \ \ E_0=\frac{1}{2}\hbar\omega\nn\ee
%
}To see that this solves the Schrodinger equation, differentiate to get {\co
%
\be \frac{d\psi_0}{dx} = -\frac{m\omega x}{\hbar}\psi_0\ \ \ {\rm and}\ \ \ \frac{d^2\psi_0}{dx^2} = -\frac{m\omega}{\hbar}\psi_0 + \frac{m^2\omega^2 x^2}{\hbar^2}\psi_0\nn\ee
%
So this obeys the Schrodinger equation with $E_0=\ft 12 \hbar \omega$.}

\para
Ok. So that's nice. We've managed to get one solution by just writing it down. And, moreover, it takes the Gaussian form. Remember that the Gaussian wasn't an energy eigenstate of the free particle -- we had to work hard in the last lecture to see how it could be a solution to the time-dependent Schrodinger equation, a wave-packet, but it definitely spread out over time. Here, however, we have a potential pinning the particle close to the origin and the Gaussian is an energy eigenstate. But, as you can see, it has a very specific width and that width is set by the strength of the potential. A steeper potential means a more sharply peaked Gaussian. 

\para

One of the things that I'm not going to prove here, but is true, is that this is the solution with the lowest possible energy. 

\para
{\co In quantum mechanics, the wavefunction with the lowest possible energy is called the {\it ground state}. For the harmonic oscillator $\psi_0(x)$ is the ground state.}

\para 
Now for the next part of the trick. And it relies on the following result:

\para
{\co {\bf Claim:} If $\psi$ is a solution to the Schrodinger equation with energy $E$ then
%
\be \left(\sqrt{\frac{m\omega}{2\hbar}}x +\sqrt{\frac{\hbar}{2m\omega}}\frac{d}{dx}\right)\psi \nn \ee
%
is a solution with energy $E+\hbar \omega$.

\vskip 3mm
\noindent
{\bf Proof:}} {\cblue Still need to write this up. Should be straightforward - just substitute the expression above into the Schrodinger equation. We could even set this as a problem perhaps?}

\para
So that's rather nice -- it's a kind of solution-generating trick. If you can find one solution, then you can find an infinite number of them. Just keep acting with this combination of $x$ and $d/dx$ and you'll get a whole tower of solutions, with increasing energy. Moreover, we have one solution to start with, a seed solution $\psi_0(x)$ with energy $E_0=\ft12 \hbar \omega$. 

\para
{\co We can build up new solutions starting from $\psi_0(x)$. The next few are
%
\para
The next few (unnormalised) wavefunctions are
%
\be \psi_1(x)  = \left(\sqrt{\frac{m\omega}{2\hbar}}x +\sqrt{\frac{\hbar}{2m\omega}}\frac{d}{dx}\right)\psi_0(x) = 2\sqrt{\frac{m\omega}{\hbar}}  x \,e^{-m\omega x^2/2\hbar}\nn\ee
%
This has energy $E_1= \frac{3}{2}\hbar\omega$.  Then
%
\be \psi_2(x)  = \left(\sqrt{\frac{m\omega}{2\hbar}}x +\sqrt{\frac{\hbar}{2m\omega}}\frac{d}{dx}\right)\psi_1(x) = \left(-2+\frac{4m\omega}{\hbar}x^2\right)e^{-m\omega x^2/2\hbar} \nn\ee
%
This has energy $E_2=\frac{5}{2}\hbar\omega$.}

\para
{\cblue Could set as a problem
\be
\psi_3(x) =  \left(\sqrt{\frac{m\omega}{2\hbar}}x +\sqrt{\frac{\hbar}{2m\omega}}\frac{d}{dx}\right)\psi_2(x) =  \sqrt{\frac{m\omega}{\hbar}}\left(-12x +\frac{8m\omega}{\hbar}x^2\right)e^{-m\omega x^2/2\hbar} \nn\ee
%
} \alyssa{This is the practice problem for Lecture 7.} \href{https://drive.google.com/file/d/1UBXGh8inxjkUpZ_Nq9do-vqrjMrkI7Xu/view?usp=drive_link}{[Lecture 7 Problem Statement]}

So what have we found here. First, we've got a way to construct an infinite tower of energy eigenstates. {\co The spectrum of the Hamiltonian} -- meaning the collection of all energy eigenvalues -- {\co is given by
%
\be E = \left(n+\frac{1}{2}\right) \hbar \omega\ \ \ {\rm with}\ \ \ n=0,1,2,3,\ldots\nn\ee
%
}
Although I won't prove it here, it turns out that this is all the normalisable solutions to the Schrodinger equation. This simple method gets everything.

\para
So what do these wavefunctions look like? Well, they all have the same exponential fall off at large $x$. In particular, that means they're all normalisable although I haven't made any attempt to actually normalise them. Then there's a polynomial in front of the exponential. For what it's worth, these are called Hermite polynomials. They're a class of very well studied functions in mathematics, starting in the 1800s. There's a massive wikipedia page that describes their various properties.


\para
{\cblue Plot the first few wavefunctions here. We have a sim that allows you to visualise any solution of the HO \href{https://capacity-simulations.github.io/simulations/harmonic-oscillator_2.html}{here} }

\para
We can plot these wavefunctions. The first one is just the Gaussian. The next one looks like this, then this, and so on. As we go up in energy, the polynomials that sit in front are higher order, and so have more zeros, and so the wavefunctions cross the axis more and more times. The ground state is the Gaussian and only vanishes asypmtotically. The first excited state $\psi_1$ crosses the axis once; the next state $\psi_2$ crosses it twice, and so on. That, it turns out, is common to any problem of this sort; whatever the shape of the potential -- provided it rises asymptotically -- the wavefunctions have more and more nodes at higher energies. 

\para
So these are the solutions. We've got the spectrum of energies -- that turns out to be a very important result. But what are we to make of these wavefunctions. The classical harmonic oscillator has a ball bouncing back and forth. But these wavefunctions don't seem to look anything like that! Even by eye, you can see that the average position of the particle is at $x=0$. (But then, if you think about it, the average position of a particle bouncing back and forth is also at $x=0$ so perhaps that's not so surprising.)

\para
Of course, in some sense these solutions couldn't ever look like a particle bouncing back and forth. We set out looking for energy eigenstates and, by definition, these don't depend on time. Nonetheless, if you know where to look, you can see some hints of the classical physics. 


\para
There's a general rule of thumb that quantum systems  tend to look more classical as you go to higher energies. That's because the discreteness in the energy becomes less obvious. If the gap between different energy levels is tiny compared to the total energy, then you're less likely to notice it -- it looks much more like a continuum of possible energy levels, which is what happens classically. So this suggests that it might be easier to interpret the wavefunctions with very high energies. Here we plot the wavefunctions  for the energy levels $n=20$ and $n=60$.

\para
{\cblue plot wavefunctions}

\para
First, you can see that as the energy $n$ increases, the wavefunction spreads further and further from the origin. That's expected classically -- the higher the energy, the further the ball rolls.  You can make that agreement precise -- the edge of the quantum mechanical wavefunction agrees with how far a classical ball with that energy will travel. 

\para
Next, look at the way the wavefunction oscillates. We know that a free particle with definite momentum $p=\hbar k$ is associated to $e^{ikx}$. The bigger $k$, the smaller the wavelength, and the higher the momentum. In the wavefunctions, you can see that the wavelength of oscillations is much smaller near $x=0$, which corresponds to the  bottom of the potential. The wavelength then  gets stretched towards the edges. This coincides with the classical expectation, where the particle is travelling much faster at the bottom of the potential and slows as it rises.



\para
Finally, look at the overall envelope of the wavefunction, a curve that peaks at the edges and dips in the middle.  This is telling us that if you do a measurement of a particle in this high-energy state, then you're more likely to find it near the edges than near the origin. But this, too, is the same as the classical picture. If you take a photograph of a ball oscillating in a harmonic potential, you're most likely to catch it sitting towards the end of its trajectory, simply because it's going slower and so spends more time in those regions. Conversely, the ball is less likely to be caught as it whizzes past the origin.  This is, again, captured by the form of the wavefunction. 
We learn that the quantum world is not completely disconnected from the classical. You just have to know where to look.

\para
{\co 
At high energies, the wavefunctions:
\begin{itemize}
\item Spread further out as $n$ increases
\item Show the particle moving faster at the bottom and slower at the top
\item Show that the particle is more likely to be found at the edge, where it classically spends more time.
 \end{itemize}
All of this agrees with our classical expectations.}

\para
For the free particle, there were no normalisable energy eigenstates, but we could sum the non-normalisable states to build the Gaussian wavepacket, which looks like it has a reasonably well-defined position and momentum. It's the closest we get to constructing something like a classical particle in quantum mechanics. We can do the same here. 


{\co
There's a solution to the full time-dependent Schrodinger equation called a {\it coherent state} that takes the form
%
\be \psi(x,t) = \exp\left(-\frac{m\omega}{2\hbar}(x-x_c(t))^2 + \frac{i}{\hbar}p_c(t)(x-\ft12 x_c(t)) + i\phi(t)\right)\nn\ee
%
where $x_c(t)$ and $p_c(t)$ are the classical position and momentum of a particle in a harmonic oscillator
%
\be x_c(t) = A\cos(\omega(t-t_0)) \ \ \ {\rm and}\ \ \ p_c(t) = -Am\omega \sin(\omega(t-t_0))\nn\ee
%
for any choice of $A$ and $t_0$, and
%
\be \phi(t) = -\frac{\omega t}{2} - \frac{mA^2}{8\hbar} \cos(2\omega(t-t_0))\nn\ee
%
} I don't know of a nice classical interpretation of this function $\phi(t)$. It's just something annoying that we have to include to get a solution to the equation. But, because it comes with a factor of $i$ in the exponent, it drops out when we compute the probability distribution of the particle (as, indeed, does that second term with $p_c(t)$.) {\co After we normalise, we have
%
\be P(x,t) = |\psi(x,t)|^2 = \sqrt{\frac{m\omega}{\pi\hbar}}\exp\left(-\frac{m\omega}{\hbar}(x-x_c(t))^2\right)\nn\ee
%
} Now we see clearly that the probability distribution is sharply peaked around the classical trajectory $x=x_c(t)$. The coherent state doesn't have a well-defined energy, but it does have something approximating a well-defined position. 

\para
For the free particle, when we built the Gaussian wavepacket, the probability spread out over time. But it doesn't happen here -- the width remains fixed. This is a special property of the harmonic oscillator that doesn't survive for other potentials. 

\para
{\cblue Probably don't need this, but the coherent states are built as
%
\be \psi_{\alpha} (x,t) = \sum_{n=0}^\infty \frac{\alpha^n}{\sqrt{n!}} e^{-iE_nt/\hbar}\psi_n(x)\nn\ee
%
with $\alpha = \sqrt{m\omega/2\hbar} A e^{i\omega t_0}$.}


\para
{\cblue I haven't mentioned creation and annihilation operators here, although they're clearly implicitly present in the construction of the higher states. We might revisit the harmonic oscillator briefly later in the course as a way of illustrating the more algebraic approach, using commutation relations. If we want to, that could be a good time to make the connection to quantum fields.}


\newpage
\section{Bound States}

So far we've met the free particle -- where there are no normalisable energy eigenstates because the wavefunction can spread to infinity -- and the harmonic oscillator, where the potential grows asymptotically, ensuring that all states are trapped. 

\para
But often we're interested in situations that have some features of both. We can have potentials $V(x)$ that look like the situation shown here. {\cblue Draw a random potential that wiggles up and down}. There are some dips and bumps somewhere in the middle, but asymptotically the potential vanishes. One very important quantum system that we'll meet at the end of these lectures is the hydrogen atom, where the electron is bound to the proton through the Coulomb force, and that has similar properties albeit in three dimensions.

\para
{\co We will consider situations with a potential that has the behaviour
%
\be V(x)\ \rightarrow\ 0\ \ \ \ \mbox{as $x\ \rightarrow \pm \infty$}\nn\ee
%
As always, we want to solve the Schrodinger equation
%
\be -\frac{\hbar^2}{2m}\frac{d^2\psi}{dx^2} + V(x)\psi = E\psi\nn\ee
%
We can start by looking at this equation far from the potential at $x\ra \pm \infty$. Since $V(x) \ra 0$, we have the equation for a free particle
%
\be -\frac{\hbar^2}{2m}\frac{d^2\psi}{dx^2}  = E\psi\nn\ee
%
There are two kinds of solutions to this equation:}
%
\begin{itemize}
{\co
\item Scattering states: These have energy  $E>0$ and wavefunction
%
\be \psi = e^{ikx}\ \ \ {\rm with}\ \ \ E=\frac{\hbar^2k^2}{2m}\nn\ee
%
with $k\in \R$. These states are non-normalisable.}

For the free particle, we threw these states away as they are non-normalisable. But we also found that they were useful because we could use them to construct the Gaussian wavepackets. We'll find another, slightly different, interpretation of these states in the next lecture. 


\item The second class of solutions wasn't available for the free particle. These are

{\co Bound states: the solutions with energy $E<0$ take the form
%
\be \psi  = Ae^{-\eta x} + B e^{+\eta x}\ \ \ {\rm with}\ \ \ E=-\frac{\hbar^2\eta^2}{2m}\nn\ee
%
with $\eta\in \R$. For a free particle, these would be badly non-normalisable as they diverge at $x\ra \infty$ or $x\ra -\infty$.  For example, we could set $B=0$ so that the wavefunction $\psi=e^{-\eta x}$ decays nicely as $x\rightarrow +\infty$, but then it will blow up at $x\rightarrow -\infty$. 



\para
However, it may be possible to find solutions like this to the full Schrodinger equation with $V(x) \neq 0$ so that the wavefunction decays at both ends, with
%
\be \psi \sim \left\{\begin{array}{lc} e^{-\eta x}\ \ \ & \mbox{as $x\rightarrow +\infty$}\\ e^{+\eta x}\ \ \ & \mbox{as $x\rightarrow -\infty$}\end{array}\right. \nn\ee
%
Note that to solve the Schrodinger equation we must have the same value of $\eta$ in the exponent} (because, far from the potential, both sides must satisfy the free Schrodinger equation with the same value of the energy $E$.) {\co These solutions are trapped in the potential. They exist (if at all) only for very specific values of $\eta$.} 
\end{itemize}

The key idea here is all about the normalisablity of the wavefunction.  For most values of $\eta$, if you choose to have an exponentially decaying wavefunction at one end, say $\psi(x) \ra e^{+\eta x}$ as $x\ra -\infty$, then the full Schrodinger equation tells you have this function evolves as you increase $x$. The potential $V(x)$ changes the function. By the time you get to very large $x$, where the potential vanishes again, you'll typically be left with a solution that has both $\psi(x)  = A e^{-\eta x} + Be^{\eta x}$ with both $A$ and $B$ non-zero. But for very special values of $\eta$ you might find a solution with $B=0$ so that the solution is normalisable on both ends. 

\para
We can see how this works by looking at an example.



\subsection{An Example: The Finite Potential Well}

{\co Consider the potential
%
\be V(x) = \left\{\begin{array}{cc} -V_0\ \ \ &\ \ \ -a<0<a\\ x \ \ \ & \ \ \ {\rm otherwise}\end{array}\right.\nn\ee
%
with $V_0>0$ so the potential is a dip.} {\cblue Plot this}

\para
We want to solve the Schrodinger equation in this potential. It's a little peculiar, because the potential itself is discontinuous. But it turns out there's a simple way to proceed. First we solve the Schrodinger equation in the middle region $-a<x<a$. Then we solve it in the outside region $|x|>a$. And then we patch the two together. 


\para
{\cblue There's a story here that's in the book, but I'm brushing under the rug, about how the wavefunction and its derivative should be continuous, but $\psi''(x)$ is discontinuous. Could explain this as extra material perhaps.} \alyssa{This will be included in the Tutorial for Lecture 8.} \href{https://drive.google.com/file/d/1fw8hOjUlY-BsJ0BJx1Xko-qa6W9G8FRz/view?usp=drive_link}{[Lecture 8, Tutorial 1]}

\para
Before we proceed, there's an observation that makes our life easy. The potential is an even function, meaning that it looks the same when reflected about the origin; it has reflection symmetry. This means that all energy eigenstates are either even functions or odd functions. (If you don't know what this means, I'll explain shortly.) Here's the proof

\para
{\co If $\psi(x)$ solves the Schrodinger equation for some energy $E$ then so too does $\psi(-x)$. (Because $V(x) = V(-x)$. If there is only one solution for each energy then we must have
%
\be \psi(x) = \alpha \psi(-x)\nn\ee
%
for some $\alpha\in \C$. But we have
%
\be \psi(x)= \alpha\psi(-x)  = \alpha^2\psi(-(-x)) =  \alpha^2\psi(x)\nn\ee
%
which tells us that $\alpha =\pm 1$, corresponding to either even ($\alpha=1$) or odd ($\alpha=-1$) wavefunctions.} 

\para
The upshot is that when looking for solutions, we can take the wavefunction is either even or odd. The calculations are generally much easier if you look for solutions of this form. (Note that to prove this, we assumed that there is only one solution with each energy level. But it's also true if there are multiple solutions with the same energy; you just need to work a little harder to prove it.)

\para
{\co We will look for even wavefunctions. Outside the potential, where $V(x) =0$, we take

%
\be \psi(x) = \left\{\begin{array}{cc} A e^{-\eta x}\ \ \ &\ \ \ {x>a}\\Ae^{+\eta x}\ \ \ & \ \ \ x<-a\end{array}\right.\nn\ee
%
with $A$  a normalisation constant. Our goal is to find the allowed values of $\eta$ which will determine the energy of the bound state $E=-\hbar^2\eta^2/2m$.

\para
Inside the potential, the Schrodinger equation is
%
\be -\frac{\hbar^2}{2m}\frac{d^2\psi}{dx^2}  = (E+V_0)\psi\ \ \ {\rm for} \ \ -a<x<a\nn\ee
%
This is again the free Schrodginer equation but with a shifted energy. } This is because the potential is just constant here..life would be harder if the potential actually changed in space in this region. 

{\co We're looking for solutions with $E<0$ but with $E+V_0>0$, so the solutions in this region are plane waves of the form $e^{\pm i kx}$ for some $k$. But because we're looking for even functions, we can take
%
\be \psi(x) = B\cos kx\ \ \ {\rm for} \ \   |x|< a\nn\ee
%
with $B$  a normalisation constant. Here $k$ determines the energy, now with
%
\be E+V_0 = +\frac{\hbar^2 k^2}{2m}\nn\ee
%
So we have
%
\be E = -\frac{\hbar^2\eta^2}{2m} = \frac{\hbar^2 k^2}{2m} - V_0\nn\ee
%
} This gives a relationship between the fall-off determined  by $\eta$ outside the potential, and the wavelength of oscillation $k$ inside the potential. But we still haven't figured out what values of $\eta$, and hence the energy, are allowed.  For this we need to see how the two solutions -- one outside and one inside -- are patched together.

\para
{\co We patch together the solutions inside and outside the potential by requiring that $\psi(x)$ and $\psi'(x)$ are continuous.}


\para
\alyssa{The below justification will be included in a Tutorial for Lecture 8.} \href{https://drive.google.com/file/d/1gpp5nltJft6RJ99iGAlz_Bhno3Qyn4yf/view?usp=drive_link}{[Lecture 8, Tutorial 2]}
{\cblue  We might want to justify that  $\psi(x)$ and $\psi'(x)$ are both continuous at $x=\pm a$ where the potential is discontinuous, whether  in the lecture or in some additional material. To see it,  start with the Schrodinger equation and integrate it over a small interval around $x=a$
%
\be && -\frac{\hbar^2}{2m} \int_{a-\epsilon}^{a+\epsilon}dx\ \frac{d^2\psi}{dx^2} =  \int_{a-\epsilon}^{a+\epsilon}dx\ (E-V(x))\psi \nn\\
\Longrightarrow \ \ \ && \left.\frac{d\psi}{dx}\right|_{a+\epsilon} - \left.\frac{d\psi}{dx}\right|_{a-\epsilon} = -\frac{2m}{\hbar^2} \int_{a-\epsilon}^{a+\epsilon}dx\ (E-V(x))\psi  \nn\ee
%
Although $V(x)$ is discontinuous at $x=a$, it is finite. This means that the integral on the right-hand side vanishes as we take limit $\epsilon \rightarrow 0$, telling us that $\psi'$ is continuous at $x=a$. But if $\psi'(x)$ is continuous, then so too is $\psi(x)$ itself. }

\para
{\co Because the wavefunction is even, we only need to do this at $x=a$. (It's then guaranteed to work at $x=-a$.) We have
%
%
\be \psi(x)\ \mbox{cts at $x=a$} \ \ \ &\Longrightarrow& \ \ \ \ \ B\cos ka = Ae^{-\eta a} \nn\\ 
\psi'(x)\ \mbox{cts at $x=a$} \ \ \ &\Longrightarrow& \ \ -kB\sin ka = -\eta Ae^{-\eta a} \nn\ee
%
}
Dividing the second condition by the first gives us 
{\co 
%
\be k\tan ka = \eta\nn\ee
%
This relates $k$ and $\eta$. It should be solved togeher with our previous condition 
%
\be k^2 + \eta^2 = \frac{2mV_0}{\hbar^2}\nn\ee
%
Finding solutions to these will then fix the energy $E$.}

\para
These are not equations that have a simple solution. But it's easy to get a feel for what's going on by just plotting these graphs of $k$ vs $\eta$. {\cblue Show plots here} The second equation is just a circle, while the first is plotted as $\eta = k\tan ka$, which are the lines with vertical asymptotes.  We get a solution whenever the two curves intersect.  This whole story is about getting a normalisable wavefunction so we better have $\eta>0$. As expected, there are  only discrete solutions, happening for specific values of $\eta$. And we see that there are only a finite number of solutions.





\para
The first solution is guaranteed because, no matter how small the radius $2mV_0/\hbar^2$, the circle will always intersect the first curve. However, the number of subsequent solutions depends on the parameters in the game. We can see that the number of solutions will grow as we increase $V_0$, the depth of the well, since this increases the radius of the circle. The number of solutions will also grow as we increase the width of the well; this is because the separation between the vertical asymptotes is determined by the divergence of the tan function and so occurs when $k\approx \pi/2a$. As we increase $a$, keeping $V_0$ fixed, the vertical-ish lines get closer together while the  circle stays the same size.







\para
{\co 
If we plot these curves graphically, we find
%
\begin{itemize}
\item There is at least one solution.
\item The get more solutions as the potential well gets deeper (increasing $V_0$) or wider (increasing $a$). 
\end{itemize}
We can be more precise. The $n^{\rm}$ solution has $k=k_n$ and energy
%
\be E_n = \frac{\hbar^2k_n^2}{2m}-V_0\ .\label{enwindy}\nn\ee
%
From the graph, the $n^{\rm th}$ crossing occurs somewhere in the region
%
\be (n-1)\pi < k_na < \left(n-\frac{1}{2}\right)\pi\nn\ee
%
This gives an estimate of the energy
}

\para
There is also a mild surprise in the wavefunction. Classically, the particle with energy $E<0$ would be trapped inside the well. But the quantum wavefunction extends to the region $|x|>0a$, although it is exponentially suppressed there. 

\para
{\co The wavefunction tells us that there is some probability of finding the particle outside the potential well, even though the classical particle with energy $E<0$ cannot possibly escape.} {\cblue Maybe plot the probability density here?}










\subsection*{Perhaps an Exercise: Odd Parity Wavefunctions}
\alyssa{This will be included as a practice problem for Lecture 8.} \href{https://drive.google.com/file/d/1mXMniOTgfhR4k8XUgV0cop9buB_VE7zs/view?usp=drive_link}{[Lecture 8 Problem Statement]}
{\cblue




We can repeat this analysis in the case of odd parity wavefunctions. We're now looking for bound state solutions which, outside the potential well, take the form
%
\be \psi(x) = \left\{\begin{array}{cc} A e^{-\eta x}\ \ \ &\ \ \ {x>a}\\-Ae^{+\eta x}\ \ \ & \ \ \ x<-a\end{array}\right.\nn\ee
%
Meanwhile, the odd parity wavefunction inside the well is
%
\be \psi(x) = B\sin kx\ \ \ {\rm for} \ \  |x|< a\nn\ee
%
Patching these two solutions at $x=a$ gives
%
\be \psi(x)\ \mbox{cts at $x=a$}\ \ \  &\Longrightarrow& \ \ \ \ \ B\sin ka = Ae^{-\eta a} \nn\\ 
\psi'(x)\ \mbox{cts at $x=a$}\ \ \  &\Longrightarrow& \ \ kB\cos ka = -\eta Ae^{-\eta a}\nn \ee
%
which now gives us
%
\be \frac{k}{\tan ka} = -\eta\nn\ee
%
Once again this should be solved in conjunction with 
%
\be k^2 + \eta^2 = \frac{2mV_0}{\hbar^2}\nn\ee
%


\para
This time there is no guarantee that a solution exists. If we plot the circle superimposed on the equation for $\eta$, we find that the first solution only exists provided 
%
\be \frac{2m V_0}{\hbar^2} > \frac{\pi^2}{4a^2}\ \ \ \Longrightarrow\ \ \ \frac{8mV_0}{\hbar^2\pi^2} > \frac{1}{a^2}\nn\ee
%
This means that the first parity odd solution exists only if the potential is deep enough or wide enough.}


\subsection*{Comment: Delta-Function Potential}
\alyssa{This will be included as a Tutorial for Lecture 8.} \href{https://drive.google.com/file/d/1gpp5nltJft6RJ99iGAlz_Bhno3Qyn4yf/view?usp=drive_link}{[Lecture 8, Tutorial 2]}
{\cblue The delta-function potential offers a nice example of how this works, because you see the delta function changes the slope of the wavefunction by a fixed amount and you need that to turnover the  $e^{\pm \eta x}$ wavefunction profile. I've included it here in case we want to add it to some online content. 

\para
The delta-function potential is given by
%
\be V(x) = -V_0\, \delta(x)\nn\ee
%
for some constant $V_0$. The delta-function is an infinitely spiky function (really a distribution....might need to include an explainer on this, if not here then in later lectures) satisfying
%
\be \delta(x) = 0\ {\rm unless}\ x=0\ \ \ {\rm and}\ \ \ \int_{-\infty}^{+\infty} \delta(x) =1\ .\ee
%
You can think of it as an infinitely narrow, but infinitely long spike localised at the origin. You might think that it's not particularly realistic as a potential, and there is some merit to that. But there are situations -- such as impurities in solids -- where the exact form of the potential is complicated and most likely unknown, and it is useful to have a simple toy model that can be easily solved. This is what the delta function offers.

\para
The discontinuity in the delta function is significantly more extreme than that of the finite potential well and, once again, we're going to have to understand how to deal with it. Our strategy is the same as before: we take the Schrodinger equation as the starting point and see what the potential means for the wavefunction. To start, we integrate the Schrodinger equation over a small interval around the origin
%
\be -\frac{\hbar^2}{2m}\int_{-\epsilon}^{+\epsilon}dx\ \frac{d^2\psi}{dx^2} = \int_{-\epsilon}^{+\epsilon}dx\ (E+V_0\,\delta(x))\psi\ .\ee
%
This time the right-hand side isn't so innocent. While the $E$ term simply vanishes in the limit $\epsilon\rightarrow 0$, the delta function doesn't. We then find
%
\be \lim_{\epsilon\rightarrow 0} \left( \left.\frac{d\psi}{dx}\right|_{+\epsilon} - \left.\frac{d\psi}{dx}\right|_{-\epsilon}\right)  = -\frac{2mV_0}{\hbar^2}\,\psi(0)\ .\label{ada}\ee
%
We learn that the delta function leaves its imprint in the derivative of the wavefunction, and $\psi'$ must now be discontinuous at the origin. The wavefunction itself, however, should be continuous.


\para
This is all the information that we need to find a bound state. Away from the origin, the negative energy, normalisable solutions of the Schrodinger equation are simply
%
%
\be \psi(x) = \left\{\begin{array}{cc} A e^{-\eta x}\ \ \ &\ \ \ {x>0}\\Ae^{+\eta x}\ \ \ & \ \ \ x<0\end{array}\right.\ee
%
for some $\eta >0$. Note that we've chosen normalisation factors that ensure the wavefunction is continuous,
%
\be \lim_{x\rightarrow 0^+}\ \psi(x) = \lim_{x\rightarrow 0^-}\ \psi(x) = A\ .\ee
%
However, the presence of the delta function means that the two solutions must be patched together with a discontinuity in the derivative
%
\be\lim_{x\rightarrow 0^+}\ \psi'(x) -\lim_{x\rightarrow 0^-}\ \psi'(x)  &=& -A\eta \left(\lim_{x\rightarrow 0^+}\ e^{-\eta x} + \lim_{x\rightarrow 0^-}\ e^{+\eta x}\right)\nn\\ &=& -2A\eta 
\ .\ee
%
This expression should be identified with the discontinuity \eqn{ada}, giving the result that we need
%
\be \eta = \frac{mV_0}{\hbar^2}\ .\ee
%
We learn that the negative-valued delta function has just a single bound state with energy
%
\be E = -\frac{\hbar^2\eta^2}{2m}  = -\frac{V_0^2m}{2\hbar^2}\ .\label{1ddeltabound}\ee
%
This calculation also makes it clear how the presence of a potential can change the asymptotic behaviour from $e^{\eta x}$ as $x\rightarrow-\infty$ to $e^{-\eta x}$ as $x\rightarrow +\infty$. The delta function does this all in one go, but can only achieve the feat for a very specific value of $\eta$. More general potentials flip the sign of the exponent more gradually, but again can only do so for specific values of $\eta$, leading to a (typically finite) discrete collection of bound states.}


\newpage
\section{Scattering}


There's one last aspect of physics that we're going to talk about for a particle moving on a line. This is scattering.

\para
The basic idea behind scattering is simple. You can take a particle and you throw it at an object and watch how it bounces back. Then, ideally, you use the information about how it bounces to infer information about the object in question.  A particularly simple example of this is called ``looking at things''. But the general principle of throwing stuff at other stuff has evolved into the most powerful experimental technique in all of physics. It's how we understand the structure of the atom, and the structure of DNA and, more recently, the populations of black holes (although in that case, where one black hole collided with another and looked at the emission of gravitational waves, someone else was doing the throwing.)


\para
In this lecture, we're going to look at some basic ideas that arise when we scatter quantum particles.  Our primary goal is to continue to build intuition for how quantum particles behave. 




\para
{\co Our set-up is the same as the previous lecture: a potential $V(x)$  that is localised in space, so that
%
\be V(x)\ \rightarrow 0 \ \ {\rm  as}\ |x|\ \rightarrow \ \infty\ .\nn\ee
% 
}
In the last lecture, we understood that such potentials typically have some finite number of negative energy bound states, trapped in the potential. Here we instead ask: what happens if we stand far from the potential and throw in a quantum particle? Will it bounce back, or will it pass through the potential? Or, because this is a  quantum particle and can do weird things, will it do both?




\para
The first thing we have to do is set up the problem mathematically. And there's a subtlety in how we do this.

\para
The most obvious approach would be to describe the particle as a kind of  Gaussian wavepacket that we met in a previous lecture -- after all, that's the closest we have to a classical particle in the quantum world. Then we could just evolve this wavepacket through the time-dependent Schrodinger equation and see what happens. That has the advantage that it's rather intuitive.

\para
{\cblue It occurs to me that it might be possible to do this numerically and show then what happens. It would be nice to cook up a simulation that does this.}

\para
While that approach is conceptually straightforward, it's not mathematically straightforward. We saw in the previous lecture that even the free particle in quantum mechanics is not particularly straightforward. But there is, it turns out, an easier way.

\para
{\co We will work with the wavefunctions
%
\be \psi_k(x) = A  e^{ikx}\nn\ee
%
for some constant $A\in \mathds{C}$. These have definite momentum $p=\hbar k$.

\para
We can't think of these as describing a single particle because they are non-normalisable. Instead we give them a different interpretation: they describe a continuous beam of particles with probability density
%
\be P(x,t) = |\psi(x,t)|^2 =|A|^2\nn\ee
%
This is now interpreted as the average density of particles. To reinforce this perspective, we can compute the probability current} (which we met in an earlier lecture)
%
{\co
\be J(x,t) &=& -\frac{i\hbar}{2m}\left(\psi_k^\star\frac{d\psi_k}{dx} - \psi_k\frac{d\psi_k^\star}{dx}\right) \nn\\&=& -\frac{i\hbar}{2m} \times |A|^2 \times 2ik = |A|^2 \times \frac{p}{m}\ \ \ \ \label{flux}\ee
%
This is the average density of particles,  multiplied by their velocity. It's the average flux of particles.}

\para
This framing of the question has the advantage that, despite the dynamic nature of scattering, nothing is changing with time. We will look at the equilibrium set-up of a beam of particles coming in, a beam bouncing band and another beam passing through. From this we will be able to extract the probability that a particle bounces back or passes through. And we can do it by solving the time-independent Schrodinger equation, which is much easier and cleaner.
 
 

\subsection{Tunnelling}

{\co We're going to look at just a single example. It's a finite potential barrier, like a bump in the road
%
\be V(x) = \left\{\begin{array}{cc} U\ \ \ &\ \ \ -a<x<a\\ 0 \ \ \ & \ \ \ {\rm otherwise}\end{array}\right.\ \  \ \ \ \  \ \ \ \  \ee
%
where $U>0$.} This is like an inverted potential well whose bound states we studied in the last lecture. It's just now the bump goes up rather than down.


We will throw a particle, or more precisely a beam of particles with energy $E$, in from the left and see what happens. Our classical expectation is that the particle would bounce back if $E<U$ and would fly over the barrier if $E>U$. We will see that quantum particles are more interesting. 

\para
We first need to  set up the problem. {\co We're looking for a solution to the Schrodinger equation that has an ingoing beam of particles
%
\be \psi(x) \sim A e^{ikx}\ \ \ \ {\rm as}\ x\rightarrow -\infty\nn\ee
%
with $A$ the density of the beam. We should take $k>0$ as this beam is travelling to the right. The  energy of each particles in the beam is
%
\be E = \frac{\hbar^2k^2}{2m}\nn\ee
%
}

\para
{\co We expect that some of the beam will bounce back. That will give a solution on the far left which looks like
%
\be \psi(x) = Ae^{ikx} + Be^{-ikx}\ \ \ {\rm as}\ x\rightarrow -\infty\nn\ee
%
The $e^{-ikx}$ means that this is a returning beam, moving to the left. It has a difference coefficient $B\in \C$ because the density of the beam can differ. } Not all particles will be reflected back. {\co We would expect $|B|<|A|$ as not everything bounces back.} That's something we can check after the calculation is done. 

\para
Next we should think about what happens on the other side of the potential as $x\ra \infty$. It's quite possible that some of the beam makes it through, so {\co We should have a term in the solution that looks like
%
\be \psi(x) = C e^{ikx} \ \ \ {\rm as}\ x\ra +\infty\nn\ee
%
with $C\in \C$.} Notice that the momentum $k$ is the same as the incoming momentum -- that's guaranteed on grounds of energy conservation. Said differently, we'll solve the Schrodinger equation with the same value of $E$ on the left and right. 

\para
We only have the $e^{ikx}$ on the right; we don't have $e^{-ikx}$ since that would correspond to a beam moving left. But we're not throwing in particles from that side, only from the other side. {\cblue Simple picture here to describe what's going on.}

\para
The solutions above hold everywhere away from the potential, so $|x|>a$. What do we do where there's the bump?

\para
{\co We'll look at the situation $E<U$, where classically the particle bounces back. In the middle of the bump, we have
%
\be -\frac{\hbar^2}{2m} \frac{d^2\psi}{dx^2} = (E-U)\psi\nn\ee
%
For $E<U$, this has solutions 
%
\be \psi(x) = F e^{-\eta x} + G e^{+\eta x}\nn\ee
%
with $F,G\in \C$ and 
%
\be \eta = \sqrt{\frac{2m(U-E)}{\hbar^2}}\nn\ee
%
The upshot is that the Schrodinger equation is solved by 
%
\be \psi(x) =\left\{\begin{array}{cc}Ae^{ikx} + Be^{-ikx} & \ \ \ x<-a\\ Fe^{-\eta x} + Ge^{+\eta x} & \ \ \ |x| <a \\ 
Ce^{ikx} & \ \ \ x>a\end{array}\right.\ .\ee
%
Our next task is to insist that $\psi(x)$ and $\psi'(x)$ are both continuous at $x=\pm a$ to figure out those coefficients $A,B,C,F$ and $G$.}

\para
We'll do this calculation shortly. But it's a little bit fiddly, so first we should figure out what we're going to do with the answer when we get it. 

\para
{\co The coefficients $A$, $B$ and $C$ tell us the flux of particles. The incoming (or incident) flux is
%
\be J_{\rm inc} = |A|^2 \,\frac{\hbar k}{m}\nn\ee
%
The reflected flux is
%
\be J_{\rm ref} =  |B|^2\,\frac{\hbar k}{m}\nn\ee
%
Finally, the transmitted flux is
%
%
\be J_{\rm trans} = |C|^2\,\frac{\hbar k'}{m}\nn\ee
%
We define reflection and transmission coefficients
%
\be R = \frac{J_{\rm ref}}{J_{\rm inc}} \ \ \ {\rm and}\ \ \  T = \frac{J_{\rm trans}}{J_{\rm inc}}\nn\ee
%
These tell us what fraction of the incident beam is reflected and what fraction makes it over to the other side. Or, equivalently, the probability of reflection vs transmission. 
%
\be R +T =1\nn\ee
%
This is the statement that we don't lose any of the beam. } Nothing gets trapped at the potential. We'll see that this indeed holds for the solution for our bump potential.

\para

\alyssa{The below proof is included in the Tutorial for Lecture 9.} \href{https://drive.google.com/file/d/1RMR5gRvQ70qLkgVSp1QA3s2LBOTQ9TER/view?usp=drive_link}{[Lecture 9, Tutorial 1]}
{\cblue
I'm not sure that we want to prove that $R+T=1$ in the lecture, but here's the proof if we want it as tangent or a sideshow or whatever.

\para
Given a solution $\psi(x)$ to the Schrodinger equation, we can construct a conserved probability current
%
\be
J(x) = -i \frac{\hbar}{2m}\left(\psi^\star\frac{d\psi}{dx} - \psi \frac{d\psi^\star}{dx}\right)\nn\ee
%
For a time-independent solution, the conservation  of probability means thiso beys $dJ/dx=0$. This means that $J(x)$ is constant. (Mathematically, this is the statement that the Wronskian is constant.) For our scattering solution, to the left of the potential at $x\ra -\infty$, we have

\be J(x) &=& \frac{\hbar k}{2m}\Big[\left(A^\star e^{-ikx} + B^\star e^{+ikx}\right)\left(Ae^{ikx} - B e^{-ikx}\right) \\ &&\ \ \ \ \ \ \ \ \ \  +\  \left(A e^{ikx} + B e^{-ikx}\right)\left(A^\star e^{-ikx} - B^\star e^{+ikx}\right) \Big] =  \frac{\hbar k}{m}\left(|A|^2-|B|^2\right)\nn\ee
%
Meanwhile, as $x\rightarrow +\infty$, we have
%
\be J(x) = \frac{\hbar k}{2m}|C|^2\ \ \  {\rm as} \ \ x\rightarrow +\infty \nn\ee
%
Equating the two gives
%
\be |A|^2-|B|^2 = |C|^2 \ \ \ \Longrightarrow \ \ \ R+T = 1\nn\ee
%
So probabilities do what probabilities are supposed to do.}

\para
Ok. So we now know what we want to compute -- the reflection and transmission probabilities. And we're going to get them by relating $B$ and $C$ to the incoming particle density $A$. And we do that by insisting that the wavefunction we which we wrote down previously is continuous and its derivative is continuous. 


\para
{\co Insisting that $\psi(x)$ and $\psi'(x)$ are continuous (cts) at 
%
\be \psi(x)\ \mbox{cts at $x=-a$}\ \ \  &\Longrightarrow& \ \ \ \ \ Ae^{-ika} + Be^{+ika}=Fe^{+\eta a} + Ge^{-\eta a} \nn\\ 
\psi'(x)\ \mbox{cts at $x=-a$}\ \ \  &\Longrightarrow& \ \ ik(Ae^{-ika}-Be^{+ika}) = -\eta(Fe^{+\eta a} - Ge^{-\eta a})\ \nn \ee
%
Meanwhile, the matching conditions at $x=a$ are
%
\be \psi(x)\ \mbox{cts at $x=+a$} \ \ \ &\Longrightarrow& \ \ \ \ \ Ce^{ika} =Fe^{-\eta a} + Ge^{+\eta a} \nn\\ 
\psi'(x)\ \mbox{cts at $x=+a$} \ \ \ &\Longrightarrow& \ \ ikCe^{ika} = -\eta(Fe^{-\eta a} - Ge^{+\eta a})\nn\ee
%
}
We have four equations. If we can solve for  $C$ in terms of $k$ and $\eta$, since this will tell us how the transmitted flux depends on both the incoming momenta $k$ and the height of the potential $U$.  It's straightforward, if a little tedious. 

\para
{\cblue If we want to give some handholding separately, here's some handholding to help you along your way. First, use the second pair of equations to write
%
\be 2\eta Fe^{-\eta a} = (\eta-ik)Ce^{ika}\ \ \ \ {\rm and} \ \ \ 2\eta G e^{\eta a} = (\eta+ik)C e^{ika}\nn\ee
%
Next, use the first pair of equations to write
%
\be 2k Ae^{-ika} = F(k+i\eta)e^{\eta a} + G(k-i\eta)e^{-\eta a} \nn\ee
%
Now substitute the first two expressions for $F$ and $G$ in the equation above. A little bit of algebra then gives
%
\be \frac{D}{A} = \frac{2k\eta e^{-2ika}}{2k\eta  \cosh 2\eta a - i(k^2-\eta^2)\sinh 2\eta a}\nn\ee
%
}
After some algebra, we find that the transmission probability is
%
\be T = \frac{|D|^2}{|A|^2} =  \frac{4k^2\eta^2}{4k^2\eta^2  + (k^2+\eta^2)^2\sinh^2(2\eta a)}\ \nn\ee
%
There's a bunch of things to say about this. First, and most importantly,  {\co We see that $T\neq 0$. The particle has a non-zero probability to get through the barrier and out the other side, even though it didn't have enough energy to do so classically. This is called {\it quantum tunnelling}.}

\para
The next thing to notice is that the transmission probability is a complicated function of the momentum $k$ and the barrier height $U$. It's even more complicated than it looks because there's dependence buried in  $\eta= \sqrt{2m(U-E)/\hbar^2}$ and $E=\hbar^2 k^2/2m$. There's no way you were going to guess this answer without going through the calculation. That's common for these kinds of the scattering problems -- the way the probability changes with momentum is complicated. {\cblue should probably plot the transmission probability as a function of $k$ or $E$.}

\para
We can get a feel for what's going on if we look at very low energies. The first question is what does low energy mean. We would naively think it means $E\ll 1$ but that can't be true because it's not dimensionally consistent. It must be small compared to something with the right dimensions.

\para
{\co On dimensional grounds, at very low energies we have
%
\be U-E \gg \frac{\hbar^2}{2ma^2}\ .\ee
%
This means that $\eta a\gg 1$. Here the transmission probability can be written more easily as a function of the incoming energy $E$, 
%
%
\be T(E) \approx f(E)\,\exp\left(-\frac{4a}{\hbar}\sqrt{(2m(U-E))}\right)\ee
%
The coefficient in front is less important but, for what it's worth is $f(E)= 16k^2\eta^2/(k^2+\eta^2)^2$} (with both $k$ and $\eta$ viewed as functions of the energy).  The key feature is the exponential suppression - this is typical of any quantum tunnelling.


\para
{\cred This feels like a long lecture to me. Should think about what calculations we do and which ones are better left hidden.}

\alyssa{The practice problem for Lecture 9 will have students calculate the transmission probability of a potential, maybe a double inverted finite potential well.} \href{https://drive.google.com/file/d/1X81MIar2-dtrg2-fvZegfL0MSUJrpkav/view?usp=drive_link}{[Lecture 9 Problem Statement]}



\newpage
\noindent
{\Large \bf Section 3: Formalism}

\vskip 3mm

So far, we've postulated the Schrodinger equation and then gone ahead and solved it in a number of examples. We interpreted the results using a couple of axioms. The first was that the wavefunction -- or more precisely the mod-square of the wavefunction -- should be interpreted as the probability of finding the particle at some point. And the second was that the allowed values of $E$ on the right-hand side of the Schrodinger equation are the possible energies of the particle.

\para
Now it's time to put these axioms on a slightly firmer footing. In the next few lectures, I'm going to spell out the mathematical formalism that underlies quantum mechanics. This will apply much more generally than just a particle moving on a line, which is what we've discussed so far. We'll understand the right way to think about quantum systems. 

\para
I should warn you that this will involve a slight level of abstraction. Part of what we're going to learn is that quantum systems can look very different from classical systems. Not necessarily harder. In fact, there are examples -- as we'll see -- where the maths is significantly easier than anything that appears in classical mechanics. But it's also stranger. 

\para
{\co There are four facets of quantum mechanics: 
%
\begin{itemize}
\item states \vskip 2mm
\item observables \vskip 2mm
\item time evolution \vskip 2mm
\item measurement. \vskip 2mm
\end{itemize}
%
We will deal with each of them.}




\section{States}

We started these lectures by saying that the state of a quantum system is described by its wavefunction -- a complex-valued function whose norm-square is interpreted as the probability.

\para
That's true. But there's a better, more general, definition that we're going to use. Which is the following:

\para
{\co The state of a quantum system is given by a complex-valued vector.}

\para
Now that sounds very different. You might be thinking that a vector is the kind of thing you met in school  -- usually you write a row of numbers in a column, like the $(x,y,z)$ coordinates of some point. Or you might remember some definition along the lines of a vector is something with both magnitude and direction. A vector points in some direction and has some length. But part of what I want to tell you is that there's more to a vector than that. In particular, a function can be viewed as  vector.

\para
To understand this, we need to look at what a mathematician means by a vector. Roughly speaking, for them, a vector is any class of objects which satisfy two properties. First, you can multiply that object by a number and you get another object in the same class. And, second, you can add two of them together and get another object in the same class. 

\para
{\co At heart, this is just the principle of superposition. If $\psi(\bx,t)$ and $\phi(\bx,t)$ are both viable states of a system, then so too is any linear combination $\alpha \psi(\bx,t) + \beta\phi(\bx,t)$ where  $\alpha, \beta\in \mathds{C}$. Mathematically, this is the statement that the functions are examples of vectors -- add two of them together, multiplied with arbitrary coefficients,  and you get another function. They are infinite-dimensional vectors}.

\para
One of the things I want to do in this lecture is give you some intuition for why you should think of wavefunctions as infinite-dimensional vectors. 

\para
{\co Strictly, we say that they live in a {\it vector space}.}



\subsection{A Qubit}

There are also interesting quantum systems whose states are defined by finite-dimensional vectors. These are technically easier to deal with than infinte-dimensional vectors -- i.e. wavefunctions -- but conceptually less familiar. So throughout this set of lectures,  we will illustrate all the ideas with both wavefunctions and with finite-dimensional systems.

\para
{\co A {\it qubit} is the name given to a two-state system. In quantum mechanics, we often refer to these two states as $\ket{0}$ and $\ket{1}$, where $\ket{.}$ is referred to as a {\it ket} -- for now just think of it as a way of denoting the different states.}

\para
There are many different ways to build a qubit in the lab, but conceptually the simplest is to look at the spin of an electron. The electron spins in a way that is not dissimilar to the way the Earth spins but with a quantum difference. As we will see in later lectures, the spin in any given direction can only take certain values. And, for the electron, those values are either ``it spins this way'' or ``it spins the other way''. That's what we call spin-up or spin-down. So there are just two options, and those two options can be thought of as a qubit. 

\para
{\co One physical manifestion of a qubit is the spin of an electron. It can either be ``spin-up'' $\ket{\uparrow}$ ``spin-down" $\ket{\downarrow}$. So we can think of this as a qubit with
%
\be \ket{0} = \ket{\uparrow}\ \ \ {\rm and}\ \ \ \ket{1} = \ket{\downarrow}\nn\ee
%
The most general state of a qubit is 
%
\be \ket{\psi} = \alpha \ket{0} + \beta\ket{1}\nn\ee
%
with $\alpha,\beta\in \C$.} That's why qubits are more interesting than classical bits -- the qubit can be a little bit 0 and a little bit 1. We'll talk about the probability interpretation of $\alpha$ and $\beta$, and how to normalise these states, in the next lecture. 

\para
We can think of these coefficients of the qubit as the components of a vector.

\para
{\co We could also write this state as the vector $\vec{\psi}= (\alpha,\beta)$ or $\left(\begin{smallmatrix} \alpha\\ \beta \end{smallmatrix}\right)$.}  (This vertical and horizontal vector mean the same thing here.) Note that the vector $\vec{\psi}$ and the state $\ket{\psi}$ are subtly different: you can construct one from the other, but they're not quite the same thing. The componentns of the vector $\vec{\psi}$ are  the components of the state. That is, $\alpha$ tells you how much of the wavefunction is weighted on the first state and $\beta$  tells you how much of the wavefunction is weighted on  the second state. {\co The vector $(\alpha,\beta)$  is analogous to the wavefunction $\psi(x)$ which tells you how much of the wavefunction is weighted at the point $x$.} The qubit has two complex numbers; a particle on a line has an infinite number of complex numbers, one for each point $x$. That's the sense in which a function is an infinite-dimensional vector. 



\subsection{The Inner Product}

The formalism of quantum mechanics needs a piece of mathematical machinery called the inner product. This is a way to measure how close two states are to each other. 

\para
Let's first talk about finite dimensional vectors. Here the inner product is  the same thing as the scalar product between two vectors.  {\co If you have two vectors $\vec{u} = (u_1,u_2,u_3)$ and $\vec{v} = (v_1,v_2,v_3)$ with $u_i,v_i \in \R$, then the inner product between them is defined to be
%
\be \langle {u}\,|\,{v}\rangle = \vec{u}\cdot\vec{v} = \sum_{i=1}^3 u_i v_i = u_1 v_1 + u_2 v_2 + u_3 v_3 \nn\ee
%
There's an obvious extension to $N$-dimensional vectors
%
\be \langle {u}\,|\,{v}\rangle = \vec{u}\cdot\vec{v} = \sum_{i=1}^N u_i v_i\nn\ee
%
} I've indulged in some fancy notation here, using the angular brackets with the vertical line to denote the inner product. This notation turns out to be very useful in more advanced applications of quantum mechanics where it's called {\it Dirac notation}. Here I'm just introducing it  as an affection in the hope that it will  encourage familiarity for things that will come in later courses.  \alyssa{Dirac notation will be covered in a Tutorial for Lecture 10.} \href{https://drive.google.com/file/d/1rXOJsJt7k3oOCHnOf-k-5kjzRSTaIWdd/view?usp=drive_link}{[Lecture 10, Tutorial 1]]


\para
{\co For complex vectors, this needs a slight amendment. In this case, the inner product is
%
\be \langle {u}\,|\,{v}\rangle = \vec{u}^{\,\star}\cdot\vec{v} = \sum_{i=1}^N u_i^\star v_i\nn\ee
%
} i.e.  we take the complex conjugation of the first vector before we take the dot product.  {\co We do this so the inner product of any vector with itself is non-negative
%
\be \langle {u}\,|\,{u}\rangle = \vec{u}^{\,\star}\cdot\vec{u} = \sum_{i=1}^N |u_i|^2\ \geq 0\nn\ee
%
This means that we get to use the inner product to define the length of a vector $||\vec{u}|| = \sqrt{\langle u|u\rangle}$.

\para
As with wavefunctions, we will ask that any quantum state $\ket{\psi}$ is {\it normalised}. This means we require
%
\be \langle \psi |\psi \rangle = 1\nn\ee
%
So  our qubit $\ket{\psi} =\alpha \ket{0} + \beta \ket{1}$ is normalised if
%
\be |\alpha|^2 = |\beta|^2 = 1\nn\ee
%
} 
That's the story for  finite-dimensional vectors. Hopefully it's mostly familiar, except for the need to take the complex conjugation on the first vector so that the inner product of a vector with itself is always real and positive. 

\para
What about for infinite-dimensional vectors like the wavefunctions we've met so far. Here there is also a straightforward generalisation
%

\para
{\co The inner product between two wavefunctions $\phi(x)$ and $\psi(x)$ is
%
\be \langle \psi|\phi\rangle = \int dx\ \psi^\star(x)\,\phi(x) \nn\ee
%
Or, if this is a wavefunction for a particle moving in 3d, rather than 1d, we have
%
\be \langle{\psi}|{\phi} \rangle= \int d^3x\ \psi^\star(\bx)\phi(\bx)\label{inner}\ee
%
}
Again, we take the complex conjugate of the first function in the inner-product for the same reason as we did for complex vectors - it ensures that we can define the ``length" of the wavefunction. 

\para
{\co A wave-function $\psi(x)$ is normalised if $\langle \psi|\psi\rangle =1$ and is normalisable if$\langle \psi|\psi\rangle <\infty$.}

\para
We could also have the normalisble argument for finite-dimensional vectors but there it's obvious that any vector is normalisable provided the individual components less than infinity. For the wavefunction it's more subtle because each value of the wavefunction could be finite but the integral could give infinity.

\para
We can make some comments about that complex conjugation in the inner product. It means that {\co We have
%
\be \langle{\psi}|{\alpha\phi_1+\phi_2}\rangle = \alpha\langle{\psi}|{\phi_1}\rangle + \langle{\psi}|{\phi_2}\rangle\nn\ee
%
for any $\alpha \in {\bf C}$. But 
%
%
\be\langle{\alpha\psi_1+\psi_2}|{\phi}\rangle = \alpha^\star\langle{\psi_1}|{\phi}\rangle + \langle{\psi_2}|{\phi}\rangle\nn\ee
%
Also we have
%
\be \langle{\psi}|{\phi}\rangle^\star = \langle{\phi}|{\psi}\rangle\nn\ee
%
} These kind of simple observations will be used as we proceed. 



\para
{\co In  fancy mathematical language, a vector space, whether finite or infinite dimensional, with an inner product like this is called a {\it Hilbert space}, and denoted ${\cal H}$. That means that ${\cal H}$ is the space of all possible states.}
(In fact this statement is only almost true: there is an extra requirement in the case of an infinite dimensional vector space that the space is ``complete" which, roughly speaking, means that as your vector evolves it can't suddenly exit the space.)

\para
Hilbert space is the stage for quantum mechanics. The first thing you should do when describing any quantum system is specify its Hilbert space -- is it finite dimensional, is it infinite dimensionsal? And there can be subtleties. The Hilbert space for a particle moving on a line is not all complex functions on the line, but all normalisable complex functions. 

\para
Now we can be more precise...{\co A state in quantum mechanics is an element of the Hilbert space, which we write as $\ket{\psi}\in {\cal H}$ or  $\vec{\psi}\in {\cal H}$ or $\psi(x)\in {\cal H}$ depending on our preference. Two elements related by multiplication describe the same state, so $\ket{\psi}$ and $\lambda \ket{\psi}$ are equivalent for all $\lambda \in C$.  Alternatively, we can always insist that pick our state to be normalised so $\langle \psi|\psi\rangle =1$. The overall complex phase of the state is unphysical.}

\subsection{Problems}

{\cblue We could get them to compute the inner product between two different wavefunctions - say two Gaussians localised at the origin but with different widths; or two Gaussians separated in space. It's a little repetitive I guess but might be useful to give intuition for how close two different wavefunctions are from each other. \alyssa{The practice problem for Lecture 10 will have students calculate the inner product of two Gaussians with different widths localised at the origin.} \href{https://drive.google.com/file/d/1ZKVPQwYaagdpIA2RNrkCAwSgk6_wUmkk/view?usp=drive_link}{[Lecture 10 Problem Statement]}

\para
There is also a slightly fiddly proof that we might go over in a tangent. If you have two normalisable wavefunctions $\psi_1(x)$ and $\psi_2(x)$ then show that any finite linear combination is also normalisable and hence a state.} \alyssa{This proof will be covered in a Tutorial for Lecture 2.} \href{https://drive.google.com/file/d/15D_utjt-mdNEDmhK1WxOPFAAx-HqOa1U/view?usp=drive_link}{[Lecture 2, Tutorial 1]}



\newpage
\section{Observables}


In the last lecture, we've seen that the  quantum state of a system is described by a vector in Hilbert space. If our system is a particle moving in 1d or in 3d, that state is the same thing as the wavefunciton that we've been working with so far -- a complex-valued funciton over space. But it could also be something simpler -- but also more abstract -- like a finite dimensional vector. For example, a 2d vector if we're describing a qubit.

\para
The next question that we'd like to ask is: how is the information encoded in the state. And how do we extract that information? If we've got a wavefunction, how can we tell the position and momentum of a particle? Or, because this is quantum mechanics, it's better to ask how we extact the probability distribution over the position and momentum. We've already partially answered this with the Born rule that says the probability to find a particle at a given position is the mod-square of the wavefunction. But what's the analogous statement for momentum? 

\para
The technical name for these kinds of properties is an {\it observable}. An observable is a physical property that a system can have.

\para
It's best to first return to the classical world. {\co Classically, the state of a particle is determined by its position $\bx$ and momentum $\bp$.} (Recall that we need to specify both position and momentum because Newton's equation $F=ma$ is a second order equation -- so these are the two initial conditions needed to solve it.) {\co The classical observables are all possible functions of $\bx$ and $\bp$. This includes the position and momentum themselves, but also the angular momentum
%
\be {\bf L} = \bx\times \bp\nn\ee
%
and energy 
%
\be E = \frac{\bp^2}{2m} + V(\bx)\nn\ee
%
}

\para
That's the classical world. What about in quantum mechanics? Here things are a little more abstract and it's going to take a few lectures before we really get to grips with what's going on, and are able to build up the full story. But here's the starting point:

\para
{\co In quantum mechanics, all observables are represented by {\it operators} on the Hilbert space.} An operator is something that takes a state in the Hilbert space -- a vector, or a wavefunction -- and spits out another state. Or, more mathematically..

\para
{\co An operator $\hat{A}$ is a linear map
%
\be \hat{A}:{\cal H}\mapsto{\cal H} \ .\nn\ee
%
This means that if we start with a state $\ket{\psi}\in {\cal H}$ then after we act with the operator, we have a (typically) different state $\hat{A}\ket{\psi} \in {\cal H}$. Here ``linear" means that, for $\ket{\psi},\ket{\phi}\in {\cal H}$ and $\alpha\in\C$, we have
%
\be \hat{A}\big(\ket{\psi}+\alpha \ket{\phi}\big)=\hat{A}\ket{\psi} +\alpha\hat{A}\ket{\phi}\nn\ee
%
} The convention is that such operators are dressed with a hat, at least when you first meet them, so that you know what sort of  object you're dealing with. This convention tends to get dropped in more advanced renderings of the subject where you're just supposed to remember what everything is. But we're going to keep it for the present case.

\para
So this is going to take a bit of explaining. We would expect the position, or momentum, or energy of a state to be a number. But what we're seeing here is that an obervable in quantum mechanics isn't a number -- it's something that takes one state and spits out another. Over the next couple of lectures we're going to understand what that means. 


\para
But first we're going to look at some examples of operators. All our examples are going to be familiar -- well, familiarish -- but they're going to look rather different depending on whether we've got a finite dimensional Hilbert space (which means states are just vectors) or an infinite dimensional Hilbert space (which means states are wavefunctions). We'll look at each in turn. Then, for the next few lectures we're going to be looking at these two situations side by side -- finite dimensional Hilbert spaces where the states are just vectors, and infinite dimensional Hilbert spaces where the states are functions. 







\subsection{Operators as Matrices}

First let's consider finite dimensional Hilbert spaces. For example, we could start with the qubit. Previously we called the two states of the qubit $\ket{\uparrow}$ and $\ket{\downarrow}$ or $\ket{0}$ and $\ket{1}$, which were simply different names for the same thing. Here I'm going to be annoying and write the basis elements as $\ket{1}$ and $\ket{2}$ (which, as we'll see, is notation that will generalise). 

\para
{\co For a qubit, a general state is
%
\be \ket{\psi} &=& \alpha \ket{1} + \beta \ket{2}\nn \\
&=&\psi_1 \ket{1} + \psi_2 \ket{2}\ee
%
with $\alpha=\psi_1$ and $\beta=\psi_2\in \C$.} {\cblue Might be best to not introduce this new notation and just stick with $\ket{0}$ and $\ket{1}$.} To understand the action of any operator on a general state, it's enough to know how it acts on the basis elements $\ket{0}$ and $\ket{1}$. Then we can use the linearity to tell us how it acts more generally. {\co Suppose we have an operator that acts as
%
\be \hat{A} \ket{1} = A_{11} \ket{1} + A_{21} \ket{2}\ \ \ {\rm and}\ \ \ \hat{A}\ket{2} = A_{12}\ket{1} + A_{22} \ket{2}\nn\ee
%
with $A_{ij} \in \C$. Then, on the general state $\ket{\psi}$ linearity means 
%
\be \hat{A} \ket{\psi} &=& \psi_1 \hat{A}\ket{1} + \psi_2 \hat{A}\ket{2} \nn\\
&=& (A_{11} \psi_1 + A_{12}\psi_2) \ket{1} + (A_{21} \psi_1 + A_{22} \psi_2)\ket{2}\nn\ee
%
We write this as 
%
\be \hat{A} \ket{\psi} = \ket{\psi'} = \psi'_1 \ket{1} + \psi_2'\ket{2}\nn\ee
%
with 
%
\be \left(\begin{array}{c}\psi_1' \\ \psi_2' \end{array}\right)  = \left(\begin{array}{cc} A_{11}  & A_{12} \\ A_{21} & A_{22}\end{array}\right)
\left(\begin{array}{c}\psi_1 \\ \psi_2 \end{array}\right)\nn\ee
%
In other words, the operator $\hat{A}$ acts  on the coefficients $\psi_i$, $i=1,2$ as the matrix 
%
\be A = \left(\begin{array}{cc} A_{11}  & A_{12} \\ A_{21} & A_{22}\end{array}\right)\nn\ee
%
This also extends to an $N$-dimensional Hilbert space. On a finite dimensional Hilbert space, an operator is just a matrix. Suppose that we have the basis of orthonormal states $\ket{1}, \ket{2},\ldots \ket{N}$.}  (If this is too abstract, just think of these as a bunch of vectors with $\ket{1} = (1,0,\ldots,0)$ and $\ket{2} = (0,1,0,\ldots,0)$ and so on.) {\co We write these as $\ket{n}$ with $n=1,\ldots,N$. The operator $\hat{A}$ is can be thought of as the $N\times N$ matrix with elements 
%
\be A_{mn} = \bra{m} \hat{A}\ket{n}\nn\ee
%
}

\para
{\cblue I wonder if there's an exercise that can be developed here to try to get them to go through the general case below. The trouble is that they need to know what an orthonormal  basis is, and that hasn't really been explained at this point so it might be too abstract. Perhaps an explainer here? \alyssa{Orthonormal bases will be covered in a Tutorial for Lecture 11.} \href{https://drive.google.com/file/d/1XUNPlwBXt7FwrGpnGGXUjOZRcmQ9BdIB/view?usp=drive_link}{[Lecture 11, Tutorial 1]}


\para
Take a state $\ket{\psi}\in{\cal H}$ and expand in some orthonormal basis $\ket{n}$, so 
%
\be \ket{\psi} = \sum_n\psi_n\ket{n}\nn\ee
%
To start, let's suppose that our Hilbert space is ${\cal H}\cong \C^N$, so the  sum runs over $n=1,\ldots,N$. The  action of a linear operator $\hat{A}$ is fully determined by its action on the basis elements $\ket{n}$,
%
\be \ket{\psi}\mapsto \hat{A}\ket{\psi} = \sum_n\psi_n \hat{A}\ket{n}\nn\ee
%
We must be able to expand out this new state $\hat{A}\ket{\psi}$ in terms of our same basis elements. This will take the form
%
\be \hat{A}\ket{\psi} =\sum_n\psi'_n\ket{n}\nn\ee
%
for some components $\psi_n'$. We'd like to be able to express $\psi'_n$ in terms of the original $\psi_n$. This is achieved in terms of the {\it matrix elements} of $\hat{A}$. These are the coefficients of the operator $\hat{A}$ sandwiched between two basis states. We write \index{matrix elements}
%
\be A_{mn} = \bra{m}\hat{A}\ket{n}\ .\label{matrixa}\ee
%
The coefficients $A_{mn}$ are just a bunch of numbers. Because everything in quantum mechanics takes place over $\C$, each $A_{mn}$ is  a complex number. Take the inner product  with some basis element $\ket{m}$ and  use the orthonormality condition  to find
%
\be \psi_m' = \sum_n A_{mn}\psi_n\ .\ee
%
This tells us that the components of the state $\ket{\psi}$ transform as if multiplied by the matrix $A$ with components $A_{mn}$. }


\subsection{Operators as Things That Multiply and Differentiate}

That was the story when states are finite dimensional vectors. But what if they're wavefunctions. We've already said that wavefunctions can be thought of as infinite dimensional vectors, which makes it sound like the operators are infinite-dimensional matrices. And that sounds scary. 

\para
In fact, the operators {\it are} infinite dimensional matrices but, at least for our purposes, these are much less scary than they sound. It's just a not-very-useful name for things that you've met before: multiplication and differentiation.  Here we're going to give the most common physical examples of operators acting on wavefunctions

\para
{\co For a particle moving in $\R^3$, with a quantum state $\psi(\bx)$, we have the following operators}


%
\begin{itemize}
\item {\co There are three position operators, one for each direction. Combining these into a vector, they are  simply $\hat{\bx} = (\hat{x},\hat{y},\hat{z})$. These operators act on a wavefunction by multiplication
%
\be \hat{x} \psi(\bx) = x\psi(\bx)\ ,\ \ \ \hat{y}\psi(\bx) = y\psi(\bx)\ ,\ \ \ \hat{z}\psi(\bx) = z\psi(\bx)\ .\ee
% 
We will write this as $\hat{\bx}=\bx$. You can also take any function of $\bx$ and turn it into an operator. This acts in the obvious way, with $V(\hat{\bx})\psi(\bx) = V(\bx)\psi(\bx)$.}  \vskip 3mm
%
\item  {\co The momentum operator for a particle moving in 3d is
%
\be \hat{\bp} = -i\hbar \nabla\ee
%
or, written out longhand,
%
\be \hat{p}_x = -i\hbar\ppp{}{x}\ ,\ \ \ \hat{p}_y = -i\hbar\ppp{}{y}\ ,\ \ \ \hat{p}_z = -i\hbar\ppp{}{z}\ .\ee
%
} This is an example of what's called a differential operator. 
{\co These should be thought of in terms of their action on the wavefunctions so, for example,
%
\be \hat{p}_x\psi(\bx) = -i\hbar\ppp{\psi}{x}\ .\ee
%
} We've seen hints previously that momentum is somehow related to differentiation. This is the precise statement. We'll see in the next lecture how we actually extract the information about the momentum of a particle from this statement.

You might also wonder {\it why} momentum is related to differentiation. There is an answer to this, but it's a little involved -- it's related to Noether's theorem and the idea that momentum comes from translation symmetry. It's a bit advanced for this course so, for now at least, you'll have to take this as an axiom of quantum mechanics -- the momentum operator is differentiation.

\vskip 2mm\item Now we can build other observables by starting from the classical expressions and just replacing the positions and momenta with their quantum counterparts. {\co We have the angular momentum operator
%
\be \hat{\bf L} =  \hat{\bx} \times {\bf bp} = -i\hbar \,\hat{\bx} \times \nabla\ .\ee
%
} We're going to look at this much more closely later in these lectures when we come to solve the hydrogen atom -- it turns out there's quite a lot to say about this. 

{\co Finally the energy operator also comes from replacing $\bx$ with $\hat{\bx}$ and $\bp$ with $\hat{\bp}$ in the classical expression giving
%
\be \hat{H}  = -\frac{\hbar^2}{2m}\nabla^2 + V(\hat{\bx})\ .\label{hampoland}\ee
%
The energy operator is the only one that gets a special name, with $\hat{H}$ for Hamiltonian rather than $\hat{E}$ for energy.} This is because of the special role it plays in time evolution of the state through the Schrodinger equation. 
\end{itemize}
%

So we have a bunch of examples of different operators. The next question that we have to ask is what do we do with these things?!

\para
{\cblue I'm not sure if there are really problems that we can ask here, or explainers. It might be easier to find something for the next lecture.}



\newpage

\section{Eigenstates and Eigenvalues}

We've seen that states in quantum mechanics are vectors in a Hilbert space -- and one example of this is the wavefunction which is really just an infinite dimensional vector.

\para
And we've also seen that observables are operators on these states. But we haven't really explained what that means.  It's all very well and good saying things like ``the momentum of the particle is related to the operator $\hat{\bp} = -i\hbar\nabla$", but what does this mean in practice?  What does this have to do with the idea of momentum that we know and love from classical mechanics? How do we use this to figure out what momentum a given wavefunction has?

\para
The answer to this lies in some properties of the operators, whether they're matrices or differential operators. 

\para
{\co
Given an operator $\hat{A}$ acting on ${\cal H}$, the {\it eigenvalue equation} is
%
\be \hat{A}\ket{\psi} = \lambda\ket{\psi}\ \label{eigen1}\ee
%
for some $\lambda\in \C$.} To solve this equation, we need to find states $\ket{\psi}\in{\cal H}$ such that, when we hit them with $\hat{A}$, we get the same state, up to the constant $\lambda$.  {\it The constant $\lambda$ is called the {\it eigenvalue}, while $\ket{\psi}$ is the {\it eigenstate}. Typically, solutions to this equation will exist only for certain values of $\lambda$.} 

\para
We can look at some examples. When we have a finite dimensional Hilbert space, we've seen that operators are just the same thing as matrices. So we can look at the eigenvalue equation for matrices.  {\co For example, consider the matrix
%
\be A  = \frac{1}{5}\left(\begin{array}{cc} 14 & -2 \\ -2 & 11 \end{array}\right)\nn\ee
%
The eigenvalues obey
%
\be A \vec{\psi} = \lambda \vec{\psi}\nn\ee
%
}In solving this, we're looking for both the eigenvalue $\lambda$ and the eigenvector $\vec{\psi}$. If you write down the most general eigenvector, plug it in, then you'll find you've got a pair of coupled linear equations to solve. {\co We find two solutions
%
\be \vec{\psi}_1 = \left(\begin{array}{c} 1 \\ 2\end{array}\right) \ \ \ {\rm with}\ \ \ \lambda_1 = 1\nn\ee
%
and 
\be \vec{\psi}_2 = \left(\begin{array}{c} 2 \\ -1\end{array}\right) \ \ \ {\rm with}\ \ \ \lambda_2 = 3\nn\ee
%
}
This is typical. For the kind of operators that we'll need for quantum mechanics -- something called Hermitian operators which we'll define below -- he number of solutions to the eigenvalue problem will always be equal to the dimension of the matrix. If we have Hilbert space of dimension $N$, then the operators are $N\times N$ matrices and there will be $N$ eigenvalues and $N$ eigenvectors. 

\para
{\co The eigenvalues are unique. The eigenvectors are not -- we can always multiply the eigenvectors by an overall constant and that will also solve the equation, giving the same eigenvalue.}

\para
{\cblue We could probably do with an explainer on eigenvectors and eigenvalues for matrices. There's a lot that could go on here -- perhaps explaining the characteristic polynomial and how that tells you the eigenvalues. It's simple to cook up questions where students figure out the eigenvectors and eigenvalues of simple matrices themselves.} \alyssa{How to find eigenvectors from eigenvalues will be covered in a Tutorial for Lecture 12, and the practice problem for Lecture 12 will require students to figure out eigenvectors and eigenvalues of matrices.} \href{https://drive.google.com/file/d/1xY1QxbzBbt0AyL5X4bBdTMAyDvIMozNq/view?usp=drive_link}{[Lecture 12, Tutorial 1]} \href{https://drive.google.com/file/d/1JFHpnr3PnQu0lyddgoqes3x1bO4-_jMk/view?usp=drive_link}{[Lecture 12 Problem Statement]}


\para
We can play the same game with differential operators. {\co For an  operator $\hat{A}$ acting on some class of functions, the eigenvalue equation is
%
\be \hat{A} \psi(\bx) = \lambda \psi(\bx)\ \ \mbox{for some function $\psi(\bx)$}\nn\ee
%
The corresponding function $\psi(\bx)$ is called the {\it eigenfunction}. 

\para
The collection of all eigenvalues is called the {\it spectrum} of the operator $\hat{A}$.}

\para
Now we can take the first important step in interpreting these operators. Given any operator -- corresponding to some physical observable -- if you make a measurement t then the answer  you get will be one of the eigenvalues of $\hat{A}$. This is one of the key axioms of quantum mechanics. 


{\co 
%
\be \boxed{\mbox{
The outcome of any measurement of $\hat{A}$ lies in the spectrum of $\hat{A}$.}}\nn\ee
}
We've already implicitly met this idea earlier in this course. {\co One example of this can be seen in time-independent Schrodinger equation is simply the eigenvalue equation for the Hamiltonian
%
\be \hat{H}\psi = E\psi\nn\ee
%
The spectrum of the Hamiltonian determines the possible energy levels of the quantum system.} In most cases this spectrum is discrete, like the harmonic oscillator. That's what the quantum in quantum mechanics means. In other simple, but mildly pathological, cases the spectrum is continuous, like a particle on a line.

\para
All physical observables have an associated eigenvalue equation. {\co For the momentum operator the eigenvalue equation is
%
\be \hat{\bp} \psi = \bp\psi\ \ \ \Longrightarrow\ \ \ -i\hbar\nabla\psi = \bp \psi\nn\ee
%
}
Eigenvalue equations in quantum mechanics often look like the first equation above -- operator with a hat acting on $\psi$ equals eigenvalule (which has the same name as the operator but without the hat) times $\psi$. It means you sometimes have to squint to see that the equation is something other than a tautology -- the difference is just that hat on the left. In later courses, when we tend to drop the hat on operators, it gets really confusing!




\para
It's easy to solve the momentum eigenvalue equation.  {\co The (unnormalised) momentum eigenstates are
%
\be \psi = e^{i\bk \cdot \bx}\ \ \ \Longrightarrow\ \ \ \ \hat{\bp}\psi = \hbar\bk \psi\nn\ee
%
These are called {\it plane wave states}} on account of the fact that the wave is in the $\bk$-direction and the function is uniform in the plane perpendicular to $\bk$. 
{\co The associated momentum eigenvalue is $\bp = \hbar \bk$.} As we mentioned previously, this is the grown-up version of de Broglie's relationship between momentum and wavelength -- or here, the wavevector ${\bf k}$. As we've seen previously, this is an example where the wavefunctions are not normalisable.



\para
What about the position operator? We'll briefly revert to a one-dimensional system. {\co A position eigenstate must then obey
%
\be  \hat{x}\psi(x) = x\psi(x) = X\psi(x)\nn\ee
%
where $X\in \mathds{R}$ is the eigenvalue.} This is a strange-looking equation. The first $\hat{x}$ is the operator, the second $x$ is what the operator does -- it multiplies by $x$ -- and here we're using $X$ to denote the eigenvalue. Also the equation itself is strange. It tells us that we must  have $\psi(x)=0$ whenever $x\neq X$, but the equation does leave open the possibility that $\psi(X)\neq 0$. The object that has this property is slightly odd when you first meet it {\co The solution  is the Dirac delta function,
%
\be \psi(x) = \delta(x-X)\nn\ee
% 
This is defined to be 
%
\be \delta(x-X) = \left\{\begin{array}{cc} 0 \ \ & x\neq X\\ \infty \ \ & x=X\end{array}\right.\nn\ee
%
} In other words, this function is an infinite spike situation at $x=X$. Strictly speaking, this isn't a function. It's too singular for that. It's an example of what mathematicians call a generalised function or a distribution. There's a rigorous way to think about these things but, as is often the case in physics, we'll get by without being rigorous -- we'll use it carelessly and see that we get senisble results. 

\para
Although the delta-function looks scary, it's actually really easy to work with. This is because it usually only appears inside integrals and has the property that it kills the integral. In particular, {\co for any function $f(x)$, we define
%
\be \int_{-\infty}^{+\infty} f(x) \delta(x-X) \,dx = f(X)\nn\ee
%
}
That makes sense if you think about it -- the delta-function vanishes everywhere except for $x=X$, so you don't get any contribution to the integral except at $x=X$. We choose the delta-function to spit out the value of the function $f(x)$ evaluated at the point where the delta-function spikes. {\co In particular, we have
%
\be \int_{-\infty}^{+\infty} \delta(x-X) \, dx=1\nn\ee
%
In other words, the area under that infinite spike is equal to one.}


\para
The momentum eigenstates are those plane waves, and we'e seen previously that they're non-normalisable, at least for a particle on the line. The position eigenstates are delta-functions and they too are non-normalisable. {\co The position eigenstates are $\psi(x) = \delta(x-X)$} (which describes a particle localised at point captial $X$) {\co and the normalisation of this state is
%
\be \int_{-\infty}^{\infty} |\psi|^2\,dx = \int_{-\infty}^{\infty} |\delta(x-X)|^2\,dx  =\delta(0) = \infty\nn\ee
%
} Here, in the second equality we've used one of the delta functions to kill the integral, and take $f(x) = \delta(x-X)$. Then we evaluate that where it gives infinity. {\co So the position eigenstates are, like momentum eigenstates, non-normalisable}. 


\para
This is just the latest way of telling the slightly annoying story that has plagued us throughout the opening lectures -- the two simplest operators for a particle, $\hat{\bp}$ and $\hat{\bx}$ are kind of screwy in the sense that their eigenfunctions aren't normalisable. It's something that doesn't happen for most operators. 


\para
{\cred I originally had this lecture together with the Hermitian operator lecture, but I think it's probably best to split them up. Otherwise it's a very long lecture. And there's different ideas going on in the two of them. Still, this one might now be on the short side.}


\newpage
\section{Hermitian Operators}


We've seen that physical observables are related to operators. And that the eigenvalues of the operator are the possible values that observable can take. There's still a lot that we haven't explained -- like how the values of observables are related to the state of the system -- we'll get there in the next few lectures. But first there's something more we need to say about the operators. Because it's not true that any old operator can be thought of as a physical observable. Instead, observables correspond to very special kind of operator that is called {\it Hermitian}. 


\para
To explain this, we first introduce the concept of an adjoint operator. {\co Given an operator $\hat{A}$, its {\it adjoint} $\hat{A}^\dagger$ is defined by the requirement that the following inner-product relation holds
%
\be \bra{\psi}{\hat{A}}\ket{\phi} =   \bra{\phi}\hat{A}^\dagger\ket{\psi}^\star\nn\ee
%
for all states $\ket{\phi},\ket{\psi}\in{\cal H}$. } Note, in particular, the complex conjugation on the right-hand side. {\co Written in terms of integrals of wavefunctions, this means
%
\be \int d^3x\ \psi^\star \hat{A}\phi = \int d^3x\ (\hat{A}^\dagger \psi)^\star \phi\nn \ee
%
for all $\phi(\bx)$ and $\psi(\bx)$. 

\para
From this definition, the product of two operators obeys $(\hat{A}\hat{B})^\dagger = \hat{B}^\dagger \hat{A}^\dagger$.

\para
An operator is said to be {\it Hermitian} or {\it self-adjoint} if\index{self-adjoint operator}
%
\be \hat{A} = \hat{A}^\dagger \nn\ee
%
} In fact, there's a tiny bit of a lie in that previous statement. The terms ``Hermitian" and ``self-adjoint" are only synonymous for finite dimensional Hilbert spaces. When we're dealing with infinite dimensional Hilbert spaces of functions, there is a subtle distinction between them that we're just going to sweep under the rug here.

\para
{\co Importantly, observables are associated to Hermitian operators.} 

\para
This should be combined with our earlier, boxed, statement that the result of any measurement is contained in the spectrum of the operator.


\para
Again, we can think about this idea for finite dimensional Hilbert spaces - where the operators are just matrices -- or for infinite dimensional Hilbert spaces where the operators are things like multiplication or differentiation. Let's do the finite dimensional case first.

\para
{\co For a finite dimensional Hilbert space, a Hermitian operator $\hat{A}$ obeys
%
\be A_{mn} = \bra{m}\hat{A}\ket{n}= \bra{n}\hat{A}\ket{m}^\star = A_{nm}^\star\nn\ee
%
This is what we mean for a matrix to be Hermitian. The {\it Hermitian conjugate} is defined to be
%
\be A^\dagger =(A^\star)^T\nn\ee
%
} That is, you take the transpose of the matrix and take the complex conjugate of all elements. Then the matrix is hermitian if the hermitian conjugate is the same as the original. {\co A Hermitian matrix obeys
%
\be A^\dagger = A\nn\ee
%
}
For an infinite dimensional Hilbert space, we've already got our two favourite operators -- the position and momentum operators. We can just check that both are Hermitian using the definition above. 
{\co For the position operator, we have
%
\be \bra{\psi}\hat{x}\ket{\phi} = \int dx\ \psi^\star \, x\phi = \int dx\ (x \psi)^\star \phi = \bra{\phi}\hat{x}\ket{\psi}^\star\nn\ee
%
So we learn that $\hat{x} = \hat{x}^\dagger$ as promised.} Here the calculation was trivial: it is just a matter of running through definitions. It really follows simply because the position operator $\hat{x}$ multiplies the wavefunction by the position $x$ and that's a real number. 


The calculation for the momentum operator is marginally less trivial. {\co For the momentum operator in 1d, we have
%
\be \bra{\psi}\hat{p}\ket{\phi} &=& \int dx\ \psi^\star \left(-i\hbar \frac{d\phi}{dx}\right)\nn\\ &=&   \int d^3x\ i\hbar (\frac{d\psi^\star}{dx})  \phi\nn\\ &=& \int d^3x\  (-i\hbar\frac{d\psi}{dx})^\star  \phi= \bra{\phi}\hat{p}\ket{\psi}^\star\ .\ee
%
} In the second line we integrated by parts, picking up one minus sign, and in the third line we took the factor of $i$ inside the brackets with the complex conjugation giving us a second minus sign. The two minus signs cancel out to give $\hat{p} = \hat{p}^\dagger$. This is one way to see why the momentum operator necessarily comes with a factor of $i$. {\co So we have $\hat{p} = \hat{p}^\dagger$. The same is true for the 3d momentum operator $\hat{\bp} = -i\hbar\nabla = \hat{\bp}^\dagger$.}

\para
Note that, when we integrated by parts, we threw away the total derivative. This is justified on the grounds that the states $\ket{\psi}$ and $\ket{\phi}$ are normalisable and so asymptote to zero at infinity. 

\para
The Hermiticity  of the Hamiltonian $\hat{H}$ then follows immediately  because it's a real function of $\hat{\bx}$ and $\hat{\bp}$. Alternatively, you check explicitly that the Hamiltonian is Hermitian by doing the kind of integration-by-parts calculation that we saw above.




\subsection{Properties of Hermitian Operators}


There are reasons why physical observables are defined by Hermitian operators rather than any other kind. It's because they come with nice mathematical properties that translate into nice physical properties.  Here's the first

\para
{\co  An eigenstate of a Hermitian operator obeys
%
\be \hat{A}\ket{\phi_n} = \lambda_n\ket{\phi_n}\nn\ee
%
where $\lambda_n$ is the eigenvalue.} We've shifted notation slightly -- we're now labelling the different eigenstates and eigenvalues now labelled by $n=1,\ldots$. We will ultimately see that the range of this label runs up to the dimension of the Hilbert space...either finite or infinite depending on the situation. 

\para
{\co For a Hermitian operator $\hat{A}$, the eigenvalues and eigenstates have the following two properties:
%
\begin{itemize}
\item The eigenvalues $\lambda_n$ are real.  \vskip 3mm
\item If two eigenvalues are distinct $\lambda_n\neq \lambda_m$, then their corresponding eigenstates are orthogonal: $\bra{\phi_n}{\phi_m}\rangle=0$. 
\end{itemize}
%
}
Recall our important boxed statement from earlier: the outcome of any measurement of $\hat{A}$ is given by one of the eigenvalues of $\hat{A}$. The first of the properties above ensures that, happily, the result of any measurement is guaranteed to be  a real number. That's a good thing. If you go into a lab, you'll see dials and needles and rulers and other complicated things. All of them measure real numbers, never complex numbers.


\para
{\co Both of these properties are simple to prove.  If $\hat{A}$ is Hermitian, then we have
%
\be \bra{\phi_n}\hat{A}\ket{\phi_m} = \bra{\phi_m}\hat{A}\ket{\phi_n}^\star\ \ \ &\Longrightarrow&\ \ \  \lambda_m \bra{\phi_n}{\phi_m}\rangle = \lambda_n^\star\bra{\phi_m}{\phi_n}\rangle^\star\nn\\
&\Longrightarrow&\ \ \ (\lambda_m -\lambda_n^\star)\,\bra{\phi_n}{\phi_m}\rangle=0\nn\ee
%
If we set  $n=m$, then we find that $\lambda_n=\lambda^\star_n$ so that the eigenvalue is real. Alternatively, pick distinct eigenvalues $\lambda_n\neq \lambda_m$ to show that the distinct eigenstates are orthogonal, with $\bra{\phi_n}{\phi_m}\rangle=0$.

\para
If we're dealing with a finite dimensional Hilbert space, this just means that the corresponding eigenvectors are orthogonal. For wavefunctions, it means eigenstates obey
%
\be  \int d^3x\ \phi_n^\star\phi_m = 0 \ \ \ {\rm if}\ m\neq n\nn\ee
%
It may well be that an operator has several different eigenstates all with the same eigenvalue. In this case that eigenvalue is said to be {\it degenerate}, with the {\it degeneracy} equal to the number of linearly independent eigenstates. Conversely, an operator with distinct eigenvalues is said to have a {\it non-degenerate} spectrum. }


\para
The calculation above doesn't tell us anything about the eigenstates when an eigenvalue is degenerate. In this case, it's still possible to pick orthogonal eigenstates. We won't prove this here. It's straightforward to show for the case of a finite dimensional Hilbert space, just by picking an orthogonal basis for the space spanned by the eigenvectors. 


\para
{\co The eigenstates are orthogonal.} (Or can be taken to be orthogonal in the case of a degenerate spectrum.) {\co If we can normalise the eigenstates, then they will be {\it orthonormal}, meaning
%
\be \bra{\phi_n}{\phi_m}\rangle = \delta_{mn}\nn\ee
%
}

\para
{\cblue Things are getting abstract here without too much grounding. That's going to continue for the next lecture as well. Hopefully this will be rectified when we get to the qubit in the lecture after that. We might think about brining the qubit forwards -- or perhaps it's better to get through the abstract stuff and then exhibit everything using the qubit as an example.}


\subsection{Eigenbases}

There's one last property of Hermitian operators that will be important. Here we're going to state it without giving a proof. 

\para
{\co  The eigenstates of any Hermitian operator are {\it complete}. This means that we can expand any state $\ket{\psi}$ in terms of eigenstates of a given operator $\hat{A}$
%
\be \ket{\psi} = \sum_n a_n\ket{\phi_n} \ \ \ \mbox{where $\hat{A}\ket{\phi_n} = \lambda_n\ket{\phi_n}$}\nn\ee
%
for some complex coefficients $a_n$.} 


\para
This means that every observable comes with its own preferred orthonormal basis, which are the eigenstates of the corresponding operator. We will soon see that a great deal of the physics of quantum mechanics is lurking in the different ways in which we can expand a state, by choosing bases associated to different operators. 

\para
Let's look at some simple examples. We'll start with a two-dimensional Hilbert space where, as we've seen, the observables are just Hermitian $2\times 2$ matrices. {\co For a 2-dimensional Hilbert space, we could consider the observable 
%
\be A = \left(\begin{array}{cc} 2 \ & 0 \\ 0\  & 7\end{array}\right)\nn\ee
%
} I chose this matrix just because it's simple. {\co The eigenvectors are just $\bphi_1 = (1,0)$ and $\bphi_2 = (0,1)$.} They have eigenvalues 2 and 7 respectively, but that's not super important for our immediate purposes. What we really care about are the eigenvectors. {\co The most general state is $\bpsi = (\alpha,\beta)$ and can be (trivially) written as sum of the two eigenvectors
%
\be \left(\begin{array}{c} \alpha \\ \beta\end{array}\right) = \alpha \bphi_1 + \beta \bphi_2\nn\ee
%
} But we could also look at an entirely different operator. For example, {\co consider the matrix
%
\be B  = \frac{1}{5}\left(\begin{array}{cc} 14 & -2 \\ -2 & 11 \end{array}\right)\nn\ee
%
We looked at this before: it has eigenvectors $\tilde{\bphi}_1 =\ft{1}{\sqrt{5}}(1,2)$ and $\tilde{\bphi}_2 =\ft{1}{\sqrt{5}}(2,-1)$.} (Those factors of $1/\sqrt{5}$ are there to normalise the eigenvectors.) {\co Any vector can be expanded as the sum of these two: 
%
\be \left(\begin{array}{c} \alpha \\ \beta\end{array}\right) = \tilde{\alpha} \tilde{\bphi}_1 + \tilde{\beta} \tilde{\bphi}_2\nn\ee
%
with
%
\be \tilde{\alpha} = \frac{1}{\sqrt{5}} (\alpha  + 2\beta) \ \ \ {\rm and}\ \ \ \tilde{\beta} = \frac{1}{\sqrt{5}}(2\alpha -\beta)\nn\ee
} 
{\cblue check I did this right!} But this is a general story. Take any Hermitian operator -- it could be a finite dimensional matrix or infinite dimensional something else -- then it has a set of eigenvectors. And you can expand any state you like in terms of those eigenvectors.

\para
{\co If $\ket{\phi_n}$ are the eigenvectors of $\hat{A}$ then we can expand any state as
%
\be \ket{\psi} = \sum_n a_n\ket{\phi_n}\nn\ee
%
for some coefficients $a_n\in \C$. For wavefunctions, that means we can write
%
\be \psi(\bx) = \sum_n a_n \phi_n(\bx)\nn\ee
%
where $\phi_n$ are eigenfunctions of some operator.}

\para
We've already seen some examples of eigenfunctions. For the first few lectures we were solving the time-independent Schrodinger equation, finding the energy eigenvalues and the energy eigenstates. Those eigenstates are the $\phi_n(\bx)$. For example, when we looked at the harmonic oscillator, we found the energy eigenstates were expressed in terms of Hermite polynomials. But those polynomials form a basis -- that means it's possible to express any function as a sum of those different Hermite polynomials. 


\para
{\co If you know both $\ket{\psi}$ and the  normalised eigenstates $\ket{\phi_n}$, then it's simple to get an expression for the coefficients $a_n$. We have
%
\be \bra{\phi_n}{\psi}\rangle = \sum_m a_m \bra{\phi_n}{\phi_m}\rangle = \sum_m a_m \delta_{mn} = a_n\nn\ee
%
where, in the second equality, we've used orthonormality of the eigenstates. If we have wavefunctions, then this becomes the integral
%
\be a_n = \int d^3x\ \phi_n^\star(\bx)\,\psi(\bx)\ .\ee
%
Relatedly, the norm of a state $\ket{\psi}$ is
%
\be ||\psi||^2 = \bra{\psi}{\psi}\rangle = \sum_{m,n} a_n^\star a_m\bra{\phi_n}{\phi_m}\rangle= \sum_n|a_n|^2\ \ee
%
where we have again used the orthonormality condition. We see that a given state is normalised only if $\sum_n|a_n|^2 = 1$.}



\subsection{Momentum Eigenstates and the Fourier Transform}

There is one very important application of the idea that any function can be expressed as a sum of eigenstates of  a Hermitian operator. Consider the momentum operator. We've seen that we have a problem because the eigenstates are plane waves which, for a particle on a line, are not normalisable. So let's get around this by thinking of a particle moving on circle.

\para
{\co For a particle moving on a circle of radius $R$, all wavefunctions obey $\psi(x + 2\pi R) = \psi(x)$.  the momentum eigenstates obey
%
\be \hat{p}\phi_n(x) =-i\hbar\frac{d\phi_n}{dx} = p_n\phi_n(x)\nn\ee
%
We saw in a previous lecture that  the eigenstates and eigenvalues are
%
\be \phi_n = \sqrt{\frac{1}{2\pi R}}\,e^{inx/R} \ \ \ {\rm and}\ \ \ p_n =\frac{\hbar n}{R} \ , \   n\in \Z\nn\ee
%
Note: these eigenstates are normalised. You can check that the eigenstates are orthonormal, obeying
%
\be \bra{\phi_n}{\phi_m}\rangle = \frac{1}{2\pi R}\int_0^{2\pi R} dx\ e^{i(m-n)x/R} =\delta_{mn} \nn\ee
%
} {\cblue Might need an explainer to do this integral} \alyssa{There will be a Tutorial for Lecture 13 to help with this integral.} \href{https://drive.google.com/file/d/1EExmeqJf8ICMpsf0w5eeV5qmQL_5w-Lx/view?usp=drive_link}{[Lecture 13, Tutorial 2]}

\alyssa{Fourier decomposition will be covered in a Tutorial for Lecture 13.} \href{https://drive.google.com/file/d/1n-YQ8dFd1MKPnm-bguGlogegpUygVlsI/view?usp=drive_link}{[Lecture 13, Tutorial 1]}
{\co Any periodic complex function on the circle, such that $\psi(x) = \psi(x+2\pi R)$, can be expanded in terms of these eigenstates, meaning that
%
\be \psi(x) = \sqrt{\frac{1}{2\pi R}}\, \sum_{n\in \Z}a_n e^{inx/R}\nn\ee
%
This is known as the Fourier expansion of the function} -- we're writing any function $\psi(x)$ as a bunch of different cos and sine waves.  {\co The Fourier coefficients $a_n$ are given by 
%
\be a_n = \int_0^{2\pi R} dx\ \phi_n^\star\psi = \sqrt{\frac{1}{2\pi R}}\,\int_0^{2\pi R} dx\ e^{-inx/R}\psi(x)\nn\ee
%
}

\para
We can also think about what happens to the particle on the line. {\co For a particle on a line, the momentum eigenstates are labelled by a continuous variable $k\in \R$} (rather than a discrete variable). {\co 
%
\be \phi_k(x) = \sqrt{\frac{1}{2\pi}} e^{ikx}\nn\ee
%
As we've seen several times, these states  are non-normalisable. Nonetheless, there is a version of the orthogonality relation  which is now
%
\be \bra{\phi_k}{\phi_{k'}}\rangle = \frac{1}{2\pi} \int dx\ e^{i(k-k')x} = \delta(k-k')\nn\ee
%
where $\delta(k)$ is the Dirac delta function. } {\cblue Expand the previous explainer to include this perhaps}. (That is that for different values of $k$, the states are orthgonal, but the norm of any single state is infinite -- that's the spike in the delta function.)  {\co We can expand a function $\psi(x)$ in these eigenstates, with the sum over states now replaced by an integral
%
\be \psi(x) = \sqrt{\frac{1}{2\pi}} \int_{-\infty}^{+\infty} dk\ \tilde{\psi}(k)\,e^{ikx}\nn\ee
%
The Fourier coefficients $a_n$ are now promoted to full function $\tilde{\psi}(k)$. We can use the orthogonality condition to invert the expression above and write $\tilde{\psi}(k)$ in terms of the original function $\psi(x)$, 
%
\be \tilde{\psi}(k) =  \sqrt{\frac{1}{2\pi}} \int{-\infty}^{+\infty} dx\ {\psi}(x)\,e^{-ikx}\nn\ee
%
The relations are the Fourier transform and inverse Fourier transform of  a function.}


\para
{\cred This is quite a long lecture I think...the next one will be much shorter, but it's probably more important to keep a coherent set of ideas in the same lecture and I think this one does this -- it's more on the maths and the next one on physics.}

\para
{\cblue There's a lot of maths going on in this lecture and probably many exercises/explainers that we could develop to get the message across.} \alyssa{Right now, there are two planned Tutorials for Lecture 13: one on Fourier Decomposition and one on delta functions. Is there another topic you think would work better in a Tutorial than in the lecture as written?}


\para
\alyssa{The practice problem for Lecture 13 will have students check the orthonormality of the eigenstates of a particle in a circle.} \href{https://drive.google.com/file/d/13x8lddNRo6ANECz4AoY_vG4KdxirXCWn/view?usp=drive_link}{[Lecture 13 Problem Statement]}


 

\newpage
\section{Time Evolution}

For the last few lectures, we spent some time building up some fairly abstract mathematics. We've understood that states in quantum mechanics are vectors in Hilbert space. In some situations,  these are finite dimensional vectors, of the kind we met in school. In other situations the states are functions -- the wavefunction of the particle.

\para
We've learned that observables are a special kind of operator on these states -- Hermitian operators. The eigenvalues of these operators are real and should be viewed as the possible result of any experiment -- we're going to get to how you relate these eigenvalues to measurements in the next lecture. Meanwhile, these operators also come with a special set of states. These are the eigenstates associated to the operator.

\para
In this lecture, we're going to look at one application of this idea of eigenstates. Of all the observables -- all the Hermitian operators -- in quantum mechanics, there is one that holds an elevated importance. It's the observable associated to energy, called the Hamiltonian. And it's important because the Hamiltonian tells us how a state changes with time. 



\para
We met the time-dependent Schrodinger equation right back at the beginning of these lectures. We can rewrite it  in our new, slightly fancy notation. {\co If we write the state as $\ket{\psi(t)}$, the time-dependent Schrodinger equation is
%
\be i\hbar\ppp{\ket{\psi(t)}}{t} = \hat{H}\ket{\psi(t)}\nn\ee
%
The energy eigenstates obey
%
\be \hat{H} \ket{n} = E_n\ket{n}\nn\ee
%
} This is the kind of equation that we spent a long time solving in various examples for a particle moving in one dimension. {\co If our system sits in an energy eigenstate at time $t=0$, so $\ket{\psi(0)} = \ket{n}$, then the solution to the time-dependent Schrodinger equation is simple. It is
%
\be \ket{\psi(t)} = e^{-iE_n t\hbar} \ket{n}\ \ \ \Longrightarrow \ \ \   i\hbar\ppp{\ket{\psi(t)}}{t} = E_n\ket{\psi(t)}\nn\ee
%
} That is, if we start in an energy eigenstate then we stay in an energy eigenstate for all time. It's just the overall complex phase of the state that changes. 

\para
But what if we don't start in an energy eigenstate? Well, the time-dependent Schrodinger equation is still straightforward to solve because it's linear. All we need to do is use the fact that we met in the previous lecture that the energy eigenstates are complete -- meaning we can expand any state as a linear combination of energy eigenstates.

\para
{\co Suppose that, at $t=0$, we sit in a combination of energy eigenstates, so 
%
\be \ket{\psi(0)} = \sum_n a_n \ket{n}\nn\ee
%
fo some constant $a_n\in\C$. Then the solution to the time-dependent Schrodinger equation is
%
\be \ket{\psi(t)} = \sum_n a_n e^{-i E_nt/\hbar}\ket{n}\nn\ee
%
We can check this by differentiating,
%
\be i\hbar\ppp{\ket{\psi(t)}}{t} = \sum_n a_n e^{-iE_nt/\hbar} E_n \ket{n} = \sum_n a_n  e^{-iE_nt/\hbar}\hat{H}\ket{n} 
= \hat{H}\ket{\psi(t)}\nn\ee
%
} We see that all we have to do to figure out the time-dependence of a general state is first decompose it into energy eigenstates. Then each energy eigenstate evolves in its own sweet way, with the complex phase ticking over as $e^{-iEt/\hbar}$. And what we get is then all of those states summed up.

\para
That kind of time evolution seems very different from what we're used to in classical mechanics. It looks nothing like $F=ma$ that Newton told us to solve.

\para
{\cblue: Still needs work here, but it would be nice to get some examples. A few obvious candidates are
\begin{itemize}
\item A particle in an infinite box. We could start with some simple state -- say a wavefunction that is a straight line $\psi(x,t=0) = x$, suitably normalised. Then decompose it into Fourier modes. Then evolve each Fourier mode, and put it back together to see how $\psi(x,t)$ changes.
\item Something similar with the Gaussian. This would reproduce the result we saw for the wavepacket in an earlier lecture
\item Perhaps more challenging, we could put together a state for the harmonic oscillator. It must be possible to construct a coherent state in this way (even if we don't say the name) and then see it evolve.
\end{itemize}
%
This is probably a good opportunity for a simulation showing some of these things. And perhaps also an exercise, either using Fourier transforms or using a finite basis of states (which we'll see in a lecture or two with qubits)} \alyssa{Fourier transforms will be covered in a Tutorial for Lecture 15.} \href{https://drive.google.com/file/d/1IoUQ40qWqa0JPbx3EQL5RoLQj3zFWbqI/view?usp=drive_link}{[Lecture 15, Tutorial 1]}


\para
There's also a more formal mathematical way of saying all this. 
{\co The solution to the time-dependent Schrodinger equation can be written formally a 
%
\be \ket{\psi(t)} = \hat{U}(t)\, \ket{\psi(0)}\nn\ee
%
where $\hat{U}(t)$ is the exponential operator
%
\be \hat{U}(t) = \exp\left(-\frac{i\hat{H}t}{\hbar}
\right)\nn\ee
%
} Now this look like a bit of a cheat. It's true that if I differentiate that object $U(t)$, and treat it as if it was any other exponential, then it looks like I pull down a factor of Hamiltonian and that solves the Schrodinger equation. But am I really allowed to do this? Also, what does that exponential of an operator mean? 

\para
Well, the short answer is that -- yes -- I am allowed to do this. It can be made mathematically precise. But we can also use some of our previous ideas to write down a different expression for that solution $U(t)$. {\co  If the Hamiltonian has energy eigenstates $\hat{H}\ket{n} = E_n\ket{n}$, then we can write
%
\be \hat{U}(t) = \sum_n e^{-iE_nt/\hbar}\ket{n}\bra{n}\nn\ee
%
} Here we've met something new -- it's an object $\ket{n}\bra{n}$ with an inverted ket facing the other way. It's a cute way of writing an operator. {\co The idea is that this can act on any state $\ket{\phi}$ to give
%
\be \ket{n}\bra{n}\,\ket{\phi} = \bra{n}\phi\rangle\,\ket{n}\nn\ee
%
} Here the first term is just a number -- the overlap between the state $\ket{n}$ and the state $\ket{\phi}$. This is then multiplying the actual state $\ket{n}$. This is part of the advantage of the Dirac notation -- you get to write things like this.

\para
{\cblue Do we need more explanation here? }

\para
{\co Then, acting on an initial state $\ket{\psi(0)} = \sum_n a_n\ket{n}$, we have
%
\be \hat{U}(t) \ket{\psi(0)} = \sum_ne^{-iE_nt/\hbar} \ket{n}\bra{n}\left(\sum_m a_m \ket{m}\right) = \sum_n a_n e^{-iE_nt/\hbar}\ket{n} = \ket{\psi(t)}\nn\ee
%
} where I've used the orthonormality $\bra{n}m\rangle = \delta_{mn}$ in the middle there. {\co That's exactly what we wanted from the operator $\hat{U}(t)$.}

\para
{\co 
The operator $\hat{U}(t)$ is {\it unitary}\index{unitary}, meaning that it obeys
%
\be \hat{U}^\dagger \hat{U}=\mathds{1}\ .\ee
%
This can be seen by plugging in 
%
\be  U^\dagger U = \left(\sum_ne^{-iE_nt/\hbar} \ket{n}\bra{n}\right) \left(\sum_me^{-iE_mt/\hbar} \ket{m}\bra{m}\right) =  \sum_n \ket{n}\bra{n} = \mathds{1}\nn\ee
%
}{\cblue Definitely need more explanation here!} {\co For this reason, evolution in quantum mechanics is sometimes said to be {\it unitary}}


\para
{\cred Current quite a short lecture, but putting in an example of two would stretch it out.}

\newpage
\section{Measurement}

Finally, we get to the last part of quantum mechanics. And to what is, in many ways, the most surprising and disturbing aspect of the theory: measurement.

\para
So far, there's nothing probabalistic about what we've said. We have a state and it evolves through the Schrodinger equation. The whole set up is just as deterministic as anything in classical mechanics. If you know what state you start with, then you know what state you end up with. 

\para
But, ultimately, quantum mechanics is all about probability. And that comes when you take a look at the system and make a measurement. We don't see a particle in a superposition of different places. We see it at a specific position. When you look at the particle -- when you observe it -- that wavefunction of probability, of potentiality, reduces to something much more definite. The particle becomes localised at a point -- or at the very least in a small region -- where you see it. \href{https://capacity-simulations.github.io/simulations/Sims-v1/wavefunction_collapse.html}{Wavefunction collapse sim}. This is the esssence of measurement.

\para
We can now describe what's happening mathematically.  {\co 
Suppose that we have a quantum system in a given state $\ket{\psi}$. We want to measure an observable that is associated to an operator $\hat{A}$ whose eigenvalue equation is
%
\be \hat{A} \ket{\phi_n} = \lambda_n\ket{\phi_n}\nn\ee
%
The possible result of any measurement is taken from the spectrum of eigenvalues $\{\lambda_n\}$. But how does this result depend on the state $\ket{\psi}$?
}

\para
Here is the answer.{\co We first decompose the state $\ket{\psi}$ in terms of orthonormal eigenstates of $\hat{A}$, 
%
\be \ket{\psi} = \sum_{n} a_n\ket{\phi_n}\nn\ee
%
We will take both $\ket{\phi_n}$ and the wavefunction $\ket{\psi}$ to be normalised, which means that $\sum_n|a_n|^2=1$. Assume that the spectrum is non-degenerate. The probability that the result of the measurement gives the answer $\lambda_n$ is given by the {\it Born rule},\index{Born rule}
%
\be {\rm Prob}(\lambda_n) = |a_n|^2\nn\ee
%
A normalised wavefunction ensures that the sum of probabilities is 
%
\be \sum_n{\rm Prob}(\lambda_n)=1\nn\ee
%
Suppose that we make a measurement and get the result $\lambda_n$ for some particular $n$.  Now there's no doubt about the value of the observable: we've just measured it and know that it's $\lambda_n$. To reflect this, there is a ``collapse of the wavefunction", which jumps spontaneously to\index{collapse of wavefunction}
%
\be \ket{\psi} \ \longrightarrow \ \ket{\phi_n}\nn\ee
%
This ensures that if we make a second  measurement of  $\hat{A}$ immediately after the first, then we get the same answer, $\lambda_n$.}

\para
The discontinuous jump of the wavefunction  is very different in character from the smooth evolution  under the guidance of the time-dependent Schr\"odinger equation. In particular, the collapse of the wavefunction is not a linear operation: if you add two wavefunctions together, you first have to normalise the sum to compute the probability, and that's a simple but non-linear calculation.

\para
Given that measurement of quantum states is performed by stuff that's made out of atoms and therefore should itself be governed by the Schrodinger equation, it seems like there should be a more unified description of how states evolve. We'll talk more about this in a few lectures time when discussing the different interpretations of quantum mechanics. 

\para
{\co There's a simple generalisation of the story above when the spectrum is degenerate. In that case, the probability that we get a given result $\lambda$ is
%
\be {\rm Prob}(\lambda) = \sum_{n\, |\, \lambda_n=\lambda} |a_n|^2\nn\ee
%
After the measurement, the wavefunction  $\ket{\psi}$ collapses to $C\sum_{n|\lambda_n=\lambda} a_n\ket{\phi_n}$, where $C$ is the appropriate normalisation factor. }


\para
{\cblue Can design lots of questions here. We're going to do the qubit in a couple of lectures, so it will have to be some kind of Fourier transform-like question here. For example, we could take a state in an infinite box in a superposition of energy states and ask the probability for the energy to be something specific. Or we could be more ambitious -- take some state and ask the probability for the particle to be in the left half of the box. Then evolve in time (which requires us first do the Fourier transform of the step function) and then ask the same question some time later.} \alyssa{The practice problem for Lecture 15 is currently the more ambitious version of the superposition state in an infinite box.} \href{https://drive.google.com/file/d/15Bvpf6_OoxxVA7U5_A12x5CQM7J7rL9H/view?usp=drive_link}{[Lecture 15 Problem Statement]}

\subsection{Expectation Values}

{\co Take a normalised state $\ket{\psi}$ and measure the observable $\hat{A}$, with eigenstates and eigenvalues given by
%
\be \hat{A} \ket{\phi_n} = \lambda_n\ket{\phi_n}\nn\ee
%
As we've seen, if we expand the original state in orthonormal eigenfunctions 
%
\be \ket{\psi} = \sum_{n} a_n\ket{\phi_n}\nn\ee
%
then the probability that we measure $\lambda_n$ is $|a_n|^2$ (assuming a non-degenerate spectrum). This means that if we have many systems, all in the same state $\ket{\psi}$, and perform the same measurement on each, then the results will differ but the average will be\index{expectation value}
%
\be \langle \hat{A}\rangle_\psi =  \sum_n {\rm Prob}(\lambda_n) = \sum_n |a_n|^2 \lambda_n \nn\ee
%
}
Note the little subscript $\psi$ on the angular brackets, reminding us that the average value depends on the state of the system. We call this average the {\it expectation value}. It has a nice expression in terms of the inner product.
%
\\ {\ } \\ 
%
{\co {\bf Claim:} For a normalised state $\ket{\psi}$, the expectation value can be written as 
%
\be \langle\hat{A}\rangle_\psi = \bra{\psi}\hat{A}\ket{\psi}\nn\ee
%
We will sometimes absorb the operator $\hat{A}$ into the state and write $\hat{A}\ket{\psi}= \ket{\hat{A}\psi}$. If the state $\ket{\psi}$ is unnormalised, then the generalisation of this expression is 
%
\be \langle\hat{A}\rangle_\psi = \frac{\bra{\psi}{\hat{A}\psi}\rangle}{\bra{\psi}{\psi}\rangle}\nn\ee
%
}
In later applications we will often drop the subscript $\psi$ when it's clear what state we're talking about.


\para
We don't need to do anything clever to prove this: just follow our previous definitions. 

\para
{\co {\bf Proof:} We start with a normalised wavefuncion
%
\be \bra{\psi}\hat{A}\ket{\psi} &=& \sum_{n,m}a_m^\star a_n\bra{\phi_m}\hat{A}\ket{\phi_n}\nn\\ &=&\sum_{n,m}a_m^\star a_n \lambda_n \bra{\phi_m}{\phi_n}\rangle = \sum_n|a_n|^2\lambda_n\nn \ee
%
}
where, in the last step, we used the orthonormality of the eigenstates.



\para
{\co  We can look at some simple examples to check that this coincides with our earlier results. The expectation value of the position operator is
%
\be \langle\hat{\bx}\rangle_\psi = \int d^3x\ \bx\,|\psi(\bx)|^2\nn\ee
%
}
This is telling us that  $|\psi(\bx)|^2$ has interpretation of a probability distribution over position. This, of course, was how we first motivated the wavefunction in our earlier lectures. 

\para
{\co
The expectation value of the momentum operator, evaluated  on a wavefunction $\psi(\bx)$, is 
%
\be \langle \hat{\bp}\rangle_\psi =  -i\hbar\int d^3x\ \psi^\star(\bx)\,\nabla\psi(\bx)\ \nn\ee
%
We can think of the  wavefunction written as the Fourier transform
%
\be \psi(\bx) = \left(\frac{1}{2\pi\hbar}\right)^{3/2} \int d^3p\ \tilde{\psi}(\bp)\,e^{i\bp\cdot\bx/\hbar}\nn\ee
%
} where the factor of $(2\pi\hbar)^{-3/2}$ is because we're now working in three dimensions.  
{\co This gives
%
\be \langle\hat{\bp}\rangle_\psi &=& \frac{1}{(2\pi\hbar)^3}  \int d^3p\int d^3p'\  \bp'\, \tilde{\psi}(\bp)^\star  \tilde{\psi}(\bp') \int d^3x\ e^{-i(\bp-\bp')\cdot\bx}
\nn\\  &=& \int d^3p\int d^3p'\   \bp'\, \tilde{\psi}(\bp)^\star \tilde{\psi}(\bp')\, \delta^3(\bp-\bp')\nn\\ &=& \int d^3p\ \bp\,|\tilde{\psi}(\bp)|^2\nn \ee
%
}
where, to go from the first to second line, we've used the Fourier representation of the Dirac delta function, essentially just three copies of the 1d result. {\cblue Need an explainer about delta functions and Fourier transforms previously.} The upshot is that we learn something rather nice: while the  mod-squared of the original wavefunction $|\psi(\bx)|^2$ has the interpretation as a probability distribution over position, {\co the mod-squared of its Fourier transform $|\tilde{\psi}(\bp)|^2$ has the interpretation  as a probability distribution over momentum.}


\para
{\cred Currently quite a short lecture -- maybe we need an example or two to illustrate what's going on.}

\newpage
\section{The Qubit}

We can illustrate many of the ideas in quantum mechanics by thinking about the simple qubit. A lot of the ideas of measurement, and time evolution, and probability are simple here -- after all, we're just dealing with $2\times 2$ matrices: nothing more complicated. 

\para
{\co Recall, a qubit  is a  system with just two states -- you can think of this as $\ket{0}$ and $\ket{1}$, or as $\upket$ and $\downket$. We'll use both.

\para
The most general state is $\ket{\psi} = \alpha\ket{0} + \beta\ket{1}$. A general operator acts on the vector $(\alpha,\beta)$ as a $2\times 2$ matrix $A$. In quantum mechanics, we want $A$ to be Hermitian so $A^\dagger = A$. The most general such matrix can be written as
%
\be A = \left(\begin{array}{cc} a_0+ a_3 & a_1 - ia_2 \\ a_1+ia_2 & a_0-a_3 \end{array}\right) \nn \ee
%
The coefficient $a_0$ multiplies the identity matrix. The other three coefficients can also be thought of as multiplying fixed matrices. We write this as
%
\be A = a_0\mathds{1} + {\bf a}\cdot \bsigma\nn \ee
%
where  ${\bf a}= (a_1,a_2,a_3)$ and  $\bsigma=(\sigma^1,\sigma^2,\sigma^3)$ is a vector of matrices, defined by

%
\be \sigma^1 = \left(\begin{array}{cc} 0 \ & 1 \\ 1 \ & 0\end{array}\right)\  ,\ \ \ \sigma^2 = \left(\begin{array}{cc} 0 \ & -i \\ i \ & 0\end{array}\right)\  ,\ \ \ \sigma^3 = \left(\begin{array}{cc} 1 \, &\, 0 \\ 0 \ & -1\end{array}\right)\nn\ee
%
These are the {\it Pauli matrices}. You will also see them written as $\sigma^1=\sigma^x$ and $\sigma^2=\sigma^y$ and $\sigma^3=\sigma^z$. In quantum information, they write them simply as  as $\sigma^1=X$, $\sigma^2=Y$, and $\sigma^3=Z$.}

\para
Together with the identity matrix, the Pauli matrices form a basis of $2\times 2$ Hermitian matrices. There aren't many areas in physics where the entire world has agreed on a set of conventions but, remarkably, this is one such place, with the choice of Pauli matrices fixed.

\para
The Pauli matrices have a number of nice properties. First, each one squares to the identity matrix. Second, if you multiply two of them together then you get the third (up to a factor of $i$).  This is summarise as: {\co The Pauli matrices obey
%
\be \sigma^i \sigma^j = \delta^{ij}\mathds{1} + i\sum_{k=1}^3\epsilon^{ijk} \sigma^k\nn\ee
%
} {\cblue Need to explain what the epsilon symbol means here. I've avoided using index summation convention as I haven't explained that previously.} \alyssa{Index summation notation will be covered in a Tutorial for Lecture 16.} \href{https://drive.google.com/file/d/16p6tLVzgNswnyzVavrRlcRbfSWl1785W/view?usp=drive_link}{[Lecture 16, Tutorial 1]}


\para
{\cblue: Maybe set the following as an exercise? Or an explainer? We can also use this result to derive expressions for multiplying together more general matrices. We have
%
\be ({\bf a}\cdot \bsigma) ({\bf b}\cdot \bsigma)   &=& a_ib_j\sigma^i\sigma^j = a_ib_j(\delta^{ij}\mathds{1} + i\epsilon^{ijk}\sigma^k)\nn\\ &=& ({\bf a}\cdot{\bf b})\mathds{1} + i({\bf a}\times {\bf b})\cdot \bsigma\nn\ee
%
\alyssa{This identity will be covered in a Tutorial for Lecture 16.} \href{https://drive.google.com/file/d/1REZQEEen-HQYo7odgWdHPsDsUAqUlPBS/view?usp=drive_link}{[Lecture 16, Tutorial 2]}
}

\subsubsection*{Eigenvalues and Eigenvectors}

Each Pauli matrix has eigenvalues $\pm 1$. It's simple to check this explicitly and find the corresponding eigenvectors. 

\para
{\co The eigenvectors are  simplest for the diagonal Pauli matrix $\sigma^3$, which has
%
\be \sigma^3\left(\begin{array}{c} 1 \\ 0 \end{array}\right) = \left(\begin{array}{c} 1 \\ 0 \end{array}\right) \ \ \ {\rm and}\ \ \ \sigma^3\left(\begin{array}{c} 0\\ 1 \end{array}\right)=-\left(\begin{array}{c} 0 \\ 1 \end{array}\right)\nn\ee
%
We can also think of the corresponding operator as acting on states in the Hilbert space. We have
%
\be \hat{\sigma}^3 = \ket{0}\bra{0}-\ket{1}\bra{1} \ \ \ \Longrightarrow\ \ \  \hat{\sigma}^3\ket{0}=\ket{0}\ \ \ {\rm and}\ \ \  \hat{\sigma}^3\ket{1} = -\ket{1}\nn\ee
% 
It's not much harder to find eigenvectors for the other two Pauli matrices: we have
%
\be \sigma^1\left(\begin{array}{c} 1 \\ 1 \end{array}\right) = \left(\begin{array}{c} 1 \\ 1 \end{array}\right) \ \ \ {\rm and}\ \ \ \sigma^1\left(\begin{array}{c} 1\\ -1 \end{array}\right)=-\left(\begin{array}{c} 1 \\ -1 \end{array}\right)\ .\nn \\ 
\sigma^2\left(\begin{array}{c} 1 \\ i \end{array}\right) = \left(\begin{array}{c} 1 \\ i \end{array}\right) \ \ \ {\rm and}\ \ \ 
\sigma^2\left(\begin{array}{c} 1\\ -i \end{array}\right)=-\left(\begin{array}{c} 1 \\ -i \end{array}\right)\nn\ee
%
Neither of these eigenvectors are normalised.  Acting on states in the Hilbert space we have 
%
\be \hat{\sigma}^1 = \ket{0}\bra{1} + \ket{1}\bra{0} \ \ \ {\rm and}\ \ \ \hat{\sigma}^2=  -i\ket{0}\bra{1} + i\ket{1}\bra{0} \nn\ee
%
 So, for example,  the eigenstates of $\hat{\sigma}^1$ are
%
\be \hat{\sigma}^1\frac{1}{\sqrt{2}}(\ket{0}\pm \ket{1}) = \pm \frac{1}{\sqrt{2}}(\ket{0}\pm \ket{1})\ \nn\ee
%
where now we've included the $1/\sqrt{2}$ normalisation. }

\para
{\cblue Don't think I need to do this in the lecture. Again, could be an exercise, or an explainer. \alyssa{This will be a practice problem for Lecture 16.} \href{https://drive.google.com/file/d/1jvtDfjA08XnCLlPbO-3-0mzkFylVRBVq/view?usp=drive_link}{[Lecture 16 Problem Statement]}

\para
We could ask: what are the eigenvectors of a general matrix \mbox{${\bf n}\cdot \bsigma$}, where ${\bf n}$ is a unit vector, so $|{\bf n}|=1$. This matrix also satisfies \mbox{${\rm Tr}\,({\bf n}\cdot\bsigma)=0$} and $({\bf n}\cdot\bsigma)^2 = \mathds{1}$, so the eigenvalues are again $\pm 1$. To construct the eigenvectors, we first decompose the unit vector in terms of angles, and write
%
\be {\bf n} = (\sin\theta\,\cos\phi,\sin\theta\,\sin\phi ,\cos\theta)\nn\ee
%
with $\theta\in [0,\pi]$ and $\phi\in [0,2\pi)$.  This  obeys ${\bf n}\cdot{\bf n}=1$. We then have
%
\be {\bf n}\cdot\bsigma = \left(\begin{array}{cc} \cos\theta & \sin\theta\, e^{-i\phi}\ \\ \ \sin\theta\,e^{i\phi} & -\cos\theta\end{array}\right)\nn\ee
%
A short calculation shows that the eigenvectors are
%
\be {\bf n}\cdot\bsigma  \left(\begin{array}{c} \cos(\theta/2) \\  e^{i\phi} \sin(\theta/2)\end{array}\right)   &=&
 \left(\begin{array}{c} \cos(\theta/2) \\  e^{i\phi} \sin(\theta/2)\end{array}\right) \nn \\ 
{\bf n}\cdot\bsigma \left(\begin{array}{c}  -\sin(\theta/2) \\    e^{i\phi}\cos(\theta/2)\end{array}\right) &=& 
- \left(\begin{array}{c} -\sin(\theta/2) \\   e^{i\phi}  \cos(\theta/2)\end{array}\right)   \nn\ee
%
These follow from the trigonometric identities
%
\be \cos\theta\cos(\theta/2) +\sin\theta\sin(\theta/2) &=& \cos(\theta/2)\nn\\
\sin\theta\cos(\theta/2) -\cos\theta\sin(\theta/2) &=& \sin(\theta/2) \nn\ee
%
There's something curious in these expressions. The vector ${\bf n}$ depends on $\theta$, while the eigenvectors depend on  $\theta/2$. Let's think about how the vectors change as we send $\theta \ra \theta + 2\pi$. For example, we could set $\phi=0$ so that ${\bf n}$ points in the $(x,z)$-plane. As $\theta$ increase from 0 to $2\pi$, the vector ${\bf n}$ rotates once in this plane. However, that factor of $\theta/2$ means that the eigenvectors come back to minus themselves upon a $2\pi$ rotation. 
}






\subsection{Spin}\index{spin}


One of the most common realisation of a qubit is a spin system. For examaple, the electron, proton and neutron all have ``spin-half". That means that, in any given direction, the spin can only take one of two values: up or down.

\para
We can explain what this means using the Pauli matrices. The spin of a particle is a kind of intrinsic angular momentum. The strength of that angular momentum is measured in terms of Planck's constant $\hbar$. This has units of energy times time, which is the same as the units of angular momentum. The operator that measures the angular momentum of a spin-half particle in either the $x$, $y$ or $z$-directions is the Pauli matrix. 

\para
{\co The {\it spin operators} are defined to be
%
\be \hat{\bf S} = \frac{\hbar}{2}\hat{\bsigma}\nn\ee
%
This is a vector's worth of operators, $\hat{\bf S}=(\hat{S}^x,\hat{S}^y,\hat{S}^z)$, also written as $\hat{\bf S}=(\hat{S}^1,\hat{S}^2,\hat{S}^3)$.} The operators measure the angular momentum in the direction $x$, $y$ and $z$ respectively. All of our previous results for the Pauli matrices get trivially rescaled by that factor of $\hbar/2$. It means that if you ask for the value of the spin in any direction, there are just two possible answers: $+\hbar/2$ or $-\hbar/2$.



\para
Because we're dealing with spin states, we will revert to the $\upket$ and $\downket$ notation for the basis of states.  These are called ``spin-up" and ``spin-down" states respectively. But, more precisely, they are states that have a well-defined spin {\it in the} $z${\it -direction}. 

\para
{\co We write the eigenstates of the operator $\hat{S}^z$ as
%
\be \hat{S}^z\upket=\frac{\hbar}{2} \upket\ \ \ {\rm and}\ \ \ \hat{S}^z\downket = -\frac{\hbar}{2}\downket\nn\ee
%
} This means that the states $\upket$ and $\downket$ only have a well-defined spin in the $z$-direction. But they don't have well-defined spin in the $x$ or $y$-directions.  In fact, it's not possible to construct a state that has a well-defined spin in both $x$- and $z$-directions simultaneously. 

\para
If we want to build a state that has well-defined spin in the $x$-direction, then we need an eigenstate of $\hat{S}^x$. {\co We can construct eigenstates of $\hat{S}^x$ as
%
\be \rightket = \frac{1}{\sqrt{2}} (\upket + \downket) \ \ \ {\rm and}\ \ \ \leftket = \frac{1}{\sqrt{2}}(\upket-\downket)\nn\ee
%
These states obey
%
\be \hat{S}^x\rightket=\frac{\hbar}{2} \rightket\ \ \ {\rm and}\ \ \ \hat{S}^x\leftket = -\frac{\hbar}{2}\leftket\nn\ee
%
But we can see what that means. If we have the state $\rightarrow$ then it's got well-defined angular momentum in the $x$-direction. But it doesn't have well-defined angular momentum in the $z$-direction -- it's a superposition of up and down states in the $z$-direction. 

\para
{\cblue Exercise: Find eigenstates in the $y$-direction? \alyssa{Finding the eigenstates in a general direction is included in the practice problem for Lecture 16.} \href{https://drive.google.com/file/d/1jvtDfjA08XnCLlPbO-3-0mzkFylVRBVq/view?usp=drive_link}{[Lecture 16 Problem Statement]} Or possibly find eigenstates in a general direction ${\bf n}$ is ${\bf n}\cdot \hat{\bf S}$ which are given by
%
\be \ket{{\bf n};+} &=& \cos(\theta/2) \upket + e^{i\phi}\sin(\theta/2)\downket\nn\\ \ket{{\bf n};-}&=& -\sin(\theta/2)\upket + e^{i\phi} \cos(\theta/2)\downket  \nn\ee
%
These have eigenvalues $\pm \hbar/2$, respectively. }


\para
We can use this to understand just one strange measuremt is in quantum mechanics. 


\para
{\co Start with a spin in a general normalised state
%
\be \ket{\psi} = \alpha \upket + \beta \downket\nn\ee
%
with $|\alpha|^2 + |\beta|^2 = 1$.  To start, we make a  measurement of the spin in the $z$-direction.} This is straightforward because our state is already written in terms of $\hat{S}^z$ eigenstates, so we know that we measure spin $s^z=\pm \hbar/2$ with probabilities
%
\be s^z = \left\{\begin{array}{cc} +{\hbar}/{2}\ \ \ &\mbox{with probability $|\alpha|^2$}  \\ -{\hbar}/{2}\ \ \ &\mbox{with probability $|\beta|^2$}\end{array}\right..\ee
%
Suppose that we found the result $s^z=+\hbar/2$. Then the wavefunction $\ket{\psi}$ collapses to the state $\upket$. If we again measure $\hat{S}^z$, then there's no doubt about what we will find: it will always give us the result $s^z=+\hbar/2$. 

\para
Next we choose to measure the spin in the $x$-direction. To understand what happens, we need to decompose our state $\upket$ in eigenstates of $\hat{S}^x$. These are the $\rightket$ and $\leftket$ states.  Inverting, we have
%
\be \upket = \frac{1}{\sqrt{2}}(\rightket + \leftket)\nn\ee
%
The Born rule now tells us that if we measure $\hat{S}^x$, there's no way to know what we will find: both possibilities are equally likely,
%
%
\be s^x = \left\{\begin{array}{cc} +{\hbar}/{2}\ \ \ &\mbox{with probability 1/2}  \\ -{\hbar}/{2}\ \ \ &\mbox{with probability 1/2}\end{array}\right..\nn\ee
%
Suppose that we again find the result $+\hbar/2$. This time, it's $s^x=+\hbar/2$, so the wavefunction collapses to $\rightket$. Now we go back and remeasure our original spin $\hat{S}^z$. It wasn't so long ago that we confirmed (repeatedly!) that $s^z=+\hbar/2$. But now we've made a measurement of $\hat{S}^x$ and this, of course, changes the story. We must decompose our state $\rightket$  in terms of $\upket$ and $\downket$ $s^z$ will yield
%
%
\be s^z = \left\{\begin{array}{cc} +{\hbar}/{2}\ \ \ &\mbox{with probability 1/2}  \\ -{\hbar}/{2}\ \ \ &\mbox{with probability 1/2}\end{array}\right..\ee
%
We see that all memory of our previous $\hat{S}^z$ measurement has been erased by the measurement of $\hat{S}^x$. There's no longer any guarantee that we will still get $s^z=+\hbar/2$ this time around. Instead, the values $\pm \hbar/2$ are equally likely. }


\para
This is the essence of quantum mechanics. Here, the calculations are simple, but the concepts are not! 

\para
There is no meaning to the spin of the particle in both $x$- and $z$-directions simultaneously. If you know one precisely, then you have no knowledge of the other. 


\para
Part of the issue here is that the measurement of a quantum system is never innocent. You can't just take a small peek and walk away pretending that you've not done anything. Instead, the collapse of the wavefunction means that any measurement of  a quantum system necessarily disturbs the state. 

\para
If the state originally had a specific value for one observable,  and you then measure a different observable,  then you destroy the original property of the state. It's tempting to say that there's no way to {\it know} the values of those two observables -- say spin in the $x$ and $z$ direction -- at the same time. But that's not what the mathematics is telling us. Instead, the quantum particle can't  {\it have} specific values of spins in both directions at the same time. Indeed, most  of the time it doesn't have a specific value for either since its wavefunction is a superposition of eigenstates. This is the essence of the Heisenberg uncertainty principle, which we will meet in the next lecture -- it's a way to quantify how much information we can have about two different observables. 



\subsection{The Stern--Gerlach Experiment}

{\cblue I wonder if this might make for a good explainer or a tangent? (Another suggestion was a tangent on the Dirac equation and the origin of spin. That could work as well.) \href{https://capacity-simulations.github.io/simulations/Sims-v2/stern-gerlach.html}{[Stern-Gerlach sim]}

\para
The {\it Stern--Gerlach experiment} is an iconic demonstration of some of the results described above. It was designed by Otto Stern, and carried out by Walther Gerlach in Frankfurt in 1922. Ironically, this was some years before the idea of spin was proposed and the striking results of the experiment were initially misinterpreted (as something that they then called ``space quantisation" but is better thought of as a quantisation of orbital angular momentum of electrons in an atom).

\para
The point of the experiment is to find a way to measure the spin. In the original experiment,  this was done not for electrons but for silver atoms. This was a  rather fortuitous choice, as it is now known that silver atoms do indeed have spin 1/2. The key idea is to translate the quantised spin of the atom into a quantised change of position.  This is achieved by putting the atom in a magnetic field ${\bf B}$.



\para
In the book, I describe how the force comes about...not entirely sure it's needed. It's really the beam splitting that's important.




\para
The Stern--Gerlach experiment (figure needed)  takes advantage of this force. A beam of spin 1/2 particles moves in the $x$-direction, and is subjected to a magnetic field that, as above, is largely pointing in the $z$-direction, and whose strength varies.  The particles then experience a force which, depending on the orientation of their spin, pushes them either up or down in the $z$-direction. Further down the line, a screen is placed that measures the position of each particle. 





\para
The  position acts as a proxy for the spin. If the particle was in the state $\upket$, the particle will have moved upwards; if the particle was in the state $\downket$, the particle will have moved downwards. But what if it's in the general state $\ket{\psi} = \alpha\upket + \beta\downket$.  One might think that the particle gets pushed a little bit up, or a little bit down, depending on the relative values of $\alpha$ and $\beta$. That's what a classical force would suggest. But that's not what's seen.  Measuring the position of the particle on the screen is tantamount to measuring the spin $\hat{S}^z$ and collapsing the wavefunction. As we've seen, the value of the observed spin can {\it only} be $s^z=+\hbar/2$ or $s^z=-\hbar/2$. This is true even in the state general state,  where $\alpha$ and $\beta$ tell us the probabilities of observing these two outcomes. 

\para
This means that the beam isn't deflected in a continuous way as $\alpha$ and  $\beta$ change. Each particle either goes up or it goes down. The beam will be split into two, with $|\alpha|^2$ and $|\beta|^2$ determining the probability of any given particle to go up or down (or, equivalently, the fraction of the beam that goes up, and the fraction that goes down).

\para
That is the striking result that Stern and Gerlach observed in their experiment. When the beam hits the screen, it forms two  distinct dots with nothing in the middle. In the quantum world, you don't get to sit in the superposition forever: when you're observed you have to make a choice. Either you're up. Or you're down. There is no in between. 

}





\subsection{Time Evolution}

So far, we've not considered the time evolution of the qubit itself. All of the discussion of measurement took place assuming the Hamiltonian vanishes, so there's no other way that the spin evolves between measurements. But we can also use the qubit to illustrate how the spin evolves due to a Hamiltonian.

\para
Physically, there's a very natural way to induce a Hamiltonian. We just put the spin in a magnetic field 

\para
{\co In a magnetic field ${\bf B}$, the Hamiltonian for the spin is given by 
%
\be \hat{H}= -\mu{\bf B}\cdot \hat{\bf S}\nn\ee
%
with $\mu$ a constant.} (There's a story about the value of this constant and how we can predict it from first principles. But we'll leave this for another day.) In fact, this is the most general form of the Hamiltonian for a qubit. For other physical realisations of a qubit, the meaning of ${\bf B}$ and $\mu$ will be different, but the Hamiltonian will still take this form.




\para
{\co Suppose we take ${\bf B} = (0,0,B)$. The energy eigenstates are then $\upket$ and $\downket$, with
%
\be \hat{H}\upket =-\frac{\hbar \mu B}{2} \upket\ \ \ {\rm and}\ \ \ \hat{H}\downket = +\frac{\hbar\mu B}{2} \downket\nn\ee
%
The time-dependent Schrodinger equation  is 
%
%
\be i\hbar\ppp{\ket{\psi}}{t} = H\ket{\psi} \nn\ee
%
This tells us that these eigenstates evolve in time as $e^{+i\mu B t/2}\upket$ and $e^{-i\mu Bt/2}\downket$. The general state $\ket{\psi}= \alpha \upket + \beta\downket$ then evolves as
%
\be \ket{\psi(t)} = \alpha e^{i\mu B t/2} \upket + \beta e^{-i\mu B t/2}\downket \nn\ee
%
If we measure $\hat{S}^z$, then  the probabilities that we get $\pm \hbar/2$ are still $|\alpha|^2$ and $|\beta|^2$, respectively, and do not change with time. But there's a different story if we measure, say, $\hat{S}^x$. Now we decompose the state as
%
\be \ket{\psi(t)} &=& \frac{\alpha}{\sqrt{2}} e^{i\mu Bt/2} \left(\rightket + \leftket\right) + \frac{\beta}{\sqrt{2}}e^{-i\mu B t/2}\left(\rightket-\leftket\right)\ .\ \ \ \ \ \ee
%


This means that the probability that $s^x=+\hbar/2$ is 
%
\be P_+ = \frac{1}{2}\left|\alpha e^{i\mu B t/2} + \beta e^{-i\mu B t/2}\right|^2\nn\ee
%
and this is very much time dependent.} {\cblue plot this.} {\co As a particular example, suppose that we pick $\alpha = \beta = 1/\sqrt{2}$. This means that we prepare our initial  state as $\ket{\psi(t=0)}=\rightket$. Over time we have
%
\be \ket{\psi(t)} = \cos \left(\frac{\mu Bt}{2}\right)\rightket + \sin\left(\frac{\mu Bt}{2}\right)\leftket\ .\label{zenoround}\ee
%
We see that the state oscillates between the two orthonormal states $\rightket$ and $\leftket$.} 






\newpage
\section{Commutation Relations}

The key idea of measurement -- and one of the key ideas of quantum mechanics -- is that there are certain properties -- for example, the spin of a particle in different directions -- that are mutually incompatible. It's not possible for a quantum state to have well-defined values for both at the same time. If you know one for sure, then the other will be uncertain - there will be some probability associated to it. 


\para
There's a way to formalise this mathematically, to determine when a measurement of one observable will necessarily mess up the other. he underlying structure is known as {\it commutation relations}.

\para
{\co 
Given two operators $\hat{A}$ and $\hat{B}$, their {\it commutator} is defined to be
%
\be [\hat{A},\hat{B}] = \hat{A}\hat{B} - \hat{B}\hat{A}\ \nn\ee
%
}
The key idea is familiar from matrices. If you multiply two matrices $A$ and $B$, then the order matters. Typically $AB$ is not the same as $BA$ and the commutator $[A,B]$ captures the difference between these.  {\co Clearly    $[\hat{A},\hat{B}]=- [\hat{B},\hat{A}]$.}

\para
When thinking about formulae like the commutator, you should remember that  the operators have a job to do: they act on states. It's important to remember this when manipulating these kind of equations.

\para
{\co A good rule of thumb is to think of operator equations  like this as meaning
%
\be [\hat{A},\hat{B}]\ket{\psi} = \hat{A}\hat{B}\ket{\psi} - \hat{B}\hat{A} \ket{\psi}\ \ \ \mbox{for all $\ket{\psi}$}\nn\ee
%
When we do calculations, at least initially,  this phrasing helps in  avoiding slip ups.}

\para
For particles roaming around,  the fundamental operators are the  position $\hat{\bx}$ and the momentum  $\hat{\bp}$. It will prove useful to compute their commutators. 
%
\\ {\ } \\ 
%
{\co {\bf Claim:} The commutation relations between position and momentum operators are
%
\be [\hat{x},\hat{x}] = [\hat{p},\hat{p}] = 0 \ \ \ {\rm and}\ \ \ \ [\hat{x},\hat{p}]= i\hbar\nn\ee
%
These are known as the {\it canonical commutation relations} on account of the important role that they play in quantum mechanics.} Note that the right-hand side of the final relation is just a constant, but this too should be viewed as an operator: it comes with an implicit unit operator $\mathds{1}$ and just multiplies any state by a constant.

\para
{\co In 1d, the commutation relation is $[\hat{x},\hat{p}]=i\hbar$. This kind of relation is not possible for finite matrices $A$ and $B$; the analogous result would be  $[A,B]=i\hbar \mathds{1}$, with $\mathds{1}$ the unit matrix, but if we take the trace of the left-hand side we find
% 
\be {\rm Tr}[A,B] = {\rm Tr}(AB-BA) = 0\nn\ee
%
because ${\rm Tr}(AB) = {\rm Tr}(BA)$. Yet clearly ${\rm Tr}\,\mathds{1} \neq 0$.} However, infinite matrices, also known as operators, can do things that finite matrices cannot.
%
\\ {\ } \\ 
%
{\co {\bf Proof:} This follows by just being careful with definitions. We have
%
\be \hat{x} = x\ \ \ {\rm and}\ \ \  \hat{p} = -i\hbar \frac{d}{dx}\nn\ee
%
We should view all operators as acting on wavefunctions $\psi(x)$. The first commutators $[\hat{x},\hat{x}] = [\hat{p},\hat{p}]=0$ are trivial: anything commutes with itself. The interesting one is
%
\be  [\hat{x},\hat{p}]\psi(x) = -i\hbar\left(x\ppp{}{x} - \ppp{}{x}{x}\right)\psi(x) \nn\ee
%
}
We can see the issue: the first derivative hits only the function $\psi$. Meanwhile, the second derivative acts on $x\psi$, giving us an extra term from the chain rule. Moreover, the $\partial \psi/\partial x$ terms cancel, leaving us only with this extra term. {\co We have
%
\be   [\hat{x},\hat{p}]\psi(\bx) = i\hbar \left(x\frac{d\psi}{dx} - \frac{d(x\psi)}{dx}\right)  = i\hbar \psi)(x) \ \ \ \mbox{for all $\psi(x)$}\nn\ee
%
This is the claimed result. \hfill $\Box$

\para
The generalisation of the canonical commutation relations to 3d, 
 $\hat{\bx} = (\hat{x}_1,\hat{x}_2,\hat{x}_3)$ and $\hat{\bp} = (\hat{p}_1,\hat{p}_2,\hat{p}_3)$, is
%
\be [\hat{x}_i,\hat{x}_j] = [\hat{p}_i,\hat{p}_j] = 0\ \ \ {\rm and}\ \ \ [\hat{x}_i,\hat{p}_j] = i\hbar \delta_{ij}\nn\ee
%
}{\cblue Possibly set this as an exercise? The novelty is in noting that partial derivatives commute.} \alyssa{Showing the generalisation of the canonical commutation relations to 3d is set as the first part of the practice problem for Lecture 17. \href{https://drive.google.com/file/d/1RcCfNfrKIEmnRbL_AyJe6mOLPKukT4Fm/view?usp=drive_link}{[Lecture 17 Practice Problem]}}



\para
The logic that we've presented above is to first introduce the position and momentum operators and subsequently derive the canonical commutation relations. This, it turns out, is inverted. In more advanced presentations of quantum mechanics, the canonical commutation relations are the starting point. There is then a theorem -- due to Stone and von Neumann -- that, roughly speaking, says that one can derive the form of the operators  from the canonical commutation relations.



\para
{\co 
Occasionally we will meet two operators that commute, meaning
%
\be [\hat{A},\hat{B}]=0\nn\ee%
%
This is interesting because \ldots
%
\\ {\ }\\
%
{\bf Claim:} $[\hat{A},\hat{B}]=0$ if and only if $\hat{A}$ and $\hat{B}$ share the same eigenstates.
%
\\ {\ }\\
%
{\bf Proof:}  First suppose that both $\hat{A}$ and $\hat{B}$ share the same eigenstates, so that
%
\be \hat{A}\ket{\phi_n} = \lambda_n\ket{\phi_n}\ \ \ {\rm and}\ \ \ \hat{B}\ket{\phi_n} = \mu_n\ket{\phi_n}\nn\ee
%
Then any state $\ket{\psi}$ can be expanded in terms of these eigenstates $\ket{\psi}= \sum_na_n\ket{\phi_n}$, giving
%
\be[\hat{A},\hat{B}] \ket{\psi} = \sum_n a_n [\hat{A},\hat{B}]\ket{\phi_n} = \sum_n a_n (\lambda_n\mu_n-\mu_n\lambda_n)\ket{\phi_n}=0\nn\ee
%
Because this equation holds for all $\ket{\psi}$, it is equivalent to saying $[A,B]=0$. Conversely, suppose that $[\hat{A},\hat{B}]=0$. We restrict  ourselves to the case where $\hat{A}$ has a non-degenerate spectrum.  If $\ket{\phi}$ is an eigenstate of $\hat{A}$, then we have
%
\be \hat{A}\ket{\phi} = \lambda \ket{\phi} \ \ \ &\Longrightarrow&\ \ \ \hat{B}\hat{A}\ket{\phi} = \lambda\hat{B}\ket{\phi} \nn\\ 
&\Longrightarrow&\ \ \ \hat{A}\hat{B}\ket{\phi} = \lambda\hat{B}\ket{\phi}\nn\ee
%
where the second line follows because $[\hat{A},\hat{B}]=0$. But this tells us that $\hat{B}\ket{\phi}$ is also an eigenstate of $\hat{A}$ with eigenvalue $\lambda$. By assumption, $\hat{A}$ has only one eigenstate with eigenvalue $\lambda$, which means that we must have
%
\be \hat{B}\ket{\phi} = \mu\ket{\phi}\nn\ee
%
for some $\mu$. But this is the statement that $\ket{\phi}$ is also an eigenstate of $\hat{B}$. The two operators need not have the same eigenvalues, but they do have the same eigenstates.
\hfill $\Box$}


\para
There's a straightforward generalisation of this proof that shows the result still holds even when $A$ and $B$ have degenerate spectra.



\para
There's an important physical consequence of this result. {\co It {\it is} possible for  a quantum system to have simultaneous values for commuting observables. If you measure $\hat{A}$ and then measure $\hat{B}$, the state will not be perturbed as long as $[\hat{A},\hat{B}]=0$.} (Actually this last statement does require that the spectrum of $\hat{A}$ is non-degenerate.) If you then measure $\hat{A}$ again, you'll get the same answer as the first time.  







\subsection{The Heisenberg Uncertainty Principle}

Quantum mechanics is all about probability. If you have some probability distribution  over a bunch of different possibilities, then there are a number of questions you can ask. The simplest is the average -- that's what we call the expectation value in quantum mechanics.  The next simples is the 
 the variance, which gives some indication of the spread of the distribution.

\para
{\co In quantum mechanics the variance  is called {\it uncertainty}. In a normalised state $\ket{\psi}$, the uncertainty of an observable $\hat{A}$ is denoted $\Delta_\psi {A}$ and is defined as
%
\be (\Delta_\psi{A})^2 &=& \langle (\hat{A}-\langle\hat{A}\rangle_\psi)^2\rangle_\psi
\nn\\ &=& \langle \hat{A}^2- 2\langle\hat{A}\rangle_\psi\hat{A} + \langle\hat{A}\rangle_\psi^2 \rangle_\psi  
\nn\\ &=&\langle\hat{A}^2\rangle_\psi - \langle \hat{A}\rangle_\psi^2\nn\ee
%
}
In statistics, the uncertainty $\Delta_\psi A$ is also  known as the {\it standard deviation} of the distribution. 

\para
{\co We know that the expectation value is $\langle \hat{A}\rangle_\psi = \bra{\psi}\hat{A}\ket{\psi}$. Using this, and assuming a normalised wavefunction, we can write the uncertainty as
%
\be (\Delta_\psi{A})^2 = \bra{\psi}\hat{A}^2\ket{\psi} - \bra{\psi}\hat{A}\ket{\psi}^2\nn\ee
%
To motivate the name ``uncertainty", suppose that the state $\ket{\phi}$ of the system is an eigenstate of $\hat{A}$, so  $\hat{A}\ket{\phi} = \lambda\ket{\phi}$. 
Then the uncertainty of $\hat{A}$ in this state is
%
\be (\Delta_\phi{A})^2 = \bra{\phi}\hat{A}^2\ket{\phi} - \bra{\phi}\hat{A}\ket{\phi}^2 = \lambda^2 -\lambda^2 = 0\nn\ee
%
}
This makes sense: in an eigenstate $\ket{\phi}$ we know exactly what the value of $\hat{A}$ is, so there is no uncertainty in the measurement of $\hat{A}$. 

\para
But the full story is more interesting. It turns out  that, for certain pairs of observables, as you reduce the uncertainty in one, the uncertainty in the other necessarily grows. One such pair is our favourite position and momentum. 

\para
{\co We work in one dimension. The canonical commutation relations 
%
\be [\hat{x},\hat{p}] = i\hbar\nn\ee
%
From this, we derive the following result:
%
\\ {\ }\\
{\bf Claim:} In any state $\ket{\psi}$, we necessarily have
%
\be \Delta_\psi x\,\Delta_\psi p \geq \frac{\hbar}{2} \ \nn
\ee
%
This is the mathematical expression of the {\it Heisenberg uncertainty relation}. The better you know the position of the particle, the less well you know its momentum.
%
\\ {\ }\\
%
{\bf Proof:} To make life simple, we will assume that we have a  normalised state $\ket{\psi}$ for which $\langle{x}\rangle=\langle{p}\rangle=0$. (If this isn't the case, you can rescue the proof below by defining new operators  $\hat{X} = \hat{x} -\langle{x}\rangle$ and $\hat{P} = \hat{p} -\langle{p}\rangle$ which do obey $\langle{X}\rangle = \langle{P}\rangle=0$.)

\para
Next, consider the one-parameter family of states defined by
%
\be \ket{\Psi_s}= (\hat{p} - is\hat{x})\,\ket{\psi}\ \ \ {\rm with}\  s\in\mathds{R}\nn\ee
%
The norm of this state is, like all others, positive definite, meaning 
%
\be \bra{\Psi_s}{\Psi_s}\rangle \geq 0\nn\ee
%
Translated back to the original state $\ket{\psi}$, this tells us that
%
\be 0 \leq \bra{(\hat{p}-is\hat{x})\psi}{(\hat{p}-is\hat{x})\psi}\rangle  = \bra{\psi}{(\hat{p}+is\hat{x})(\hat{p}-is\hat{x})\psi}\rangle\nn\ee
%
where the equality uses the fact that both $\hat{x}$ and $\hat{p}$ are Hermitian. Expanding out, we have
%
\be 
 0\leq \bra{\psi}{(\hat{p}^2+is[\hat{x},\hat{p}] + s^2\hat{x}^2)\psi}\rangle =\bra{\psi}{(\hat{p}^2 -s\hbar + s^2\hat{x}^2)\psi}\rangle\nn  \ee
 %
 where, this time, the equality uses the commutation relation $[\hat{x},\hat{p}]=i\hbar$.  The $\hat{p}^2$ and $\hat{x}^2$ terms in this expression are just the uncertainties because we've chosen $\langle{x}\rangle=\langle{p}\rangle=0$. So we have
 %
 \be (\Delta_\psi p)^2 -s\hbar + s^2(\Delta_\psi x)^2 \geq 0 \ \ \ \mbox{for all $s\in\mathds{R}$}\nn\ee
 %
 This is a quadratic in $s$ and the inequality tells us that it has either one  root or no roots. This can only be true if the discriminant ``$b^2-4ac$" is non-positive -- that is,
 %
 \be 4(\Delta_\psi x)^2(\Delta_\psi p)^2\geq \hbar^2\nn\ee
 %
 This is the Heisenberg uncertainty relation, \hfill $\Box$
 }
 
 \para
 {\cblue Could do this in the lecture, or as an aside:}
 You can have some fun trying to come up with experiments that might evade the Heisenberg uncertainty relation, and then see how nature finds a way to  ensure that the result is satisfied after all. 
 
 \para
 For example, the obvious way to figure out where a particle is sitting is to look at it. But if you want to resolve something on a distance scale $\Delta x$, then you have to use light of wavelength $\lambda \lesssim \Delta x$.
 Ultimately, light is made up of particles called photons. We haven't said anything about photons so far.  However, all we need to know is that, like all other quantum particles, photons have a momentum given by the de Broglie formula 
 %
 \be p = \frac{2\pi\hbar}{\lambda}\nn\ee
 %
 Clearly there's a pay-off: if you want to be sure that the particle is sitting in a smaller and smaller region $\Delta x$ then you have to hit it with photons that carry higher and higher momentum, and these will impart some recoil on the particle of order  $\Delta p \sim 2\pi\hbar/ \Delta x$. The intrusive nature of measurement means that we can't  know both the position and momentum to better than  $\Delta p \Delta x \sim 2\pi \hbar$. The Heisenberg uncertainty relation provides the more accurate bound.

 


\para
\alyssa{The proof of the generalised Heisenberg uncertainty relation is part of the practice problem for Lecture 17.} \href{https://drive.google.com/file/d/1RcCfNfrKIEmnRbL_AyJe6mOLPKukT4Fm/view?usp=drive_link}{[Lecture 17 Practice Problem]}
{\cblue Perhaps set this as a problem: It it straightforward to repeat the proof above for general observables $\hat{A}$ and $\hat{B}$. One then finds that, in any state $\ket{\psi}$, 
%
\be (\Delta_\psi A)(\Delta_\psi B) \geq \frac{1}{2}|\,\langle [\hat{A},\hat{B}]\rangle\,|\ .\nn\ee
%
This is the generalised Heisenberg uncertainty relation. }
 
\subsubsection*{The Gaussian Wavepacket Revisited}

{\cblue Do this in the lecture, or set as an extended exercise?} \alyssa{This section will be covered in a Tutorial for Lecture 17.} \href{https://drive.google.com/file/d/1So8tapOEUkljUpnU8ERyn7qTcuPAARaS/view?usp=drive_link}{[Lecture 17, Tutorial 1]}

\para
We already met a baby version of the Heisenberg uncertainty relation when we studied the Gaussian wavepacket. We can return to this example now that we have a better idea of what we're doing.

\para
We'll forget about time dependence  and look at the normalised Gaussian wavefunction in one dimension
%
\be \psi(x) = \left(\frac{a}{\pi}\right)^{1/4} e^{-ax^2/2}\nn\ee
%
It's not hard to check that $\langle{x}\rangle = \langle{p}\rangle=0$. What about the uncertainty?

\para
The uncertainty in position is given by 
%
\be (\Delta_\psi x)^2 = \langle \hat{x}^2\rangle_\psi = \sqrt{\frac{a}{\pi}}\int_{-\infty}^{+\infty} dx\ x^2 e^{-ax^2}\nn\ee
%
This integral is straightforward. We start from the Gaussian integral
%
\be \int_{-\infty}^{+\infty} dx\ e^{-ax^2} = \sqrt{\frac{\pi}{a}}\nn\ee
%
Then we differentiate both sides with respect to $a$. This gives
%
\be  \ppp{}{a}\int_{-\infty}^{+\infty} dx\ e^{-ax^2} &=& - \int_{-\infty}^{+\infty} dx\ x^2e^{-ax^2} = \ppp{}{a} \sqrt{\frac{\pi}{a}} = -\frac{1}{2}\sqrt{\frac{\pi}{a^3}}\nn\ee
%
So we have the result
%
\be \int_{-\infty}^{+\infty} dx\ x^2e^{-ax^2} =  \frac{1}{2}\sqrt{\frac{\pi}{a^3}}\nn\ee
%
Substituting, we get the uncertainty in position to be
%
\be (\Delta_\psi x)^2 = \frac{1}{2a}\nn\ee
%
This is to be expected: it is the usual variance of a Gaussian distribution.


\para
For the uncertainty in momentum, we have
%
\be (\Delta_\psi p)^2 &=&\langle \hat{p}^2\rangle_\psi =  \sqrt{\frac{a}{\pi}}\int_{-\infty}^{+\infty} dx\ e^{-ax^2/2}\left[-\hbar^2 \frac{d^2}{dx^2}e^{-ax^2/2}\right]\nn\\ &=& 
\sqrt{\frac{a}{\pi}} \hbar^2 \int_{-\infty}^{+\infty} dx\ (a-a^2x^2) e^{-ax^2} \nn\\ 
&=& \sqrt{\frac{a}{\pi}} \hbar^2 \left[\sqrt{\frac{\pi}{a}}a-\frac{a^2}{2}\sqrt{\frac{\pi}{a^3}}\right] = \frac{1}{2}\hbar^2 a
\nn\ee
%
Multiplying these results together gives
%
\be \Delta_\psi x\,\Delta_\psi p = \frac{\hbar}{2}\nn\ee
%
We see that the Gaussian wavepacket is rather special: it saturates the bound from the Heisenberg uncertainty relation. The class of Gaussian wavefunctions, parameterised by $a$, does the best job possible of balancing the twin requirements of localising in both position and momentum.







\subsection{Ehrenfest Theorem}


Commutation relations are good for many things in quantum mechanics. Here's another: they give us a nice expression for how the the expectation value of some operator $\hat{A}$ changes with time. 
We'll assume that the operator itself has no time dependence, so all change happens because the state is evolving. 

\para
{\co We have
%
\be \ppp{}{t}\langle \hat{A}\rangle_\psi = -\frac{i}{\hbar}\langle[\hat{A},\hat{H}]\rangle_\psi\ .\nn\ee
%
This is known as the {\it Ehrenfest theorem}.} The proof follows from putting together the various pieces that we've learned so far. The expectation value $\langle\hat{A}\rangle_\psi = \bra{\psi}\hat{A}\ket{\psi}$ involves two copies of the state, one a bra and one a ket. The chain rule means that we differentiate both in turn. 

\para
{\co To prove this, we use the Schrodinger equation, together with its adjoint version
%
\be i\hbar\ppp{\ket{\psi(t)}}{t} = \hat{H}\ket{\psi(t)}\ \ \ \Longrightarrow\ \ \ -i\hbar\ppp{\bra{\psi}}{t} =\bra{\psi}\hat{H}\nn\ee
%
We then have
%
\be \ppp{}{t}\langle \hat{A}\rangle_\psi&=& \left(\ppp{}{t}\bra{\psi}\right)\hat{A}\ket{\psi} + \bra{\psi}\hat{A}\ppp{}{t}\ket{\psi} \nn\\
&=&\frac{i}{\hbar}\bra{\psi}\hat{H}\hat{A}\ket{\psi} - \frac{i}{\hbar}\bra{\psi}\hat{A}\hat{H}\ket{\psi}=-\frac{i}{\hbar}\bra{\psi}[\hat{A},\hat{H}]\ket{\psi}\nn\ee
%
as promised.}

\para
Note, also, that any operator $\hat{A}$ that commutes with the Hamiltonian, so that $[\hat{A},\hat{H}]=0$, will have a constant expectation value $\langle \hat{A}\rangle_\psi$ regardless of the state $\ket{\psi}$. This sounds very much like a quantity that is conserved in classical mechanics, and we will make this analogy precise in Chapter 


\para
{\co 
Ehrenfest's theorem gives a bridge between the classical and quantum worlds. In particular, suppose that we take the Hamiltonian for a particle moving in a potential in 3d,
%
\be \hat{H} = \frac{1}{2m}\hat{\bp}^2 + V(\hat{\bx})\nn\ee
%
Then, in any state, we have
%
\be \frac{d\langle \hat{\bx}\rangle}{dt} &=& -\frac{i}{\hbar} \langle [\hat{\bx},\hat{H}]\rangle = \frac{1}{m}\langle \hat{\bp}\rangle\nn\\ \frac{d\langle \hat{\bp}\rangle}{dt} &=& -\frac{i}{\hbar} \langle [\hat{\bp},\hat{H}]\rangle =  - \langle \nabla V (\hat{\bx})\rangle
\nn\ee
%
These are Newton's equations of motion, here rendered as Hamilton's equations. }



\para
{\cblue There's a lot of possible exercises that we can do here. One is to revisit the creation and annihilation operators of the harmonic oscillator, $\hat{a}$ and $\hat{a}^\dagger$. We introduced these previously but without really highlighiting what they were. Now we can redo the whole harmonic oscillator story but using $[\hat{a},\hat{a}^\dagger]=1$. } \alyssa{We will revisit the Quantum Harmonic Oscillator with creation and annhilation operators in a Tutorial for Lecture 17} \href{https://drive.google.com/file/d/14k03OegKK7EO5YkkT6GodccZlMBsELQO/view?usp=drive_link}{[Lecture 17, Tutorial 2]}

\para
{\cred This feels like a long lecture to me -- maybe 45 minutes, maybe longer. The qubit one similar. I have a nagging feeling that things have slowly been increasing in length.}

\newpage
\section{Interpretations of Quantum Mechanics}

{\cblue This lecture needs thinking about! :)}


\newpage
\noindent{\Large \bf Section 4: Particles in Three Dimensions}

\para
{\cred I think that the material in this section takes about 5 lectures in Cambridge. Lectures are, weirdly, 50 minutes each. So although I've put it into just four lectures, it may be that we have to rethink the timing.} 


\para
Now we have a better idea of the basic framework of quantum mechanics, we can return to using it to solve some things. And our goal in this section is to use quantum mechanics to understand one of the basic problems, the one that first drove Heisenberg to formulate the theory: the structure of atoms. Or, since we're going to take things easier, the structure of the hydrogen atom in particular.

\para
{\co We're going to solve a class of more general problems: a particle, moving in three dimensions with a potential. This means that we will find eigenstates of the Hamiltonian
%
\be \hat{H} = -\frac{\hbar^2}{2m}\nabla^2 + V(\bx)\ee
%
where
%
\be \nabla^2 = \ppp{}{x^2} + \ppp{}{y^2} + \ppp{}{z^2}\ .\ee
%
} This is still pretty limiting. We're still only talking about one particle, not many. That means that we won't be able to talk about atoms that are more complicated than hydrogen because they contain more electrons. 

\para
{\co For simplicity, we will focus only on  {\it central potentials}, meaning that
potential
%
\be V(\bx) = V(r)\ \ \ {\rm with}\ r=|\bx|\ .\ee
%
These are both the simplest potentials to study and, happily, the kinds of potentials that are most useful for physics. }

\para
The Schrodinger equation 
%
\be\hat{H}\psi = E\psi\ee
%
is a partial differential equation. These are usually very hard to solve. But, as will see, it's possible to use the rotational symmetry of the central potential to our advantage so that we can reduce the Schrodinger equation to somehing that is very much solvable.





\section{Angular Momentum}

The trick to solving these kinds of problems is to use the rotational symmetry to our advantage. And the way this manifests itself is that the angular momentum of the particle is conserved. The strategy to solving the problem is to first figure out what the angular momentum is doing, and then deal with everything else.

\para
{\co In classical mechanics, the angular momentum is
%
\be {\bf L} = \bx\times \bp\nn\ee
%
The quantum angular momentum operator is
%
\be \hat{\bf L} = \hat{\bx} \times \hat{\bp} = -i\hbar \hat{\bx} \times \nabla\nn\ee
%
In components, with $\bx = (x_1,x_2,x_3)$, these operators are
%
\be \hat{L}_1 &=& -i\hbar\left(x_2\ppp{}{x_3} -x_3\ppp{}{x_2}\,\right)\ 
\nn\\  \hat{L}_2&=& -i\hbar\left(x_3\ppp{}{x_1} -x_1\ppp{}{x_3}\,\right) \ \\ 
 \hat{L}_3 &=& -i\hbar\left(x_1\ppp{}{x_2} -x_2\ppp{}{x_1}\,\right)\ .\nn\ee
%
Alternatively, these can be summarised in index notation as
%
\be \hat{L}_i = -i\hbar\,\epsilon_{ijk}\, x_j\ppp{}{x_k}\nn\ee
%
}
It's not hard to check that the angular momentum operators are Hermitian. Indeed, this really follows immediately from the fact that both $\hat{\bx}$ and $\hat{\bp}$ are Hermitian.


\para
{\co The three angular momentum operators do not mutually commute. They obey the relations
%
\be  [\hat{L}_i,\hat{L}_j] = i\hbar \epsilon_{ijk} \hat{L}_k\nn\ee
%
This  that a quantum particle cannot have a well-defined angular momentum in all three directions simultaneously. } If, for example, you know the angular momentum $L_1$ then there will necessarily be some uncertainty in the angular momentum in the other two directions.

\para
\alyssa{Showing the commutation relations of spin operators will be part of the practice problem for Lecture 19} \href{https://drive.google.com/file/d/1rzUg-9TFLG4pKcETaEh5WDlRMdxFGVMX/view?usp=drive_link}{[Lecture 19 Problem Statement]}
{\cblue Possible exercise: show the commutation relations of spin operators are
%
\be [\hat{S}^i,\hat{S}^j] = i\hbar \epsilon^{ijk}\hat{S}^k \nn\ee
%
Comment also that whenever these same commutation relations crop up, there's typically something to do with angular momentum behind them. What's really going on is that this is the structure associated to the group $SO(3)$ of rotations.}





\para
{\co Let's prove the commutation relations. It will suffice to prove just one case, say
%
\be [\hat{L}_1,\hat{L}_2] = i\hbar \hat{L}_3\nn\ee
%
To check this, we recall that operator equations of this kind should be viewed in terms of their action on functions and then just plug in the appropriate definitions,
%
\be [\hat{L}_1,\hat{L}_2]\psi &=& -\hbar^2\left(x_2\ppp{}{x_3} -x_3\ppp{}{x_2}\,\right)\left(x_3\ppp{}{x_1} -x_1\ppp{}{x_3}\,\right) \psi \nn\\ &&\ \ \ \  +\hbar^2 \left(x_3\ppp{}{x_1} -x_1\ppp{}{x_3}\,\right)  \left(x_2\ppp{}{x_3} -x_3\ppp{}{x_2}\,\right)  \psi \ \nn\ee
%
The important point is to remember that the derivatives in the first brackets act on the $x_i$ factors in the second, as well as on $\psi$. 
A little algebra will convince you that nearly all terms cancel. The only ones that survive are
%
\be [\hat{L}_1,\hat{L}_2]\psi = \hbar^2\left(x_1\ppp{}{x_2} -x_2\ppp{}{x_1}\,\right)\psi =i\hbar\hat{L}_3\psi\nn\ee
%
Because this holds for all  functions $\psi(\bx)$, it is equivalent to the claimed result.} The other commutation relations  follow in the same vein.



{\co We can also look at the magnitude of the  angular momentum. It's more useful to look at the ${\bf L}^2$ rather than $|{\bf L}|$. This is associated to the {\it total angular momentum} operator
%
\be \hat{\bf L}^2 = \hat{L}_1^2 +\hat{L}_2^2 + \hat{L}_3^2  \nn\ee
%
This is again a Hermitian operator. It commutes with all $\hat{L}_i$
%
\be [\hat{\bf L}^2,\hat{L}_i] = 0\nn\ee
%
} This is important. As we saw,  operators that commute can be simultaneously diagonalised. So while quantum mechanics
doesn't allow us to assign simultaneous values to the angular momentum in all directions, it is possible for a  quantum state to have a specific total angular momentum -- meaning that it is an eigenstate of $\hat{\bf L}^2$ -- and also a specific value of angular momentum in {\it one} chosen direction -- meaning that it is also an eigenstate of, say, $\hat{L}_3$.

\para
{\co To prove this we will first need a simple lemma. For any operators $\hat{A}$ and $\hat{B}$ 
%
\be [\hat{A}^2,\hat{B}] = \hat{A}[\hat{A},\hat{B}] + [\hat{A},\hat{B}]\hat{A}\nn\ee
%
This is simple to check:  we start from the right-hand side and expand
%
\be \hat{A}[\hat{A},\hat{B}] + [\hat{A},\hat{B}]\hat{A} &=& \hat{A}^2\hat{B} - \hat{A}\hat{B}\hat{A} + \hat{A}\hat{B}\hat{A}-\hat{B}\hat{A}^2 \nn\\ &=& \hat{A}^2\hat{B} - \hat{B}\hat{A}^2 = [\hat{A}^2,\hat{B}]\nn\ee
%
We now just use this result repeatedly to prove what we want. We want to show that, for example
%
\be [\hat{L}_1^2,\hat{L}_1] +  [\hat{L}_2^2,\hat{L}_1] + [\hat{L}_3^2,\hat{L}_1] =0\nn\ee
%
}
The first of these is trivial: all operators commute with themselves so $[\hat{L}_1^2,\hat{L}_1]=0$. The next term can be written as 
%
{\co \be [\hat{L}_2^2,L_1] &=& \hat{L}_2[\hat{L}_2,L_1] + [\hat{L}_2,\hat{L}_1]\hat{L}_2 \nn\\ &=&
-i\hbar( \hat{L}_2\hat{L}_3 +\hat{L}_3\hat{L}_2)\label{l22}\ee
%
where, to get to the second line, we've used the commutation relations. The final term is
%
\be [\hat{L}_3^2,L_1] &=& \hat{L}_3[\hat{L}_3,L_1] + [\hat{L}_3,\hat{L}_1]\hat{L}_3 \nn\\ &=&
i\hbar( \hat{L}_3\hat{L}_2 + \hat{L}_2\hat{L}_3)\ . \label{l23}\ee
%
Adding these gives us the result we wanted.}






\subsection{The Eigenfunctions are Spherical Harmonics}\label{sphericalharmsec}




Our next task is to figure out the eigenfunctions and eigenvalues of the angular momentum operators $\hat{\bf L}$. Things become significantly simpler if we work in spherical polar coordinates. 

\para
{\co We work in spherical polar coordinatesbe
%
\be x_1 &=& r\sin\theta\cos\phi \nn\\  x_2 &=& r\sin\theta\sin\phi \nn\\  x_3 &=& r\cos\theta\ .\nn\ee
%
with ranges $0\leq \theta\leq \pi$ and $0\leq \phi < 2\pi$.} {\cblue include figure here.}






\para
{\co We want to write out angular momenntum operators in polar coordinates. For this, we need the chain rule. For example,
%
\be \ppp{}{x_1} &=& \ppp{r}{x_1}\ppp{}{r} + \ppp{\theta}{x_1}\ppp{}{\theta} + \ppp{\phi}{x_1}\ppp{}{\phi}\nn\\ 
&=& \sin\theta\cos\phi\ppp{}{r} + \cos\theta\cos\phi \frac{1}{r}\ppp{}{\theta} -\frac{\sin\phi}{\sin\theta} \frac{1}{r}\ppp{}{\phi}\ \nn\ee
%
with similar results for $\partial/\partial x_2$ and $\partial/\partial x_3$. 

\para
Combining these, we can write the angular momentum operators as
%
\be \hat{L}_1 &=& i\hbar\left(\cot\theta\cos\phi \ppp{}{\phi} + \sin\phi\ppp{}{\theta}\right)\ \nn\\
\hat{L}_2 &=& i\hbar\left(\cot\theta\sin\phi \ppp{}{\phi} - \cos\phi \ppp{}{\theta}\right)\ \nn\\
\hat{L}_3 &=& -i\hbar \ppp{}{\phi}\nn\ee
%
}
Note that the angular momentum operators care nothing for how the wavefunctions depend on the radial coordinate $r$. The angular momentum of a state is, perhaps unsurprisingly, encoded in how the wavefunction varies in the angular directions  $\theta$ and $\phi$. Furthermore, $\hat{L}_3$ takes a  particularly straightforward form, simply because $x_3$ is singled out as a special direction in spherical polar coordinates.


\para
{\co The total angular momentum operator becomes
%
\be \hat{\bf L}^2 = \hat{L}_1^2+\hat{L}_2^2 + \hat{L}_3^2 = -\frac{\hbar^2}{\sin^2\theta} \left[\sin\theta\ppp{}{\theta}\left(\sin\theta\ppp{}{\theta}\right) + \ppp{^2}{\phi^2}\right]\nn\ee
%
}
Again, we should be careful in interpreting what this means. In particular, the first $\partial/\partial\theta$ acts both on the $\sin\theta$ in round brackets and  on whatever function the operator $\hat{\bf L}^2$ hits. 

\para
\alyssa{Deriving the total angular momentum operator is a part of the practice problem for Lecture 19} \href{https://drive.google.com/file/d/1rzUg-9TFLG4pKcETaEh5WDlRMdxFGVMX/view?usp=drive_link}{[Lecture 19 Problem Statement]}
{\cblue We could ask them to derive this as a problem. Would probably need us to walk it through slowly though -- I think there's a subtlety in getting the chain rule thing correct so it might be worth it.}

\para
Our goal is to find simultaneous eigenfunctions for $\hat{L}_3$ and $\hat{\bf L}^2$. Because nothing depends on the radial direction $r$, it's natural to look at  wavefunctions that depend only on $\theta$ and $\phi$. (We'll put the radial $r$ dependence back in in the next lecture.). You can think of them as wavefunctions defined on a unit sphere.

\para
{\co We look for eigenfunctions $\psi(\theta,\phi)$. 
 The $\hat{L}_3$ eigenfunction is straightforward: we can take 
%
\be \psi(\theta,\phi) = Y(\theta) e^{im\phi}\nn\ee
%
for any function $Y(\theta)$. The wavefunction is single-valued in $\phi$, meaning $\psi(\theta,\phi) = \psi(\theta,\phi+2\pi)$, only if $m$ is an integer. Using $\hat{L}_3$, we see that the corresponding eigenvalues are
%
\be \hat{L}_3 e^{im\phi} = m\hbar\,e^{im\phi}\ \ \ \ m\in \Z\nn\ee
%
Angular momentum in the $z=x_3$ direction is quantised in units of $\hbar$.} So too, of course, is the angular momentum in any other direction. But, as we've stressed above, if we sit in an eigenstate of $\hat{L}_3$, then we can't also sit in an eigenstate of any other direction. 

\para
{\co For the  $\hat{\bf L}^2$ eigenvalue equation, we look for solutions
%
\be  \hat{\bf L}^2 \,Y(\theta) e^{im\phi} = \lambda Y(\theta) e^{im\phi} \nn\ee
%
with $\lambda$ the eigenvalue. Using our expression for $\hat{\bf L}^2$, this becomes
% 
\be
 && -\frac{\hbar^2}{\sin^2\theta} \left[\sin\theta\ppp{}{\theta}\left(\sin\theta\ppp{}{\theta}\right) + \ppp{^2}{\phi^2}\right] Y(\theta)e^{im\phi} = \lambda Y(\theta) e^{im\phi}\nn\\
 \Longrightarrow \ \ \ && -\frac{\hbar^2}{\sin^2\theta} \left[\sin\theta\ppp{}{\theta}\left(\sin\theta\ppp{}{\theta}\right) -m^2\right] Y(\theta) = \lambda Y(\theta) \nn\ee
 %
Solutions to this equation are a well-studied class of mathematical functions called  {\it associated Legendre polynomials}.}  Here we give a brief review of their properties. \alyssa{We will do this review of Legendre polynomials in a Tutorial for Lecture 19} \href{https://drive.google.com/file/d/1xhpN1CYT8GnAk_H8X7QlE6gxoBscc3MD/view?usp=drive_link}{[Lecture 19, Tutorial 1]}


\para
{\cblue This is fairly involved...maybe we don't go through this?

\para
The construction of these solutions proceeds in two steps. First, we introduce a set of functions known as (ordinary) {\it  Legendre polynomials}, $P_l(w)$. These obey the differential equation\index{associated Legendre polynomial}
%
\be \frac{d}{dw}\left[(1-w^2)\frac{dP_l}{dw}\right] + l(l+1)P_l(w) =0\ \ \ {\rm with}\ \  l=0,1,2,\ldots\nn \ee
%
The functions $P_l(w)$, with $l=0,1,2,\ldots$, are polynomials in $w$ of degree $l$. 

\para
What does this have to do with our eigenvalue equation? First make the substitution $w=\cos\theta$. Then, after some chain-ruley things (using $d/dw = -(1/\sin\theta) d/d\theta$), you will find that the equation above coincides with our previous equation with $m=0$. This means that $Y(\theta) = P_l(\cos\theta)$ is an eigenfunction for $m=0$, with the eigenvalue 
%
 \be \lambda = l(l+1) \hbar^2\nn\ee
 %
 It's also straightforward to generate solutions with $m\neq 0$ by diffentiating $P_l(w)$. The general eigenfunction solution is  the so-called associated Legendre polynomials $P_{l,m}(\cos\theta)$, defined by 
 %
 \be Y(\theta) = P_{l,m}(\cos\theta) = (\sin\theta)^{|m|}\,\frac{d^{|m|}}{d (\cos\theta)^{|m|}} P_l(\cos\theta)\nn\ee
 %
 This too has the same eigenvalue $\lambda = l(l+1)\hbar^2$  independent of $m$. 

\para
Because $P_l(\cos\theta)$ is a polynomial of degree $l$, you only get to differentiate it $l$ times before it vanishes.  This means that the $\hat{L}_3$ angular momentum is constrained to take values that lie in the range
 %
\be -l \leq m \leq l\nn\ee
 %
Need to think more about how much of the above we do explicitly. We certainly need that last equation $-l\leq m\leq m$. }

 \para
 {\co 
The upshot of this analysis is that the simultaneous eigenstates of $\hat{\bf L}^2$ and $\hat{L}_3$ are given by functions that are labelled by two integers, $l$ and $m$, with $|m|\leq l$
%
\be \psi(\theta,\phi) = Y_{l,m}(\theta,\phi)\ .\ee
%
These functions are called {\it spherical harmonics}  and are given by
%
\be Y_{l,m}(\theta,\phi) = P_{l,m}(\cos\theta)e^{im\phi}\nn\ee
%
The corresponding eigenvalues are 
%
\be \hat{\bf L}^2Y_{l,m}(\theta,\phi) &=& l(l+1)\hbar^2\,Y_{l,m}(\theta,\phi)\nn\\ \hat{L}_3 Y_{l,m}(\theta,\phi) &=&m\hbar \,Y_{l,m}(\theta,\phi)\nn\ee
%
The integer $l$ is called the {\it total angular momentum quantum number} (even though, strictly speaking, the total angular momentum is really $\sqrt{l(l+1)}\hbar$.) The integer $m$ is called the {\it azimuthal angular momentum}.
}





\para
The requirement that $|m|\leq l$ is simply telling us that the angular momentum in any given direction (in this case $z$) can't be more than the total angular momentum. When $m=+l$ (or $m=-l$) you should think of the object as having angular momentum maximally aligned (or anti-aligned) with the $z$-axis. When $|m|<l$, you should think of the angular momentum as aligned in some intermediate direction. Of course, all of these thoughts should be quantum thoughts: even when $m=l$, if you measure the angular momentum in a perpendicular direction like $\hat{L}_1$, then you're not guaranteed to get zero. Instead there will be a probability distribution of answers centred around zero.




\para
To build some intuition, here are some of the first spherical harmonics. {\co The $l=0$ case is just a constant\index{s-wave}
%
\be Y_{0,0} = 1\nn\ee
%
This is the situation where a particle has no angular momentum and, correspondingly, its wavefunction has no angular dependence. In the context of atomic physics, it is referred to as an {\it s-wave}. The next collection of spherical harmonics all have total angular momentum $l=1$
%
\be Y_{1,-1} = \sin\theta\, e^{-i\phi}\ ,\ \ \ Y_{1,0} = \cos\theta\  ,\ \ \ Y_{1,1} = \sin\theta \, e^{i\phi}\nn\ee
%
In the context of atomic physics, these are called {\it p-wave} states. The next, $l=2$, collection are {\it d-waves}
%
\be Y_{2,\pm 2} = \sin^2\theta \, e^{\pm 2i\phi}  ,\ \ \ Y_{2,\pm 1}  = \sin\theta\,\cos\theta\,e^{\pm i\phi}  ,\ \ \ Y_{2,0} = 3\cos^2\theta -1\nn\ee
%
}
There are various ways to visualise spherical harmonics. {\cblue show different plots of spherical harmonics here...two different visualisations are shown it the book.}


\para
{\cred Yeah -- this lecture needs a lot of writing. I'm pretty sure it's closer to 45 mins to an hour rather than 30 mins.}






\newpage
\section{Solving the 3d Schrodinger Equation}

So we've understood angular momentum, and have an expression for the eigenstates. We'll now see that this goes a long way to helping us solve the Schrodinger equation for a particle moving in 3d.


\para
{\co With a central potential with Hamiltonian is 
%
\be \hat{H} = \frac{\hat{\bf p}^2}{2m}+ V(r)\nn\ee
%
}
In classical mechanics, the central potential implies that angular momentum is conserved. The corresponding statement in quantum mechanics is that: {\co For a central potential, we have
%
\be [\hat{H},\hat{L}_i] = [\hat{H},\hat{\bf L}^2] = 0\nn\ee
%
}
This means that eigenstates of the Hamiltonian can also be chosen to be eigenstates of the total angular momentum and of angular momentum in {\it one} chosen direction. A quantum particle can have simultaneous values for all three of these quantities. This is the result that will allow us to solve the 3d Schrodinger equation.

\para
Before we can do that, we first have to prove that they commute. {\co This follows from some straightforward algebra that is very similar in spirit to what we've done so far. You can first show that
%
\be [\hat{L}_i,\hat{x}_j] = i\hbar\,\epsilon_{ijk} \hat{x}_k\ \ \ {\rm and}\ \ \  [\hat{L}_i,\hat{p}_j] = i\hbar\,\epsilon_{ijk} \hat{p}_k\nn\ee
%
From this, it's a small step to show that 
%
\be [\hat{L}_i,\hat{\bf x}^2] =  [\hat{L}_i,\hat{x}_1^2+\hat{x}_2^2+\hat{x}_3^2] &=& 0\nn\\ {\rm and}\ \ \ 
 [\hat{L}_i,\hat{\bf p}^2] =  [\hat{L}_i,\hat{p}_1^2+\hat{p}_2^2+\hat{p}_3^2] &=& 0 \nn\ee
 %
But this is all we need. The Hamiltonian  is a function only of $\hat{\bf p}^2$ and $\hat{\bf x}^2$ where $r^2 = \bx^2$.  So we're guaranteed to get 
%
\be [\hat{H},\hat{L}_i]=0 \nn\ee
%
Since the Hamiltonian commutes with each angular momentum operator, it also commutes with the total angular momentum operator $\hat{\bf L}^2$.}

\para
The upshot of this analysis is that: {\co A particle moving in a central potential can sit in  simultaneous eigenstates of three operators, which we  take to be
%
\be \hat{H},\ \hat{\bf L}^2,\ {\rm and}\ \hat{L}_3\nn\ee
%
The associated eigenvalues are energy, total angular momentum, and the angular momentum in the $z=x_3$ direction. In what follows, all energy eigenstates are labelled by these three numbers. }

\subsection*{From PDE to ODE, Using Angular Momentum}

Now we can see explicitly how these ideas help us. {\co  Using $\hat{\bf p} = -i\hbar\nabla$, the Hamiltonian is
%
\be \hat{H} = -\frac{\hbar^2}{2m}\nabla^2 + V(r)\nn\ee
%
}
The symmetry of the problem strongly suggests that we work in spherical polar coordinates. {\co In spherical polars, the Laplacian takes the form
%
\be \nabla^2 = \frac{1}{r^2}\ppp{}{r}\left(r^2\ppp{}{r}\right) + \frac{1}{r^2\sin\theta}\ppp{}{\theta}\left(\sin\theta \ppp{}{\theta} \right)+ \frac{1}{r^2\sin^2\theta}\ppp{}{\phi^2} \nn\ee
%
The angular parts of the Laplacian are identical to the total angular momentum operator $\hat{\bf L}^2$. This means that we can equally well write the Laplacian as 
%
\be \nabla^2 = \frac{1}{r^2}\ppp{}{r}\left(r^2\ppp{}{r}\right) -\frac{\hat{\bf L}^2}{\hbar^2r^2} \nn\ee
}
Now we can see the utility of working with angular momentum eigenstates. {\co  We look for solutions to the Schrodinger equation of the form
%
\be \psi(r,\theta,\phi) = R(r) Y_{l,m}(\theta,\phi)\ee
%
where $R(r)$ is the radial wavefunction. 
All of the angular dependence now sits in an eigenstate $Y_{l,m}(\theta,\phi)$ of $\hat{\bf L}^2$. }

\para
Substituting this into the Schrodinger equation, what looks like a complicated partial differential equation (or PDE) turns into a more straightforward ordinary differential equation  (ODE). {\co We have
%
\be -\frac{\hbar^2}{2mr^2}\frac{d}{dr}\left(r^2\frac{dR}{dr}\right) + V_{\rm eff}(r) R = ER\nn\ee
%
where the potential $V(r)$ is replaced with the effective potential
%
\be V_{\rm eff}(r) = V(r) +\frac{l(l+1)\hbar^2}{2mr^2}\ .\nn\ee
%
}
If you've done a course on classical mechanics, you seen something like this before. It's exactly the same as the effective potential that emerges when solving problems in a central potential.  The additional term is called the {\it angular momentum barrier}. The  only difference from our earlier, classical result is that the  angular momentum (squared) is now quantised in \mbox{units of $\hbar^2$.}  

\para
Note that the azimuthal angular momentum quantum number $m$ doesn't appear in the Hamiltonian. (Sadly for us, there is an ``$m$" sitting in both the kinetic term and in the effective potential, but this is the mass. Not the angular momentum quantum number. Sorry! The canonical name for both is $m$. It's rare that it leads to confusion precisely because the azimuthal quantum number doesn't appear in the Hamiltonian.)

\para
{\co Because the energy doesn't depend on the $\hat{L}_3$ eigenvalue,  the energy spectrum for any particle moving in a central potential will always be degenerate, with multiple eigenstates all having the same energy. In contrast, the energy will usually depend on the total angular momentum $l$. This means that we expect each energy eigenvalue to have 
%
\be {\rm degeneracy} = 2l+1\nn\ee
%
which comes from the different $Y_{l,m}$ states with $m=-l,-l+1,\ldots, l-1,l$.}


\para
The effective Schrodinger equation similar form to those that we solved earlier in these lecturs.  The only slight difference is the more complicated kinetic term, together with some associated issues about the need for the wavefunction to be finite at $r=0$. 



\para
{\cblue There were sections on the 3d harmonic oscillator and on bound states in 3d wells in the book. I've deleted them from the lectures. Maybe we could reinstate one as an exercise? Or do you want to keep them? It would be another 20-30 minutes I think to do the 3d harmonic oscillator in spherical polar, although it's a useful exercise.} \alyssa{I think we can keep the 3d harmonic oscillator as a Tutorial for Lecture 20 and the bound states in 3d wells as a practice problem. Could you put your notes of what you would like included back here or email them to me? alyssa.rudelis@gmail.com}



\newpage
\section{The Hydrogen Atom}

The hydrogen atom is the simplest of all atoms. {\co It consists of a single electron orbiting a proton, held in place by the Coulomb force
%
\be F(r) = -\ppp{V}{r}= -\frac{e^2}{4\pi \epsilon_0r^2}\nn\ee
%
Here the electron has charge $-e$ and the proton has charge $+e$. The constant $\epsilon_0$ is a way to characterise the strength of the electrostatic force.} For now, we'll give values for neither  $e$ nor $\epsilon_0$, but we'll return to this after we've computed the energy spectrum.  {\co The associated Coulomb potential is 
%
\be V(r) = -\frac{e^2}{4\pi\epsilon_0 r}\nn\ee
%
}
Our task in this section is to solve the 3d Schrodinger equation with this potential.

\para
{\co We know that solutions will be characterised by three integers, $n$, $l$, and $m$, such that:
%
\begin{itemize}
\item $n$ labels the energy state. We will take it to run from $n=1,2,3,\ldots$. In the context of the hydrogen atom, it is sometimes called the {\it principal quantum number}. \vskip 3mm
\item $l$ labels the total angular momentum. It appears as  the eigenvalue $l(l+1)\hbar^2$ of  $\hat{\bf L}^2$. \vskip 3mm
\item $m$ labels the azimuthal angular momentum, also known as the angular momentum in the $z$-direction. It appears as the eigenvalue $m\hbar$ of $\hat{L}_3$.
\end{itemize}
%
We look for wavefunctions of the form
%
\be\psi (r,\theta,\phi) = R(r) \,Y_{l,m}(\theta,\phi)\nn\ee
%
The Schrodinger equation  for the hydrogen atom becomes
%
\be -\frac{\hbar^2}{2m_e}\left(\frac{d^2 R}{dr^2} + \frac{2}{r}\frac{dR}{dr}\right) +\left(-\frac{e^2}{4\pi\epsilon_0r} +\frac{l(l+1)\hbar^2}{2m_er^2}\right)R = ER\nn\ee
%
where $m_e$ is the electron mass} (and the subscript helps avoid confusion with the azimuthal angular momentum). 
Before we get going, let's introduce some new variables so that we don't have to lug around all those constants on our journey. {\co We write
%
\be \beta =\frac{e^2m_e}{2\pi\epsilon_0\hbar^2}\ \ \ {\rm and}\ \ \ \frac{1}{a^2} = -\frac{2m_eE}{\hbar^2}\ .\nn\ee
%
} In the second of these redefinitions, we've anticipated the fact that we're looking for  bound states of the Coulomb potential so that $E<0$ and, correspondingly,  $a^2>0$. {\co You can check that $a$ has dimension of length. \alyssa{We will have students check that $a$ has units of length in the practice problem for Lecture 21.} \href{https://drive.google.com/file/d/1NqP1CT6zZzU0gjhPj_C3M_ct2S-bQVNS/view?usp=drive_link}{[Lecture 21 Problem Statement]} This cleans up the Schrodinger equation a little, which now reads
%
\be  \frac{d^2 R}{dr^2} + \frac{2}{r}\frac{dR}{dr} - \frac{l(l+1)}{r^2} R  +\frac{\beta}{r}R = \frac{R}{a^2}\ . \nn\ee
%
}
All we have to do now is solve it.

\para
{\co First, let's look at the asymptotic behaviour of solutions as  $r\rightarrow \infty$. Here, any term in the Schrodinger equation   that has a $1/r$ or $1/r^2$ suppression will be sub-leading, so we're left only with the terms 
%
\be \frac{d^2R}{dr^2} \approx \frac{R}{a^2}\ \ \ \Longrightarrow\ \ \ R(r) \approx e^{-r/a}\ \ \ {\rm as} \ r\rightarrow \infty\nn\ee
%
}  We see that $a$, which  scales as $a\sim 1/\sqrt{-E}$,    sets the  characteristic size of the wavefunction. You can check, retrospectively, that the terms we dropped from the Schrodinger equation are indeed unimportant in the $r\ra \infty$ limit when evaluated on the solution $R\approx e^{-r/a}$.

\para
What happens near the origin at $r=0$? Here the $R/r^2$ term in the Schrodinger equation will clearly dominate the $R/r$ and $R$ terms. But what about the derivatives? If we anticipate that the solution has power-law behaviour, then taking a derivative with respect to $r$ is, as far as scaling goes, equivalent to dividing by $r$. This means that {\co Near $r=0$, the dominant terms in the Schrodinger equation  are
%
\be \frac{d^2R}{dr^2} + \frac{2}{r}\frac{dR}{dr} - \frac{l(l+1)}{r^2}R \approx 0\ \ \ {\rm for}\ r\ll 1\nn\ee
%
We make the power-law ansatz $R\sim r^\alpha$, we find
%
\be \alpha(\alpha-1) + 2\alpha - l(l+1) = 0\ \ \ \Longrightarrow \ \ \ R \sim r^l\ \ {\rm as}\ r\rightarrow 0\ee
%
We discard $\alpha =-(l+1)$ because it results in a wavefunction that diverges at the origin. } 

\para
{\cblue There is a subtlety here about why exactly we throw away the wavefunctions that diverge at the origin. It is treated properly in a section in the book, but it might be worth brushing under the rug here. Or perhaps having an online aside?} \alyssa{We will cover this in a Tutorial for Lecture 21.} \href{https://drive.google.com/file/d/1jBoanjhSDT81TeGmQopiHUfDKjkdjTqB/view?usp=drive_link}{[Lecture 21, Tutorial 1]}


\para
{\co All of this motivates us to look for solutions of the form
%
\be R(r) = r^l\, f(r) e^{-r/a}\nn\ee
%
where $f(r)$ is a polynomial
%
\be f(r) = \sum_{k=0}c_k r^k\nn\ee
%
and where we must have $c_0\neq 0$ so that we get the right behaviour at small $r$. }

\para
Now we've got to roll up our sleeves and substitute this into the Schrodinger equation to get a recurrence relation for the coefficients $c_k$. It's best to go slowly. First you can check that {\co If you plug the ansatz  into the Schrodinger equation then the function $f(r)$ must satisfy
%
\be \frac{d^2f}{dr^2} + 2\left(\frac{l+1}{r} - \frac{1}{a}\right)\frac{df}{dr} - \frac{1}{ar}\left(2(l+1) - \beta a\right)f=0\nn\ee
%
Next, substitute the power-law ansatz  into this differential equation to find
%
\be \sum_{k=0}^\infty c_k\left[ (k(k-1) + 2k(l+1))r^{k-2} - \frac{1}{a}(2k + 2(l+1) - \beta a) r^{k-1}\right]=0\ .\ \ \ \ \ \ \, \nn\ee
 %
Finally, shift the dummy summation variable $k\rightarrow k-1$ in the second term to get a uniform power of $r^{k-2}$. This gives us a recurrence relation
%
\be c_k = \frac{1}{ak} \frac{2(k+l) -\beta a}{k+2l+1}\,c_{k-1}\ .\nn\ee
%
}
Now there's a clever argument {\cblue which I probably need to flesh out here, not least because I ducked it when doing the harmonic oscillator in 1d. Perhaps we should think about revisiting that lecture and doing the recurrence relation proof as well so that it's more familiar when it appears here. I don't think it's avoidable here unless we just give the answer without deriving it.} \alyssa{We will put the recurrence relation proof in an explainer for Lecture 7.} \href{https://drive.google.com/file/d/1-7a5lutZlVQALBmbVPMVMWuM6sbvbxFO/view?usp=drive_link}{[Lecture 7, Tutorial 2]}

\para
{\co Suppose that this recurrence relation fails to terminate. Then, as $k\rightarrow \infty$, we have 
%
\be c_{k} \rightarrow \frac{2}{ak} c_{k-1}\ .\ee
%
But this is the kind of expansion that we get from exponentially divergent functions of the form $f(r) = e^{+2r/a}$.} ({\cblue I think this is the step that needs fleshing out...maybe it's straightforward to just do it here.}) On a mathematical level, there's nothing wrong with these solutions: they just change the asymptotic form  from $e^{-r/a}$ to $e^{+r/a}$. But we've no use of such solutions for physics because they're badly non-normalisable.

\para
{\co This means that we  need to restrict attention to series that terminate. Which, in turn, means that there should be a positive integer $q$ for which $c_q=0$, while $c_k\neq 0$ for all $k<q$. Clearly this holds only if $a$ takes special values,
%
\be a = \frac{2}{\beta} (q+l)\ \ \ \ {\rm with}\ q=1,2,\ldots\ .\nn\ee
%
Alternatively, since both $q$ and $l$ are integers, we usually define the integer
%
\be n=q+l\nn\ee
%
which obviously obeys  $n>l$. We then have
%
\be a = \frac{2}{\beta} n\nn\ee
%
}
All that's left is to chase through the definition of the various variables in this expression.  {\co The length scale $a$ is related to the energy scale $E$, which is given by
%
\be E = -\frac{e^4 m_e}{32\pi^2\epsilon_0^2 \hbar^2} \frac{1}{n^2}\ \ \ {\rm with}\  \  n=1,2,\ldots \nn\ee
%
This is our final result for the energy spectrum of the hydrogen atom. The integer $n$ coincides with what we previously called the {\it principal quantum number}. }

\para
There's a whole lot to unpick here. First, we can substitute this expression for $E$ into the rescaling that we performed in before to get an expression for the length scale $a$ that appears in the equations. {\co Rewriting, we have
%%
\be a = nr_B \ \ \ {\rm with}\ \ \ r_B = \frac{4\pi \epsilon_0\hbar^2}{e^2 m_e}  \approx 5.3\times 10^{-11}\ {\rm m}\ .\nn\ee
%
}
 We've factored out the state-dependent integer $n$ to get a length scale $r_B$, written purely in terms of fundamental constants. {\co Here $r_B$ is the {\it Bohr radius}. It sets the characteristic size of the hydrogen atom.}

\para
Next, there's the ugly morass of constants sitting on the right-hand side.  Clearly they set the energy scale of the bound states. There's actually a better way to organise them to see what's going on.  {\co We introduce the {\it fine structure constant}
%
\be \alpha = \frac{e^2}{4\pi \epsilon_0 \hbar c}\nn\ee
%
}
At first glance, you may think that this doesn't help matters because, rather than simplifying things, we've introduced  a new physical constant $c$, the speed of light. However, the advantage of the fine structure constant is that it's a dimensionless quantity. As you can see, it includes the expression  $e^2/4\pi \epsilon_0$ familiar from the  Coulomb force and it should be viewed as the convention-independent way to characterise the strength of electromagnetism. {\co It takes a value that's easy to remember
%
\be \alpha \approx \frac{1}{137} \approx 0.0073\nn\ee
%
Written in terms of the fine structure constant, the energy eigenvalues are
%
\be E = -\alpha^2 m_ec^2\frac{1}{2n^2}\nn\ee
%
}
This makes it clear what's really setting the energy scale of the hydrogen atom: {\co We see that the  binding energy is really set by the rest mass energy of the electron! This is given by
%
\be m_ec^2 \approx 0.51 \times 10^6 \ {\rm eV}\nn\ee
%
}
A million electronvolts is an  energy scale appropriate for  relativistic particle physics. It is  much larger than the typical energy scales of atomic physics, but we can see why: there are two powers of the fine structure constant that interpolate between the two scales. {\co The energy scale of the hydrogen atom is set by
%
\be {\rm Ry} = \frac{\alpha^2 m_ec^2}{2} \approx 13.6\ {\rm eV}\ .\nn\ee
%
This is called the {\it Rydberg}.} As you can see, the (negative) ground state energy of the hydrogen atom is exactly one Rydberg. {\co The higher excited states have energy $E = -{\rm Ry}/n^2$.} Note that the hydrogen atom has an infinite number of negative energy bound states.

\para
{\co
We can also write the Bohr radius  in terms of the fine structure constant
%
\be r_B  =\frac{\hbar}{m_e c\alpha} \nn\ee
%
The length scale $\hbar/m_e c$ is called the {\it Compton wavelength} of the electron.} Later, when we study quantum field theory, we will see that it should be viewed, in some sense, as the size of the electron itself. The Bohr radius is 137 times larger because of that factor of $1/\alpha$. 

\para
Next there's a surprise. The energy spectrum of hydrogen does not depend on the angular momentum $l$. At least, that's almost true: it depends only in the sense that the states for a given energy $E=-{\rm Ry}/n^2$ only include angular momentum $l<n$. 

\para
{\co The energy does not depend on the angular momentum $l$, except that at energy level $n$ we only have $l<n$. 



\para
The angular momentum does affect the degeneracy of the spectrum. Recall that there are $2l+1$ states  associated to each total angular momentum $l$. These are the states labelled by $m=-l,\ldots,+l$. For a fixed energy level $n$, we can have any total angular momentum $l=0,\ldots n-1$. This means that the degeneracy  of the  $n^{\rm th}$ energy level  is
%
\be \mbox{total degeneracy} = \sum_{l=0}^{n-1} (2l+1) = n^2\ \nn\ee
%
}
The $n=1$ ground state is unique: it sits in the  $l=0$ s-wave. There are four $n=2$ excited states corresponding to the s-wave and three $l=1$ states of the p-wave. There are nine $n=3$ excited states, corresponding to the s-wave, the three $l=1$ p-wave states, and five $l=2$ d-wave states. And so on.


\subsection{Wavefunctions}



The recurrence relation allows us to easily construct the corresponding wavefunctions. {\co A state with energy $E= -{\rm Ry}/n^2$ has a wavefunction
%
\be \psi_{n,l,m}(r,\theta,\phi) = r^l f_{n,l}(r) e^{-\beta r/2n} \,Y_{l,m}(\theta,\phi)\nn\ee
%
where  $\beta ={e^2m_e}/{2\pi\epsilon_0\hbar^2}$ and $f_{n,l}(r)$ is a  polynomial of degree $n-l-1$, defined by
%
\be f_{n,l}(r) = \sum_{k=0}^{n-l-1} c_k r^k\nn\ee
%
with 
\be c_k = \frac{2}{ak} \frac{k+l-n}{k+2l+1}\,c_{k-1} \nn\ee
%
These are known as {\it generalised Laguerre polynomials}. }
{\cblue They are more conventionally written as $f_{n,l}(r):= L_{n-l-1}^{2l+1}(2r/nr_B)$ but we don't need to mention this.} We still have to normalise these wavefunctions. {\co The correct normalisation turns out to give
%
\be \psi_{n,l,m}(r,\theta,\phi) = N_{n,l,m}r^le^{-r/nr_B} L^{2l+1}_{n-l-1}\left(\frac{2r}{nr_B}\right)\,Y_{l,m}(\theta,\phi)\ .\nn\ee
%
Here the normalisation factor is the rather messy
%
\be N_{n,l,m}=   \sqrt{\left(\frac{2}{nr_B}\right)^3\frac{(n-l-1)!}{2n(n+l)!}} \sqrt{\frac{(2l+1)}{4\pi}\frac{(l-m)!}{(l+m)!}} \nn\ee
%
}
 The wavefunctions for the lowest energy states are, thankfully, somewhat simple. The $n=1$ ground state wavefunction necessarily has vanishing angular momentum. {\co The ground state is
%
\be \psi_{1,0,0}(r) = \sqrt{\frac{1}{\pi r_B^3}} \,e^{-r/r_B}\  .\nn\ee
%
%
We now see clearly that the Bohr radius $r_B$ does indeed set the size of the ground state of the hydrogen atom.  The first excited states are
%
\be \psi_{2,0,0} &=&\sqrt{\frac{1}{32\pi r_B^3}}\left(2-\frac{r}{r_B}\right) e^{-r/2r_B}
\nn\\
\psi_{2,1,0} &=&\sqrt{\frac{1}{32\pi r_B^3}}\,\frac{r}{r_B}e^{-r/2r_B}\cos\theta
\nn\\
\psi_{2,1,\pm 1} &=& \mp \sqrt{\frac{1}{64\pi r_B^3}}\,\frac{r}{r_B} e^{-r/2r_B}\sin\theta\,e^{\pm i\phi}\nn\ee
%
} {\cblue plot wavefunctions here}

\para
{\cred This is another long lecture -- possibly closer to two lectures. Maybe we should split into two, although I'm not sure what were the split should lie -- I don't think it makes sense to do spectrum in one lecture and wavefunctions in another.}

\para
{\cblue Haven't thought about problems or explainers for this lecture -- sorry! There was a nice suggestion to make a tanget about uses of spectroscopy in astronomy. } 


\newpage
\section{Pre-Quantum Quantum Mechanics}


The formalism of quantum mechanics that we have seen in this book was laid down in  a flurry of creativity in the years 1925 and 1926. The primary builders of this new theory were Werner Heisenberg, Erwin Schr\"odinger, Max Born, Pascual Jordan, and Paul Dirac,  several of them  under the watchful grandfatherly eye of Niels Bohr, who acted variously as mentor, landlord, and spirit guide for the young generation of quantum pioneers.





\para
The emergence of the quantum theory that we know today was preceded by a quarter-century of confusion, a time in which it was known that some old classical ideas must be jettisoned to explain various experiments, but it was far from clear what should replace them.



\para {\cblue Do we do this in the very first lecture, or do it here?}
With hindsight, the first  hint of the need for  quantum mechanics came from the spectral lines of hydrogen. Take a tube filled with hydrogen gas and pass a current through it. The gas will glow but emit light only at very particular frequencies. In 1885, a Swiss school teacher called Johann Balmer noted that the wavelengths of the most prominent visible lines  could be fitted to the formula
%
\be \lambda = \frac{4}{Ry} \frac{m^2}{m^2-4}\nn \ee
%
for $m=3$ 4, 5, 6, and 7. These are known as  the {\it Balmer series}.

\para
Later, Rydberg generalised this to the formula
%
\be \frac{1}{\lambda} = Ry \left(\frac{1}{n^2} - \frac{1}{m^2}\right)\nn \ee
%
which coincides with the Balmer formula when $n=2$. Here, $Ry$ is known as the {\it Rydberg constant}. The more prominent $n=1$ lines appear only in the ultraviolet and were not known to Balmer. They are now referred to as the {\it Lyman series}.

\para
We now know that whenever integers appear in the laws of physics, they signify some underlying quantum effect. Indeed, the Rydberg formula follows straightforwardly from the spectrum of the hydrogen atom which, as we've seen, takes the form
%
\be E_n = -\frac{e^4 m_e}{32\pi^2\epsilon_0^2 \hbar^2} \frac{1}{n^2}\ \ \ \ n=1,2,\ldots \nn\ee
%
If an electron sits in an excited state $E_m$ and drops to a lower state $E_n$, $n<m$, it will emit a photon with energy
%
\be \hbar \omega  = E_m-E_n = \frac{e^4 m_e}{32\pi^2\epsilon_0^2 \hbar^2} \left(\frac{1}{n^2} - \frac{1}{m^2}\right)
\nn\ee
%
The wavelength $\lambda$ is related to the frequency $\omega$ by $\omega = 2\pi c/\lambda$, so this immediately reproduces the Rydberg formula with
%
\be Ry = \frac{e^4m_e}{64\pi^3 \epsilon_0^2 c \hbar^3}\nn\ee
%
Of course, going backwards from the answer is easy.  But, in the early 1900s, it was far from clear what to make of the Rydberg formula.



\para
The characteristics of these spectral lines also lend their name to the different angular momentum modes. As we have seen, states are labelled by the angular momentum quantum number $l=0,1,2,3,\ldots$. These are often called $s$, $p$, $d$, and $f$, respectively. The names are old fashioned and come from the observed quality of spectral lines; they stand for {\it sharp}, {\it principal}, {\it diffuse}, and {\it fundamental}, but they remain standard, both when describing atomic structure and beyond.

\subsubsection*{The Bohr Model}

An important halfway house between the classical and quantum worlds was provided by Bohr in 1913. This was only two years after Rutherford had shown that the atom consists of a tiny nucleus, with the electrons in remote orbits. Bohr took a collection of simple classical results and 
meshed them together with some quantum handwaving in a forced and somewhat unnatural manner to arrive at a remarkably accurate picture of the hydrogen atom. Here is his argument.

\para
{\co 
We start with the Coulomb force between an electron and proton
%
\be F = -\frac{e^2}{4\pi\epsilon_0 r^2}\nn\ee
%
}
We assume that the electron orbits the proton on a classical trajectory that we take to be a circle. {\co The Coulomb force must give the requisite centripetal force
%
\be F = -\frac{m_ev^2}{r}\nn\ee
%
From this, we can determine the radius of the orbit $r$ in terms of the angular momentum $L= mvr$ 
%
\be r= \frac{4\pi\epsilon_0 L^2}{e^2m_e} \nn\ee
%
Similarly, the total energy can also be expressed in terms of the angular momentum
%
\be E = \frac{1}{2}m_ev^2 - \frac{e^2}{4\pi\epsilon_0r} = -\frac{e^4m_e}{32\pi^2 \epsilon_0 L^2}\nn\ee
%
Now comes the quantum handwaving. Bohr postulated that angular momentum cannot take any value, but only integer amounts
%
\be L = n\hbar\ \ {\rm with}\  n=1,2,3,\ldots\nn\ee
%
This implies that the energy is quantised to take values
%
\be E = -\frac{e^4m_e}{32\pi^2\epsilon_0\hbar^2}\frac{1}{n^2}\ \nn\ee
%
}
Remarkably, this is the same as our exact result from solving the Schrodinger equation. {\co Furthermore, the minimum orbit of the electron is $r_B$, given by
%
\be r_B = \frac{4\pi \epsilon_0 \hbar^2}{e^2m_e}\nn\ee
%
This is the Bohr radius that we met before.}


\para
What's going on here? There is clearly much that is wrong with the Bohr model. It assumes a classical trajectory that  is far from the picture we now have of electrons as smeared clouds of probability (which we'll see more directly in the next lecture when we discuss wavefunctions.) It makes the wrong assumption for the quantisation of angular momentum, guessing that $L^2 = n^2\hbar^2$ rather than, as we now know, $L^2 = l(l+1)\hbar^2$. It also gets the connection between angular momentum and energy $l=n$ wrong; the correct answer should be $l<n$. And yet, the end result  is miraculously exactly right! How can this be?

\para
A pretty good answer to this question is simply: dumb luck. The hydrogen atom is unique among quantum systems in that these simple-minded half-classical, half-quantum approaches give the right answers to certain important questions, most notably the energy spectrum. It's not, however, a lesson that we can take and use elsewhere. In all other cases, to get the exact quantum result you have to do the hard work of solving the Schrodinger equation. The one exception to this is  the highly excited states, where Bohr's  kind of arguments do often provide a useful caricature of the physics. (For example, when the angular momentum is large there is not much difference between $L^2 = n^2\hbar^2$ and $L^2 = l(l+1)\hbar^2$.)


\para
{\cred Quite a short lecture. It seems a little strange to have the last lecture (or perhaps the penultimate lecture) on history. But perhaps that's ok}

\newpage
\section{Looking Forwards}

{\cblue As suggested, we need a wrap-up lecture here. Presumably advertising things to come. The UV-catastrophe and entanglement are obvious suggestions...I'll try to think of some other things and a way in which this could be done in the spirit of the lectures.}











\end{document}